% !TEX program = lualatex
% ---------------------------------------------------------------------------------------------
%  Readback note 32: Short Selmer groups after one anticyclotomic twist, balanced rows, and Ribet's lemma to first order
%
%  Build:  python3 tools/rb.py note <slug>
%
%  One unit of a Lean artifact, rewritten as ordinary mathematics. The prose is hand-written. Lean names,
%  line ranges and counts come from data.tex; quoted Lean statements are copied from the artifact by the build.
%
%  Rules of the text: standard terminology; every symbol is explained where it first appears and listed again
%  in the notation at the end; no letter has two meanings; a proof is complete or says what it takes for granted;
%  Lean names appear only in the LEAN lines and in the concordance.
%
%  The signature rule (the build enforces it): a function, relation or predicate that the artifact defines is
%  written only through its notation macro, which the build makes from the spec (`notation`) with one slot for
%  each explicit argument of the Lean definition. No argument is fixed and dropped, also not "for brevity" and
%  not for one section. Every argument and every hypothesis of a stated declaration stands inside its
%  statement, or in a Setting that the statement cites by number; only implicit and instance arguments (the
%  ambient objects) may stand in the conventions.
% ---------------------------------------------------------------------------------------------
\documentclass[11pt,a4paper]{article}
\ifdefined\pdfvariable\pdfvariable suppressoptionalinfo \numexpr 1+2+4+8\relax\fi
\usepackage[a4paper,hmargin=35mm,top=22mm,bottom=27mm,footskip=12mm,nohead]{geometry}

\usepackage{fontspec}
\directlua{
  luaotfload.add_fallback("gfallback", {"STIXTwoMath-Regular.otf:mode=node;", "DejaVuSans.ttf:mode=node;"})
  luaotfload.add_fallback("gmonofallback", {"DejaVuSansMono.ttf:mode=node;", "STIXTwoMath-Regular.otf:mode=node;"})
}
\setmainfont{NewCM10}[Extension=.otf,UprightFont=*-Book,ItalicFont=*-BookItalic,
  BoldFont=*-Bold,BoldItalicFont=*-BoldItalic,SmallCapsFont=*-Book,SmallCapsFeatures={Letters=SmallCaps},
  RawFeature={fallback=gfallback}]
\setsansfont{NewCMSans10}[Extension=.otf,UprightFont=*-Book,ItalicFont=*-BookOblique,
  BoldFont=*-Bold,BoldItalicFont=*-BoldOblique,RawFeature={fallback=gfallback}]
\setmonofont{JetBrainsMono}[Extension=.otf,UprightFont=*-Regular,ItalicFont=*-Italic,
  BoldFont=*-SemiBold,BoldItalicFont=*-SemiBoldItalic,Scale=MatchLowercase,
  RawFeature={-calt;-liga;fallback=gmonofallback}]
\usepackage{amsmath}
\usepackage{amsthm}
\usepackage[math-style=TeX]{unicode-math}
\setmathfont{NewCMMath-Book.otf}
\emergencystretch=2.5em\relax

\usepackage{xcolor}
\usepackage{tikz}
\usetikzlibrary{calc,arrows.meta,positioning,backgrounds}
\usepackage[protrusion=true,expansion=true]{microtype}
\usepackage{enumitem}
\usepackage{needspace}
\usepackage{booktabs}
\usepackage{array}
\usepackage{longtable}
\usepackage{fancyhdr}
\usepackage{lastpage}
\usepackage{multicol}
\usepackage[hidelinks,pdfusetitle]{hyperref}

\definecolor{ink}{HTML}{0F172A}
\definecolor{inksoft}{HTML}{475569}
\definecolor{inkfaint}{HTML}{94A3B8}
\definecolor{rulegrey}{HTML}{CBD5E1}
\definecolor{arith}{HTML}{C2410C}
\definecolor{est}{HTML}{7C3AED}
\definecolor{ana}{HTML}{0891B2}
\definecolor{ok}{HTML}{047857}
\definecolor{leancol}{HTML}{3730A3}

\title{Short Selmer groups after one anticyclotomic twist, balanced rows, and Ribet's lemma to first order}
\author{grh-compact}
\date{}

% ------------------------------------------------------------------------------------- macros
\newcommand{\leanfont}{\ttfamily\fontsize{8.6}{10.5}\selectfont\color{leancol}}
\newcommand{\lean}[1]{\csname LN@#1\endcsname}
\newcommand{\fillin}[1]{{\color{arith}\itshape[#1]}}  % a place of the template that is not written yet; the build reports every one that is left
\newcommand{\leanquote}[1]{\csname LQ@#1\endcsname}
% generated by rb.py note; do not edit
\newcommand{\ConcDecls}{58}
\newcommand{\Pin}{72fead10}
\newcommand{\PinText}{that of commit {\leanfont 72fead10}}
\newcommand{\ArtifactName}{grh-compact}
\expandafter\newcommand\csname LN@crudeUnit\endcsname{\href{https://github.com/lsernau/grh-lean/blob/72fead102f663d5c7a407964a13dc0cf9d977207/grh-compact/src/New/Selmer/AptTupleDefs.lean\#L27-L27}{{\leanfont crudeUnit}}}
\expandafter\newcommand\csname LN@IsColumn\endcsname{\href{https://github.com/lsernau/grh-lean/blob/72fead102f663d5c7a407964a13dc0cf9d977207/grh-compact/src/New/Selmer/AptTupleDefs.lean\#L43-L45}{{\leanfont IsColumn}}}
\expandafter\newcommand\csname LN@IsRow\endcsname{\href{https://github.com/lsernau/grh-lean/blob/72fead102f663d5c7a407964a13dc0cf9d977207/grh-compact/src/New/Selmer/AptTupleDefs.lean\#L50-L59}{{\leanfont IsRow}}}
\expandafter\newcommand\csname LN@obeyingCocycles\endcsname{\href{https://github.com/lsernau/grh-lean/blob/72fead102f663d5c7a407964a13dc0cf9d977207/grh-compact/src/New/Selmer/BlocksAndCocycles.lean\#L16-L17}{{\leanfont obeyingCocycles}}}
\expandafter\newcommand\csname LN@AuxPrimeData\endcsname{\href{https://github.com/lsernau/grh-lean/blob/72fead102f663d5c7a407964a13dc0cf9d977207/grh-compact/src/New/Selmer/Corepresentable.lean\#L38-L40}{{\leanfont AuxPrimeData}}}
\expandafter\newcommand\csname LN@IsAnticyclotomicAt\endcsname{\href{https://github.com/lsernau/grh-lean/blob/72fead102f663d5c7a407964a13dc0cf9d977207/grh-compact/src/New/Selmer/Corepresentable.lean\#L46-L51}{{\leanfont IsAnticyclotomicAt}}}
\expandafter\newcommand\csname LN@IsAuxiliary\endcsname{\href{https://github.com/lsernau/grh-lean/blob/72fead102f663d5c7a407964a13dc0cf9d977207/grh-compact/src/New/Selmer/Corepresentable.lean\#L55-L55}{{\leanfont IsAuxiliary}}}
\expandafter\newcommand\csname LN@IsFamily\endcsname{\href{https://github.com/lsernau/grh-lean/blob/72fead102f663d5c7a407964a13dc0cf9d977207/grh-compact/src/New/Selmer/Corepresentable.lean\#L60-L68}{{\leanfont IsFamily}}}
\expandafter\newcommand\csname LN@evenClasses\endcsname{\href{https://github.com/lsernau/grh-lean/blob/72fead102f663d5c7a407964a13dc0cf9d977207/grh-compact/src/New/Selmer/Corepresentable.lean\#L136-L138}{{\leanfont evenClasses}}}
\expandafter\newcommand\csname LN@IsNewPlaces\endcsname{\href{https://github.com/lsernau/grh-lean/blob/72fead102f663d5c7a407964a13dc0cf9d977207/grh-compact/src/New/Selmer/Corepresentable.lean\#L142-L143}{{\leanfont IsNewPlaces}}}
\expandafter\newcommand\csname LN@withConjInv\endcsname{\href{https://github.com/lsernau/grh-lean/blob/72fead102f663d5c7a407964a13dc0cf9d977207/grh-compact/src/New/Selmer/Corepresentable.lean\#L147-L147}{{\leanfont withConjInv}}}
\expandafter\newcommand\csname LN@IsPrimeTriangle\endcsname{\href{https://github.com/lsernau/grh-lean/blob/72fead102f663d5c7a407964a13dc0cf9d977207/grh-compact/src/New/Selmer/Corepresentable.lean\#L152-L162}{{\leanfont IsPrimeTriangle}}}
\expandafter\newcommand\csname LN@IsAuxTwist\endcsname{\href{https://github.com/lsernau/grh-lean/blob/72fead102f663d5c7a407964a13dc0cf9d977207/grh-compact/src/New/Selmer/Corepresentable.lean\#L166-L167}{{\leanfont IsAuxTwist}}}
\expandafter\newcommand\csname LN@unionWithConj\endcsname{\href{https://github.com/lsernau/grh-lean/blob/72fead102f663d5c7a407964a13dc0cf9d977207/grh-compact/src/New/Selmer/Corepresentable.lean\#L206-L206}{{\leanfont unionWithConj}}}
\expandafter\newcommand\csname LN@pointChar\endcsname{\href{https://github.com/lsernau/grh-lean/blob/72fead102f663d5c7a407964a13dc0cf9d977207/grh-compact/src/New/Selmer/Corepresentable.lean\#L210-L211}{{\leanfont pointChar}}}
\expandafter\newcommand\csname LN@twistAtPoint\endcsname{\href{https://github.com/lsernau/grh-lean/blob/72fead102f663d5c7a407964a13dc0cf9d977207/grh-compact/src/New/Selmer/Corepresentable.lean\#L215-L216}{{\leanfont twistAtPoint}}}
\expandafter\newcommand\csname LN@reduceModPow\endcsname{\href{https://github.com/lsernau/grh-lean/blob/72fead102f663d5c7a407964a13dc0cf9d977207/grh-compact/src/New/Selmer/Corepresentable.lean\#L219-L220}{{\leanfont reduceModPow}}}
\expandafter\newcommand\csname LN@IsBalanced\endcsname{\href{https://github.com/lsernau/grh-lean/blob/72fead102f663d5c7a407964a13dc0cf9d977207/grh-compact/src/New/Selmer/Corepresentable.lean\#L227-L238}{{\leanfont IsBalanced}}}
\expandafter\newcommand\csname LN@IsPencil\endcsname{\href{https://github.com/lsernau/grh-lean/blob/72fead102f663d5c7a407964a13dc0cf9d977207/grh-compact/src/New/Selmer/Corepresentable.lean\#L243-L250}{{\leanfont IsPencil}}}
\expandafter\newcommand\csname LN@tieRelations\endcsname{\href{https://github.com/lsernau/grh-lean/blob/72fead102f663d5c7a407964a13dc0cf9d977207/grh-compact/src/New/Selmer/Corepresentable.lean\#L254-L256}{{\leanfont tieRelations}}}
\expandafter\newcommand\csname LN@matrixLattice\endcsname{\href{https://github.com/lsernau/grh-lean/blob/72fead102f663d5c7a407964a13dc0cf9d977207/grh-compact/src/New/Selmer/CrystallineRank.lean\#L22-L23}{{\leanfont matrixLattice}}}
\expandafter\newcommand\csname LN@coeffEmb\endcsname{\href{https://github.com/lsernau/grh-lean/blob/72fead102f663d5c7a407964a13dc0cf9d977207/grh-compact/src/New/Selmer/DeterminantBound.lean\#L34-L35}{{\leanfont coeffEmb}}}
\expandafter\newcommand\csname LN@ZfAt\endcsname{\href{https://github.com/lsernau/grh-lean/blob/72fead102f663d5c7a407964a13dc0cf9d977207/grh-compact/src/New/Selmer/DeterminantBound.lean\#L39-L39}{{\leanfont ZfAt}}}
\expandafter\newcommand\csname LN@crysRels\endcsname{\href{https://github.com/lsernau/grh-lean/blob/72fead102f663d5c7a407964a13dc0cf9d977207/grh-compact/src/New/Selmer/DeterminantBound.lean\#L43-L43}{{\leanfont crysRels}}}
\expandafter\newcommand\csname LN@complexEmbOf\endcsname{\href{https://github.com/lsernau/grh-lean/blob/72fead102f663d5c7a407964a13dc0cf9d977207/grh-compact/src/New/Selmer/DeterminantBound.lean\#L59-L60}{{\leanfont complexEmbOf}}}
\expandafter\newcommand\csname LN@AbsGal\endcsname{\href{https://github.com/lsernau/grh-lean/blob/72fead102f663d5c7a407964a13dc0cf9d977207/grh-compact/src/New/Selmer/LevelPlacesChosen.lean\#L25-L25}{{\leanfont AbsGal}}}
\expandafter\newcommand\csname LN@placesAboveG\endcsname{\href{https://github.com/lsernau/grh-lean/blob/72fead102f663d5c7a407964a13dc0cf9d977207/grh-compact/src/New/Selmer/LevelPlacesChosen.lean\#L28-L28}{{\leanfont placesAboveG}}}
\expandafter\newcommand\csname LN@badPlaces\endcsname{\href{https://github.com/lsernau/grh-lean/blob/72fead102f663d5c7a407964a13dc0cf9d977207/grh-compact/src/New/Selmer/LevelPlacesChosen.lean\#L32-L32}{{\leanfont badPlaces}}}
\expandafter\newcommand\csname LN@cardletorsionofisotropic\endcsname{\href{https://github.com/lsernau/grh-lean/blob/72fead102f663d5c7a407964a13dc0cf9d977207/grh-compact/src/New/Selmer/LinePointPowers.lean\#L19-L53}{{\leanfont card\_le\_torsion\_of\_isotropic}}}
\expandafter\newcommand\csname LN@existssharpfittingpoly\endcsname{\href{https://github.com/lsernau/grh-lean/blob/72fead102f663d5c7a407964a13dc0cf9d977207/grh-compact/src/New/Selmer/LinePointPowers.lean\#L64-L86}{{\leanfont exists\_sharp\_fitting\_poly}}}
\expandafter\newcommand\csname LN@dvdoneaddmulpowsubone\endcsname{\href{https://github.com/lsernau/grh-lean/blob/72fead102f663d5c7a407964a13dc0cf9d977207/grh-compact/src/New/Selmer/LinePointPowers.lean\#L88-L95}{{\leanfont dvd\_one\_add\_mul\_pow\_sub\_one}}}
\expandafter\newcommand\csname LN@blockFlattenAlgEquiv\endcsname{\href{https://github.com/lsernau/grh-lean/blob/72fead102f663d5c7a407964a13dc0cf9d977207/grh-compact/src/New/Selmer/RibetCramer.lean\#L16-L18}{{\leanfont blockFlattenAlgEquiv}}}
\expandafter\newcommand\csname LN@amitsurcharpolymemmatrix\endcsname{\href{https://github.com/lsernau/grh-lean/blob/72fead102f663d5c7a407964a13dc0cf9d977207/grh-compact/src/New/Selmer/RibetCramer.lean\#L27-L35}{{\leanfont amitsur\_charpoly\_mem\_matrix}}}
\expandafter\newcommand\csname LN@existsribetelementcramer\endcsname{\href{https://github.com/lsernau/grh-lean/blob/72fead102f663d5c7a407964a13dc0cf9d977207/grh-compact/src/New/Selmer/RibetCramer.lean\#L40-L75}{{\leanfont exists\_ribet\_element\_cramer}}}
\expandafter\newcommand\csname LN@ribetNormaliser\endcsname{\href{https://github.com/lsernau/grh-lean/blob/72fead102f663d5c7a407964a13dc0cf9d977207/grh-compact/src/New/Selmer/RibetCramer.lean\#L79-L80}{{\leanfont ribetNormaliser}}}
\expandafter\newcommand\csname LN@HasTriangleCost\endcsname{\href{https://github.com/lsernau/grh-lean/blob/72fead102f663d5c7a407964a13dc0cf9d977207/grh-compact/src/New/Selmer/RibetFirstOrder.lean\#L20-L24}{{\leanfont HasTriangleCost}}}
\expandafter\newcommand\csname LN@existsmemfittingIdealblock\endcsname{\href{https://github.com/lsernau/grh-lean/blob/72fead102f663d5c7a407964a13dc0cf9d977207/grh-compact/src/New/Selmer/RibetFirstOrder.lean\#L30-L55}{{\leanfont exists\_mem\_fittingIdeal\_block}}}
\expandafter\newcommand\csname LN@IsTraceIdeal\endcsname{\href{https://github.com/lsernau/grh-lean/blob/72fead102f663d5c7a407964a13dc0cf9d977207/grh-compact/src/New/Selmer/RibetFirstOrder.lean\#L67-L71}{{\leanfont IsTraceIdeal}}}
\expandafter\newcommand\csname LN@IsRibetElement\endcsname{\href{https://github.com/lsernau/grh-lean/blob/72fead102f663d5c7a407964a13dc0cf9d977207/grh-compact/src/New/Selmer/RibetFirstOrder.lean\#L75-L77}{{\leanfont IsRibetElement}}}
\expandafter\newcommand\csname LN@ribetfirstorder\endcsname{\href{https://github.com/lsernau/grh-lean/blob/72fead102f663d5c7a407964a13dc0cf9d977207/grh-compact/src/New/Selmer/RibetFirstOrder.lean\#L91-L245}{{\leanfont ribet\_first\_order}}}
\expandafter\newcommand\csname LN@existsrootofunityevallarge\endcsname{\href{https://github.com/lsernau/grh-lean/blob/72fead102f663d5c7a407964a13dc0cf9d977207/grh-compact/src/New/Selmer/RootsOfUnityPoints.lean\#L21-L64}{{\leanfont exists\_root\_of\_unity\_eval\_large}}}
\expandafter\newcommand\csname LN@existsrootsofunityellValle\endcsname{\href{https://github.com/lsernau/grh-lean/blob/72fead102f663d5c7a407964a13dc0cf9d977207/grh-compact/src/New/Selmer/RootsOfUnityPoints.lean\#L72-L105}{{\leanfont exists\_roots\_of\_unity\_ellVal\_le}}}
\expandafter\newcommand\csname LN@isTameTypeofroots\endcsname{\href{https://github.com/lsernau/grh-lean/blob/72fead102f663d5c7a407964a13dc0cf9d977207/grh-compact/src/New/Selmer/RowBalanced.lean\#L22-L32}{{\leanfont isTameType\_of\_roots}}}
\expandafter\newcommand\csname LN@eqoneofmeminertiaAbove\endcsname{\href{https://github.com/lsernau/grh-lean/blob/72fead102f663d5c7a407964a13dc0cf9d977207/grh-compact/src/New/Selmer/RowBalanced.lean\#L36-L38}{{\leanfont IsColumn.\allowbreak{}eq\_one\_of\_mem\_inertiaAbove}}}
\expandafter\newcommand\csname LN@inClassofisLatticeIn\endcsname{\href{https://github.com/lsernau/grh-lean/blob/72fead102f663d5c7a407964a13dc0cf9d977207/grh-compact/src/New/Selmer/RowBalanced.lean\#L43-L66}{{\leanfont inClass\_of\_isLatticeIn}}}
\expandafter\newcommand\csname LN@relLengthobeyingCocyclesle\endcsname{\href{https://github.com/lsernau/grh-lean/blob/72fead102f663d5c7a407964a13dc0cf9d977207/grh-compact/src/New/Selmer/RowBalanced.lean\#L73-L115}{{\leanfont relLength\_obeyingCocycles\_le}}}
\expandafter\newcommand\csname LN@eventuallyisBalanced\endcsname{\href{https://github.com/lsernau/grh-lean/blob/72fead102f663d5c7a407964a13dc0cf9d977207/grh-compact/src/New/Selmer/RowBalanced.lean\#L124-L175}{{\leanfont IsRow.\allowbreak{}eventually\_isBalanced}}}
\expandafter\newcommand\csname LN@addconjGalmemevenClasses\endcsname{\href{https://github.com/lsernau/grh-lean/blob/72fead102f663d5c7a407964a13dc0cf9d977207/grh-compact/src/New/Selmer/TateTwistSteps.lean\#L17-L23}{{\leanfont add\_conjGal\_mem\_evenClasses}}}
\expandafter\newcommand\csname LN@hasTriangleCosttateTwist\endcsname{\href{https://github.com/lsernau/grh-lean/blob/72fead102f663d5c7a407964a13dc0cf9d977207/grh-compact/src/New/Selmer/TateTwistSteps.lean\#L28-L60}{{\leanfont hasTriangleCost\_tateTwist}}}
\expandafter\newcommand\csname LQ@hasTriangleCosttateTwist\endcsname{\leanquoteline{0}{lemma hasTriangleCost\_tateTwist [Fintype J] [LinearOrder J] \{ν w τ φ t γ h c\} \{e : Fin t\} \{ty : J →o}\leanquoteline{2}{\_\} (tri : IsPrimeTriangle ν w τ φ γ t ty h c) \{Z : Submodule O \_\} (N : Z ≤ selmerModule (tateTwist}\leanquoteline{2}{ℓ E O e) (placesDividing E ℓ ∪ .conjClosure (.range w)) ∅ (QuotSpan (ℓ \textasciicircum{} ν : ↑O))) (H : ∀ c ∈ Z,}\leanquoteline{2}{∀ m, (∀ g, c g = (cycInv ℓ E O g \textasciicircum{} e.1 - 1) • m) → ∀ g, (cycInv ℓ E O γ \textasciicircum{} e.1 - 1) • c g = 0) \{b\}}\leanquoteline{2}{(hb : relLength Z (obeying (crysRels (tateTwist ℓ E O e)) \_) + length O (shaOdd ℓ E O e ν) ≤ b) :}\leanquoteline{2}{HasTriangleCost Z (withConjInv τ) (withConjInv φ) \{j | (ty j).1 = e\}}\leanquoteline{2}{(ℓ \textasciicircum{} h) b}}
\expandafter\newcommand\csname LN@evalConjPair\endcsname{\href{https://github.com/lsernau/grh-lean/blob/72fead102f663d5c7a407964a13dc0cf9d977207/grh-compact/src/New/Selmer/TiltingLemma.lean\#L18-L19}{{\leanfont evalConjPair}}}
\expandafter\newcommand\csname LN@memselmerZofmemselmerModule\endcsname{\href{https://github.com/lsernau/grh-lean/blob/72fead102f663d5c7a407964a13dc0cf9d977207/grh-compact/src/New/Selmer/TiltingLemma.lean\#L21-L25}{{\leanfont mem\_selmerZ1\_of\_mem\_selmerModule}}}
\expandafter\newcommand\csname LN@chartamegen\endcsname{\href{https://github.com/lsernau/grh-lean/blob/72fead102f663d5c7a407964a13dc0cf9d977207/grh-compact/src/New/Selmer/TiltingLemma.lean\#L33-L35}{{\leanfont IsPencil.\allowbreak{}char\_tame\_gen}}}
\expandafter\newcommand\csname LN@pointCharidOrConjInv\endcsname{\href{https://github.com/lsernau/grh-lean/blob/72fead102f663d5c7a407964a13dc0cf9d977207/grh-compact/src/New/Selmer/TiltingLemma.lean\#L40-L44}{{\leanfont IsPencil.\allowbreak{}pointChar\_idOrConjInv}}}
\expandafter\newcommand\csname LN@lengthleofdropties\endcsname{\href{https://github.com/lsernau/grh-lean/blob/72fead102f663d5c7a407964a13dc0cf9d977207/grh-compact/src/New/Selmer/TiltingLemma.lean\#L50-L87}{{\leanfont IsPencil.\allowbreak{}length\_le\_of\_drop\_ties}}}
\expandafter\newcommand\csname LN@lengthleofunseenclass\endcsname{\href{https://github.com/lsernau/grh-lean/blob/72fead102f663d5c7a407964a13dc0cf9d977207/grh-compact/src/New/Selmer/TiltingLemma.lean\#L94-L198}{{\leanfont IsPencil.\allowbreak{}length\_le\_of\_unseen\_class}}}
\expandafter\newcommand\csname LN@tilting\endcsname{\href{https://github.com/lsernau/grh-lean/blob/72fead102f663d5c7a407964a13dc0cf9d977207/grh-compact/src/New/Selmer/TiltingLemma.lean\#L206-L224}{{\leanfont tilting}}}
\expandafter\newcommand\csname LN@existsisPencil\endcsname{\href{https://github.com/lsernau/grh-lean/blob/72fead102f663d5c7a407964a13dc0cf9d977207/grh-compact/src/New/Selmer/TwistTheorem.lean\#L31-L94}{{\leanfont exists\_isPencil}}}
\expandafter\newcommand\csname LN@existsanticyclotomictwist\endcsname{\href{https://github.com/lsernau/grh-lean/blob/72fead102f663d5c7a407964a13dc0cf9d977207/grh-compact/src/New/Selmer/TwistTheorem.lean\#L103-L195}{{\leanfont exists\_anticyclotomic\_twist}}}
\expandafter\newcommand\csname LQ@existsanticyclotomictwist\endcsname{\leanquoteline{0}{lemma exists\_anticyclotomic\_twist \{κ T : Type\} [Finite κ] [Finite T] \{O : CoeffRing ℓ\} \{ρ : ∀ i,}\leanquoteline{2}{\_ →* Matrix \_ (n i) O\} \{μ : κ → T → \_\} \{L : ∀ \_, \_\} \{B : ι → \_\} (hρ : ∀ i, IsBalanced ℓ S (ρ i)}\leanquoteline{2}{(L i) (γ i) (β i) M) \{val : κ → \_\} (hf : ∀ k, IsFamily S₀ (μ k) (val k)) (hB : ∀ i, B i ≠ ⊤) (hβ :}\leanquoteline{2}{∀ i, O.normLength (β i + ∑ b, addVal \_ (ρ i (idOrConjInv b z) - 1).det) ≤ B i) \{c η : ℝ≥0\} (hc :}\leanquoteline{2}{∀ k t, ellVal ((μ k t z).1 - 1) ≤ c) (hη : 0 < η) m : ∃ I ψ, ∀ d ∈ Set.Icc 1 m,}\leanquoteline{2}{(∀ k, ellVal (val k I (ψ \textasciicircum{} d)) ≤ ellVal (val k ∅ 1) + 8 * c + η) ∧ IsAuxTwist S₀ I (ψ \textasciicircum{} d) η ∧}\leanquoteline{2}{∀ i, ∀ᶠ O' in atTop, ∃ hO : O ≤ O', ∃ ψ', (∀ g, (ψ' g).1 = (ψ \textasciicircum{} d) g) ∧}\leanquoteline{2}{O'.normLength (length O' (selmerModule (twistRep (inclusion hO).1 ψ' (ρ i)) (unionWithConj S I)}\leanquoteline{2}{(mapRel (inclusion hO).1 '' (L i ∪ \{.single (γ i) 1\})) (QuotSpan (ℓ \textasciicircum{} M : ↑O')))) ≤ B i + η}}
\newcommand{\ConcordanceRows}{\crow{crudeUnit}{\href{https://github.com/lsernau/grh-lean/blob/72fead102f663d5c7a407964a13dc0cf9d977207/grh-compact/src/New/Selmer/AptTupleDefs.lean\#L27-L27}{{\leanfont crudeUnit}}}{AptTupleDefs.lean}{27--27}
\crow{IsColumn}{\href{https://github.com/lsernau/grh-lean/blob/72fead102f663d5c7a407964a13dc0cf9d977207/grh-compact/src/New/Selmer/AptTupleDefs.lean\#L43-L45}{{\leanfont IsColumn}}}{AptTupleDefs.lean}{43--45}
\crow{IsRow}{\href{https://github.com/lsernau/grh-lean/blob/72fead102f663d5c7a407964a13dc0cf9d977207/grh-compact/src/New/Selmer/AptTupleDefs.lean\#L50-L59}{{\leanfont IsRow}}}{AptTupleDefs.lean}{50--59}
\crow{obeyingCocycles}{\href{https://github.com/lsernau/grh-lean/blob/72fead102f663d5c7a407964a13dc0cf9d977207/grh-compact/src/New/Selmer/BlocksAndCocycles.lean\#L16-L17}{{\leanfont obeyingCocycles}}}{BlocksAndCocycles.lean}{16--17}
\crow{AuxPrimeData}{\href{https://github.com/lsernau/grh-lean/blob/72fead102f663d5c7a407964a13dc0cf9d977207/grh-compact/src/New/Selmer/Corepresentable.lean\#L38-L40}{{\leanfont AuxPrimeData}}}{Corepresentable.lean}{38--40}
\crow{IsAnticyclotomicAt}{\href{https://github.com/lsernau/grh-lean/blob/72fead102f663d5c7a407964a13dc0cf9d977207/grh-compact/src/New/Selmer/Corepresentable.lean\#L46-L51}{{\leanfont IsAnticyclotomicAt}}}{Corepresentable.lean}{46--51}
\crow{IsAuxiliary}{\href{https://github.com/lsernau/grh-lean/blob/72fead102f663d5c7a407964a13dc0cf9d977207/grh-compact/src/New/Selmer/Corepresentable.lean\#L55-L55}{{\leanfont IsAuxiliary}}}{Corepresentable.lean}{55--55}
\crow{IsFamily}{\href{https://github.com/lsernau/grh-lean/blob/72fead102f663d5c7a407964a13dc0cf9d977207/grh-compact/src/New/Selmer/Corepresentable.lean\#L60-L68}{{\leanfont IsFamily}}}{Corepresentable.lean}{60--68}
\crow{evenClasses}{\href{https://github.com/lsernau/grh-lean/blob/72fead102f663d5c7a407964a13dc0cf9d977207/grh-compact/src/New/Selmer/Corepresentable.lean\#L136-L138}{{\leanfont evenClasses}}}{Corepresentable.lean}{136--138}
\crow{IsNewPlaces}{\href{https://github.com/lsernau/grh-lean/blob/72fead102f663d5c7a407964a13dc0cf9d977207/grh-compact/src/New/Selmer/Corepresentable.lean\#L142-L143}{{\leanfont IsNewPlaces}}}{Corepresentable.lean}{142--143}
\crow{withConjInv}{\href{https://github.com/lsernau/grh-lean/blob/72fead102f663d5c7a407964a13dc0cf9d977207/grh-compact/src/New/Selmer/Corepresentable.lean\#L147-L147}{{\leanfont withConjInv}}}{Corepresentable.lean}{147--147}
\crow{IsPrimeTriangle}{\href{https://github.com/lsernau/grh-lean/blob/72fead102f663d5c7a407964a13dc0cf9d977207/grh-compact/src/New/Selmer/Corepresentable.lean\#L152-L162}{{\leanfont IsPrimeTriangle}}}{Corepresentable.lean}{152--162}
\crow{IsAuxTwist}{\href{https://github.com/lsernau/grh-lean/blob/72fead102f663d5c7a407964a13dc0cf9d977207/grh-compact/src/New/Selmer/Corepresentable.lean\#L166-L167}{{\leanfont IsAuxTwist}}}{Corepresentable.lean}{166--167}
\crow{unionWithConj}{\href{https://github.com/lsernau/grh-lean/blob/72fead102f663d5c7a407964a13dc0cf9d977207/grh-compact/src/New/Selmer/Corepresentable.lean\#L206-L206}{{\leanfont unionWithConj}}}{Corepresentable.lean}{206--206}
\crow{pointChar}{\href{https://github.com/lsernau/grh-lean/blob/72fead102f663d5c7a407964a13dc0cf9d977207/grh-compact/src/New/Selmer/Corepresentable.lean\#L210-L211}{{\leanfont pointChar}}}{Corepresentable.lean}{210--211}
\crow{twistAtPoint}{\href{https://github.com/lsernau/grh-lean/blob/72fead102f663d5c7a407964a13dc0cf9d977207/grh-compact/src/New/Selmer/Corepresentable.lean\#L215-L216}{{\leanfont twistAtPoint}}}{Corepresentable.lean}{215--216}
\crow{reduceModPow}{\href{https://github.com/lsernau/grh-lean/blob/72fead102f663d5c7a407964a13dc0cf9d977207/grh-compact/src/New/Selmer/Corepresentable.lean\#L219-L220}{{\leanfont reduceModPow}}}{Corepresentable.lean}{219--220}
\crow{IsBalanced}{\href{https://github.com/lsernau/grh-lean/blob/72fead102f663d5c7a407964a13dc0cf9d977207/grh-compact/src/New/Selmer/Corepresentable.lean\#L227-L238}{{\leanfont IsBalanced}}}{Corepresentable.lean}{227--238}
\crow{IsPencil}{\href{https://github.com/lsernau/grh-lean/blob/72fead102f663d5c7a407964a13dc0cf9d977207/grh-compact/src/New/Selmer/Corepresentable.lean\#L243-L250}{{\leanfont IsPencil}}}{Corepresentable.lean}{243--250}
\crow{tieRelations}{\href{https://github.com/lsernau/grh-lean/blob/72fead102f663d5c7a407964a13dc0cf9d977207/grh-compact/src/New/Selmer/Corepresentable.lean\#L254-L256}{{\leanfont tieRelations}}}{Corepresentable.lean}{254--256}
\crow{matrixLattice}{\href{https://github.com/lsernau/grh-lean/blob/72fead102f663d5c7a407964a13dc0cf9d977207/grh-compact/src/New/Selmer/CrystallineRank.lean\#L22-L23}{{\leanfont matrixLattice}}}{CrystallineRank.lean}{22--23}
\crow{coeffEmb}{\href{https://github.com/lsernau/grh-lean/blob/72fead102f663d5c7a407964a13dc0cf9d977207/grh-compact/src/New/Selmer/DeterminantBound.lean\#L34-L35}{{\leanfont coeffEmb}}}{DeterminantBound.lean}{34--35}
\crow{ZfAt}{\href{https://github.com/lsernau/grh-lean/blob/72fead102f663d5c7a407964a13dc0cf9d977207/grh-compact/src/New/Selmer/DeterminantBound.lean\#L39-L39}{{\leanfont ZfAt}}}{DeterminantBound.lean}{39--39}
\crow{crysRels}{\href{https://github.com/lsernau/grh-lean/blob/72fead102f663d5c7a407964a13dc0cf9d977207/grh-compact/src/New/Selmer/DeterminantBound.lean\#L43-L43}{{\leanfont crysRels}}}{DeterminantBound.lean}{43--43}
\crow{complexEmbOf}{\href{https://github.com/lsernau/grh-lean/blob/72fead102f663d5c7a407964a13dc0cf9d977207/grh-compact/src/New/Selmer/DeterminantBound.lean\#L59-L60}{{\leanfont complexEmbOf}}}{DeterminantBound.lean}{59--60}
\crow{AbsGal}{\href{https://github.com/lsernau/grh-lean/blob/72fead102f663d5c7a407964a13dc0cf9d977207/grh-compact/src/New/Selmer/LevelPlacesChosen.lean\#L25-L25}{{\leanfont AbsGal}}}{LevelPlacesChosen.lean}{25--25}
\crow{placesAboveG}{\href{https://github.com/lsernau/grh-lean/blob/72fead102f663d5c7a407964a13dc0cf9d977207/grh-compact/src/New/Selmer/LevelPlacesChosen.lean\#L28-L28}{{\leanfont placesAboveG}}}{LevelPlacesChosen.lean}{28--28}
\crow{badPlaces}{\href{https://github.com/lsernau/grh-lean/blob/72fead102f663d5c7a407964a13dc0cf9d977207/grh-compact/src/New/Selmer/LevelPlacesChosen.lean\#L32-L32}{{\leanfont badPlaces}}}{LevelPlacesChosen.lean}{32--32}
\crow{cardletorsionofisotropic}{\href{https://github.com/lsernau/grh-lean/blob/72fead102f663d5c7a407964a13dc0cf9d977207/grh-compact/src/New/Selmer/LinePointPowers.lean\#L19-L53}{{\leanfont card\_le\_torsion\_of\_isotropic}}}{LinePointPowers.lean}{19--53}
\crow{existssharpfittingpoly}{\href{https://github.com/lsernau/grh-lean/blob/72fead102f663d5c7a407964a13dc0cf9d977207/grh-compact/src/New/Selmer/LinePointPowers.lean\#L64-L86}{{\leanfont exists\_sharp\_fitting\_poly}}}{LinePointPowers.lean}{64--86}
\crow{dvdoneaddmulpowsubone}{\href{https://github.com/lsernau/grh-lean/blob/72fead102f663d5c7a407964a13dc0cf9d977207/grh-compact/src/New/Selmer/LinePointPowers.lean\#L88-L95}{{\leanfont dvd\_one\_add\_mul\_pow\_sub\_one}}}{LinePointPowers.lean}{88--95}
\crow{blockFlattenAlgEquiv}{\href{https://github.com/lsernau/grh-lean/blob/72fead102f663d5c7a407964a13dc0cf9d977207/grh-compact/src/New/Selmer/RibetCramer.lean\#L16-L18}{{\leanfont blockFlattenAlgEquiv}}}{RibetCramer.lean}{16--18}
\crow{amitsurcharpolymemmatrix}{\href{https://github.com/lsernau/grh-lean/blob/72fead102f663d5c7a407964a13dc0cf9d977207/grh-compact/src/New/Selmer/RibetCramer.lean\#L27-L35}{{\leanfont amitsur\_charpoly\_mem\_matrix}}}{RibetCramer.lean}{27--35}
\crow{existsribetelementcramer}{\href{https://github.com/lsernau/grh-lean/blob/72fead102f663d5c7a407964a13dc0cf9d977207/grh-compact/src/New/Selmer/RibetCramer.lean\#L40-L75}{{\leanfont exists\_ribet\_element\_cramer}}}{RibetCramer.lean}{40--75}
\crow{ribetNormaliser}{\href{https://github.com/lsernau/grh-lean/blob/72fead102f663d5c7a407964a13dc0cf9d977207/grh-compact/src/New/Selmer/RibetCramer.lean\#L79-L80}{{\leanfont ribetNormaliser}}}{RibetCramer.lean}{79--80}
\crow{HasTriangleCost}{\href{https://github.com/lsernau/grh-lean/blob/72fead102f663d5c7a407964a13dc0cf9d977207/grh-compact/src/New/Selmer/RibetFirstOrder.lean\#L20-L24}{{\leanfont HasTriangleCost}}}{RibetFirstOrder.lean}{20--24}
\crow{existsmemfittingIdealblock}{\href{https://github.com/lsernau/grh-lean/blob/72fead102f663d5c7a407964a13dc0cf9d977207/grh-compact/src/New/Selmer/RibetFirstOrder.lean\#L30-L55}{{\leanfont exists\_mem\_fittingIdeal\_block}}}{RibetFirstOrder.lean}{30--55}
\crow{IsTraceIdeal}{\href{https://github.com/lsernau/grh-lean/blob/72fead102f663d5c7a407964a13dc0cf9d977207/grh-compact/src/New/Selmer/RibetFirstOrder.lean\#L67-L71}{{\leanfont IsTraceIdeal}}}{RibetFirstOrder.lean}{67--71}
\crow{IsRibetElement}{\href{https://github.com/lsernau/grh-lean/blob/72fead102f663d5c7a407964a13dc0cf9d977207/grh-compact/src/New/Selmer/RibetFirstOrder.lean\#L75-L77}{{\leanfont IsRibetElement}}}{RibetFirstOrder.lean}{75--77}
\crow{ribetfirstorder}{\href{https://github.com/lsernau/grh-lean/blob/72fead102f663d5c7a407964a13dc0cf9d977207/grh-compact/src/New/Selmer/RibetFirstOrder.lean\#L91-L245}{{\leanfont ribet\_first\_order}}}{RibetFirstOrder.lean}{91--245}
\crow{existsrootofunityevallarge}{\href{https://github.com/lsernau/grh-lean/blob/72fead102f663d5c7a407964a13dc0cf9d977207/grh-compact/src/New/Selmer/RootsOfUnityPoints.lean\#L21-L64}{{\leanfont exists\_root\_of\_unity\_eval\_large}}}{RootsOfUnityPoints.lean}{21--64}
\crow{existsrootsofunityellValle}{\href{https://github.com/lsernau/grh-lean/blob/72fead102f663d5c7a407964a13dc0cf9d977207/grh-compact/src/New/Selmer/RootsOfUnityPoints.lean\#L72-L105}{{\leanfont exists\_roots\_of\_unity\_ellVal\_le}}}{RootsOfUnityPoints.lean}{72--105}
\crow{isTameTypeofroots}{\href{https://github.com/lsernau/grh-lean/blob/72fead102f663d5c7a407964a13dc0cf9d977207/grh-compact/src/New/Selmer/RowBalanced.lean\#L22-L32}{{\leanfont isTameType\_of\_roots}}}{RowBalanced.lean}{22--32}
\crow{eqoneofmeminertiaAbove}{\href{https://github.com/lsernau/grh-lean/blob/72fead102f663d5c7a407964a13dc0cf9d977207/grh-compact/src/New/Selmer/RowBalanced.lean\#L36-L38}{{\leanfont IsColumn.\allowbreak{}eq\_one\_of\_mem\_inertiaAbove}}}{RowBalanced.lean}{36--38}
\crow{inClassofisLatticeIn}{\href{https://github.com/lsernau/grh-lean/blob/72fead102f663d5c7a407964a13dc0cf9d977207/grh-compact/src/New/Selmer/RowBalanced.lean\#L43-L66}{{\leanfont inClass\_of\_isLatticeIn}}}{RowBalanced.lean}{43--66}
\crow{relLengthobeyingCocyclesle}{\href{https://github.com/lsernau/grh-lean/blob/72fead102f663d5c7a407964a13dc0cf9d977207/grh-compact/src/New/Selmer/RowBalanced.lean\#L73-L115}{{\leanfont relLength\_obeyingCocycles\_le}}}{RowBalanced.lean}{73--115}
\crow{eventuallyisBalanced}{\href{https://github.com/lsernau/grh-lean/blob/72fead102f663d5c7a407964a13dc0cf9d977207/grh-compact/src/New/Selmer/RowBalanced.lean\#L124-L175}{{\leanfont IsRow.\allowbreak{}eventually\_isBalanced}}}{RowBalanced.lean}{124--175}
\crow{addconjGalmemevenClasses}{\href{https://github.com/lsernau/grh-lean/blob/72fead102f663d5c7a407964a13dc0cf9d977207/grh-compact/src/New/Selmer/TateTwistSteps.lean\#L17-L23}{{\leanfont add\_conjGal\_mem\_evenClasses}}}{TateTwistSteps.lean}{17--23}
\crow{hasTriangleCosttateTwist}{\href{https://github.com/lsernau/grh-lean/blob/72fead102f663d5c7a407964a13dc0cf9d977207/grh-compact/src/New/Selmer/TateTwistSteps.lean\#L28-L60}{{\leanfont hasTriangleCost\_tateTwist}}}{TateTwistSteps.lean}{28--60}
\crow{evalConjPair}{\href{https://github.com/lsernau/grh-lean/blob/72fead102f663d5c7a407964a13dc0cf9d977207/grh-compact/src/New/Selmer/TiltingLemma.lean\#L18-L19}{{\leanfont evalConjPair}}}{TiltingLemma.lean}{18--19}
\crow{memselmerZofmemselmerModule}{\href{https://github.com/lsernau/grh-lean/blob/72fead102f663d5c7a407964a13dc0cf9d977207/grh-compact/src/New/Selmer/TiltingLemma.lean\#L21-L25}{{\leanfont mem\_selmerZ1\_of\_mem\_selmerModule}}}{TiltingLemma.lean}{21--25}
\crow{chartamegen}{\href{https://github.com/lsernau/grh-lean/blob/72fead102f663d5c7a407964a13dc0cf9d977207/grh-compact/src/New/Selmer/TiltingLemma.lean\#L33-L35}{{\leanfont IsPencil.\allowbreak{}char\_tame\_gen}}}{TiltingLemma.lean}{33--35}
\crow{pointCharidOrConjInv}{\href{https://github.com/lsernau/grh-lean/blob/72fead102f663d5c7a407964a13dc0cf9d977207/grh-compact/src/New/Selmer/TiltingLemma.lean\#L40-L44}{{\leanfont IsPencil.\allowbreak{}pointChar\_idOrConjInv}}}{TiltingLemma.lean}{40--44}
\crow{lengthleofdropties}{\href{https://github.com/lsernau/grh-lean/blob/72fead102f663d5c7a407964a13dc0cf9d977207/grh-compact/src/New/Selmer/TiltingLemma.lean\#L50-L87}{{\leanfont IsPencil.\allowbreak{}length\_le\_of\_drop\_ties}}}{TiltingLemma.lean}{50--87}
\crow{lengthleofunseenclass}{\href{https://github.com/lsernau/grh-lean/blob/72fead102f663d5c7a407964a13dc0cf9d977207/grh-compact/src/New/Selmer/TiltingLemma.lean\#L94-L198}{{\leanfont IsPencil.\allowbreak{}length\_le\_of\_unseen\_class}}}{TiltingLemma.lean}{94--198}
\crow{tilting}{\href{https://github.com/lsernau/grh-lean/blob/72fead102f663d5c7a407964a13dc0cf9d977207/grh-compact/src/New/Selmer/TiltingLemma.lean\#L206-L224}{{\leanfont tilting}}}{TiltingLemma.lean}{206--224}
\crow{existsisPencil}{\href{https://github.com/lsernau/grh-lean/blob/72fead102f663d5c7a407964a13dc0cf9d977207/grh-compact/src/New/Selmer/TwistTheorem.lean\#L31-L94}{{\leanfont exists\_isPencil}}}{TwistTheorem.lean}{31--94}
\crow{existsanticyclotomictwist}{\href{https://github.com/lsernau/grh-lean/blob/72fead102f663d5c7a407964a13dc0cf9d977207/grh-compact/src/New/Selmer/TwistTheorem.lean\#L103-L195}{{\leanfont exists\_anticyclotomic\_twist}}}{TwistTheorem.lean}{103--195}}
% the notations of the defined functions: one slot for each explicit argument of the Lean declaration
\newcommand{\crudeUnit}[2]{\mathrm{cr}_{#1}(#2)}
\newcommand{\IsColumn}[5]{\mathsf{ColRep}_{#1}(#2,#3,#4;\,#5)}
\newcommand{\IsRow}[4]{\mathsf{RowRep}_{#1}(#2,#3;\,#4)}
\newcommand{\obeyingCocycles}[3]{\mathcal{Z}_{#1}(#2;\,#3)}
\newcommand{\AuxPrimeData}[2]{\mathrm{Aux}(#1,#2)}
\newcommand{\IsAnticyclotomicAt}[3]{\mathsf{Acyc}(#1,#2;\,#3)}
\newcommand{\IsAuxiliary}[2]{\mathsf{Auxil}(#1;\,#2)}
\newcommand{\IsFamily}[3]{\mathsf{Fam}(#1;\,#2,#3)}
\newcommand{\evenClasses}[4]{\mathrm{Ev}(#1,#2;\,#3,#4)}
\newcommand{\IsNewPlaces}[4]{\mathsf{New}_{#1}(#2,#3,#4)}
\newcommand{\withConjInv}[2]{#1^{[#2]}}
\newcommand{\IsPrimeTriangle}[9]{\mathsf{Tri}_{#1}(#2,#3,#4;\,#5;\,#6,#7;\,#8,#9)}
\newcommand{\IsAuxTwist}[4]{\mathsf{AuxTw}(#1,#2;\,#3,#4)}
\newcommand{\unionWithConj}[2]{#1\cup #2\cup\overline{#2}}
\newcommand{\twistAtPoint}[5]{\operatorname{tp}_{#1,#2}(#3,#4)(#5)}
\newcommand{\reduceModPow}[3]{#3\bmod #1^{#2}}
\newcommand{\IsBalanced}[7]{\mathsf{Bal}_{#1}(#2;\,#3,#4;\,#5,#6,#7)}
\newcommand{\IsPencil}[8]{\mathsf{Pencil}_{#1}(#2;\,#3,#4;\,#5,#6,#7;\,#8)}
\newcommand{\tieRelations}[3]{\operatorname{Ties}_{#3}(#1,#2)}
\newcommand{\matrixLattice}[1]{\mathrm{L}(#1)}
\newcommand{\coeffEmb}[2]{\operatorname{emb}_{#1}(#2)}
\newcommand{\ZfAt}[2]{Z_{\mathrm{cris}}(#1;\,#2)}
\newcommand{\crysRels}[1]{\mathcal{R}_{\mathrm{cris}}(#1)}
\newcommand{\complexEmbOf}[2]{\operatorname{ce}(#1,#2)}
\newcommand{\AbsGal}[1]{G_{E(#1)}}
\newcommand{\placesAboveG}[1]{\Sigma_{g}(#1)}
\newcommand{\badPlaces}[2]{\Sigma_{\mathrm{bad}}(#1,#2)}
\newcommand{\blockFlattenAlgEquiv}[1]{\operatorname{fl}_{#1}}
\newcommand{\ribetNormaliser}[6]{\operatorname{rn}_{#1,#2}(#3,#4,#5)(#6)}
\newcommand{\HasTriangleCost}[6]{\mathsf{TriCost}(#1;\,#2,#3;\,#4;\,#5,#6)}
\newcommand{\IsTraceIdeal}[7]{\mathsf{TrId}(#1;\,#2;\,#3;\,#4,#5,#6,#7)}
\newcommand{\IsRibetElement}[4]{\mathsf{Rib}(#1;\,#2;\,#3,#4)}
\newcommand{\evalConjPair}[1]{\operatorname{ev}^{\pm}_{#1}}
\newcommand{\UnitMembers}{205}
\newcommand{\UnitInNote}{58}
\newcommand{\UnitProofs}{96}
\newcommand{\UnitSteps}{975}
\newcommand{\UnitAllSteps}{1091}


\linespread{1.08}                       % 11 pt on 14.7 pt
\setlength{\parindent}{0pt}
\setlength{\parskip}{5.5pt plus 1.5pt minus 0.5pt}
\setlength{\abovedisplayskip}{7.5pt plus 2pt minus 3pt}
\setlength{\belowdisplayskip}{7.5pt plus 2pt minus 3pt}
\setlength{\abovedisplayshortskip}{3pt plus 2pt}
\setlength{\belowdisplayshortskip}{6pt plus 2pt minus 2pt}
\emergencystretch=1.5em
\clubpenalty=10000 \widowpenalty=10000 \displaywidowpenalty=10000
\relpenalty=7000 \binoppenalty=9000
\raggedbottom
\pagestyle{fancy}\fancyhf{}\renewcommand{\headrulewidth}{0pt}
\fancyfoot[C]{\sffamily\fontsize{7.6}{9}\selectfont\color{inkfaint}\thepage\,/\,\pageref*{LastPage}}
\setlist{topsep=2pt,partopsep=0pt,parsep=0pt,itemsep=4pt,leftmargin=2.1em,labelsep=0.55em}

\newcommand{\Q}{\mathbb{Q}}
\newcommand{\Z}{\mathbb{Z}}
\newcommand{\F}{\mathbb{F}}
\newcommand{\len}{\operatorname{len}}
\newcommand{\cok}{\operatorname{cok}}
\newcommand{\End}{\operatorname{End}}
\newcommand{\Hom}{\operatorname{Hom}}
\newcommand{\Gal}{\operatorname{Gal}}
\newcommand{\diag}{\operatorname{diag}}
\newcommand{\kw}[1]{\emph{#1}}
\newcommand{\m}[1]{\mbox{$#1$}}        % a short formula that must not be broken across lines
\newcommand{\head}[1]{\par\addvspace{17pt plus 4pt minus 2pt}\Needspace{6\baselineskip}%
  \noindent{\fontsize{13.2}{16}\selectfont\bfseries\color{ink}#1}\par\nopagebreak\vspace{-1pt}}
\newcommand{\sect}[2]{\refstepcounter{section}\label{#2}\par\addvspace{19pt plus 4pt minus 2pt}\Needspace{7\baselineskip}%
  \noindent\llap{{\fontsize{13.2}{16}\selectfont\bfseries\color{inkfaint}\thesection}\hspace{9pt}}%
  {\fontsize{13.2}{16}\selectfont\bfseries\color{ink}#1}\par\nopagebreak\vspace{-1pt}}
\newcommand{\subhead}[1]{\par\addvspace{9pt plus 2pt}\Needspace{4\baselineskip}\noindent\textbf{#1}\enspace\ignorespaces}
% where a statement lives in the Lean archive
\newcommand{\leantag}{{\sffamily\bfseries\fontsize{6.3}{8}\selectfont\color{leancol}\addfontfeatures{LetterSpace=7}LEAN}}
\newcommand{\sep}{\unskip\nobreak\hspace{5pt}{\color{inkfaint}\textperiodcentered}\hspace{5pt plus 2pt}\ignorespaces}
\newcommand{\leanline}[1]{\par\nopagebreak\vspace{-2.5pt}{\fontsize{8.9}{12.4}\selectfont\color{inksoft}\raggedright
  \noindent\llap{\leantag\hspace{9pt}}#1\par}\vspace{1pt}}
% a displayed statement in a tinted box: #1 colour, #2 text
\newcommand{\statement}[2]{\par\addvspace{7pt}\begingroup\setlength{\fboxsep}{8.5pt}\setlength{\fboxrule}{0.6pt}%
  \noindent\fcolorbox{#1!30!white}{#1!3!white}{\parbox{\dimexpr\linewidth-2\fboxsep-2\fboxrule\relax}{#2}}%
  \endgroup\par\addvspace{7pt}}
\newcommand{\nt}[2]{\par\noindent\hangindent=1.1em\hangafter=1 #1\enspace{\color{inksoft}#2}\par\vspace{1pt}}
\newcommand{\legdot}[1]{\tikz[baseline=-0.62ex]\fill[#1] (0,0) circle (2.5pt);\hspace{3pt}}
\newcommand{\crow}[4]{\csname CS@#1\endcsname & #2 & \textcolor{inksoft}{#3} & \textcolor{inksoft}{#4}\\}
\newcommand{\leanquoteline}[2]{\par\noindent\hangindent=\dimexpr#1\fontdimen2\font+2.2em\relax\hangafter=1
  \hspace*{#1\fontdimen2\font}#2}
\newcommand{\quoted}[1]{\par\vspace{3pt}{\ttfamily\fontsize{7.5}{10.4}\selectfont\color{leancol}\raggedright
  \setlength{\parskip}{0pt}\leanquote{#1}\par}\vspace{4pt}}
\newcommand{\bx}[3]{{\sffamily\bfseries\fontsize{7.3}{9.2}\selectfont\color{#1!82!black}#2}\\[1.5pt]#3}
\newcommand{\cs}[2]{\expandafter\newcommand\csname CS@#1\endcsname{#2}}

% theorems
\newtheoremstyle{plainx}{7pt plus 2pt}{7pt plus 2pt}{\itshape}{}{\bfseries}{.}{0.5em}{}
\newtheoremstyle{uprightx}{7pt plus 2pt}{7pt plus 2pt}{}{}{\bfseries}{.}{0.5em}{}
\theoremstyle{plainx}
\newtheorem{theorem}{Theorem}[section]
\newtheorem{lemma}[theorem]{Lemma}
\newtheorem{proposition}[theorem]{Proposition}
\newtheorem{corollary}[theorem]{Corollary}
\theoremstyle{uprightx}
\newtheorem{definition}[theorem]{Definition}
\newtheorem{setting}[theorem]{Setting}          % hypotheses shared by several statements; each of them cites it by number
\newtheorem{remark}[theorem]{Remark}
\newtheorem{example}[theorem]{Example}
\renewcommand{\qedsymbol}{{\setlength{\fboxsep}{0pt}\setlength{\fboxrule}{0.5pt}\fbox{\rule{0pt}{0.46em}\hspace{0.46em}}}}
% what holds throughout the note: the ambient objects only (implicit and instance arguments), never a hypothesis
\newenvironment{conventions}{\par\addvspace{5pt}\noindent\textbf{Conventions.}\enspace\ignorespaces}{\par\addvspace{3pt}}
\makeatletter
\renewenvironment{proof}[1][Proof]{\par\pushQED{\qed}\normalfont\topsep0pt\partopsep0pt
  \trivlist\item[\hskip\labelsep\itshape #1.]\ignorespaces}{\popQED\endtrivlist\@endpefalse}
\makeatother

% ---- macros of the three notes of this unit (background notions; the notions that the unit defines are in data.tex)
\newcommand{\CC}{\mathbb{C}}
\newcommand{\NN}{\mathbb{N}}
\newcommand{\Qb}{\bar{\mathbb{Q}}}
\newcommand{\Ql}{\mathbb{Q}_\ell}
\newcommand{\Zl}{\mathbb{Z}_\ell}
\newcommand{\Qlb}{\bar{\mathbb{Q}}_\ell}
\newcommand{\Zlb}{\bar{\mathbb{Z}}_\ell}
\newcommand{\Cl}{\mathbb{C}_\ell}
\newcommand{\Gl}{G_{\mathbb{Q}_\ell}}
\newcommand{\GE}[1]{G_{#1}}
\newcommand{\cyc}{\operatorname{cyc}}
\newcommand{\vl}{v_\ell}
\newcommand{\Sel}[4]{\operatorname{Sel}_{#2}^{#3}(#1;\,#4)}
\newcommand{\Ob}[2]{\operatorname{Ob}(#1;\,#2)}
\newcommand{\Coc}[1]{\operatorname{Coc}(#1)}
\newcommand{\Rel}[1]{\operatorname{Rel}(#1)}
\newcommand{\pd}[2]{\Sigma_{#2}(#1)}
\newcommand{\inab}[1]{I_{#1}}
\newcommand{\cj}[1]{\bar{#1}}
\newcommand{\Nm}{\operatorname{N}}
\newcommand{\nlen}{\operatorname{nlen}}
\newcommand{\cO}{\mathcal{O}}
\newcommand{\fP}{\mathfrak{P}}
\newcommand{\fq}{\mathfrak{q}}
\newcommand{\HT}{\mathcal{H}}
\newcommand{\GL}{\mathrm{GL}}
\newcommand{\tr}{\operatorname{tr}}
\newcommand{\rk}{\operatorname{rk}}
\newcommand{\Art}{\operatorname{Art}}
\newcommand{\av}{\operatorname{av}}
\newcommand{\Zf}{Z_{\mathrm{f}}}
% the data (as in the note on admissible data)
\newcommand{\Dm}[1]{m(#1)}
\newcommand{\Dr}[1]{r(#1)}
\newcommand{\kD}[1]{k(#1)}
\newcommand{\tD}[1]{t(#1)}
\newcommand{\fD}[1]{f(#1)}
\newcommand{\ED}[1]{E(#1)}
\newcommand{\Dg}[2]{g_{#1}(#2)}
\newcommand{\gD}[1]{g(#1)}
\newcommand{\lvl}[1]{N(#1)}
\newcommand{\hhD}[1]{h_2(#1)}
\newcommand{\hD}[1]{h(#1)}
\newcommand{\bone}[1]{b_1(#1)}
\newcommand{\btwo}[1]{b_2(#1)}
\newcommand{\wt}[3]{w_{#1,#2}(#3)}
\newcommand{\ordch}[1]{\operatorname{ord}_{\mathrm{ch}}(#1)}
\newcommand{\Dmu}[1]{\mu(#1)}
\newcommand{\ND}[1]{N_1(#1)}
\newcommand{\thick}[1]{\operatorname{thick}(#1)}
\newcommand{\inp}[2]{\statement{ana}{\textbf{#1.}\enspace #2}}

\newcommand{\Sha}{\mathrm{Sha}}

% where each declaration of the concordance is treated
\cs{crudeUnit}{Definition~\ref{d:recall}}
\cs{IsColumn}{Definition~\ref{d:recall}}
\cs{IsRow}{Definition~\ref{d:recall}}
\cs{obeyingCocycles}{Definition~\ref{d:recall}}
\cs{AuxPrimeData}{Definition~\ref{d:recall}}
\cs{IsAnticyclotomicAt}{Definition~\ref{d:recall}}
\cs{IsAuxiliary}{Definition~\ref{d:recall}}
\cs{IsFamily}{Definition~\ref{d:recall}}
\cs{evenClasses}{Definition~\ref{d:recall}}
\cs{IsNewPlaces}{Definition~\ref{d:recall}}
\cs{withConjInv}{Definition~\ref{d:recall}}
\cs{IsPrimeTriangle}{Definition~\ref{d:recall}}
\cs{IsAuxTwist}{Definition~\ref{d:recall}}
\cs{unionWithConj}{Definition~\ref{d:recall}}
\cs{pointChar}{Definition~\ref{d:recall}}
\cs{twistAtPoint}{Definition~\ref{d:recall}}
\cs{reduceModPow}{Definition~\ref{d:recall}}
\cs{IsBalanced}{Definition~\ref{d:recall}}
\cs{IsPencil}{Definition~\ref{d:recall}}
\cs{tieRelations}{Definition~\ref{d:recall}}
\cs{matrixLattice}{Definition~\ref{d:recall}}
\cs{coeffEmb}{Definition~\ref{d:recall}}
\cs{ZfAt}{Definition~\ref{d:recall}}
\cs{crysRels}{Definition~\ref{d:recall}}
\cs{complexEmbOf}{Definition~\ref{d:recall}}
\cs{AbsGal}{Definition~\ref{d:recall}}
\cs{placesAboveG}{Definition~\ref{d:recall}}
\cs{badPlaces}{Definition~\ref{d:recall}}
\cs{cardletorsionofisotropic}{Lemma~\ref{l:mirror}}
\cs{existssharpfittingpoly}{Lemma~\ref{l:fitting}}
\cs{dvdoneaddmulpowsubone}{Lemma~\ref{l:linepoint}}
\cs{blockFlattenAlgEquiv}{Definition~\ref{d:flat}}
\cs{amitsurcharpolymemmatrix}{Lemma~\ref{l:amitsur}}
\cs{existsribetelementcramer}{Lemma~\ref{l:cramer}}
\cs{ribetNormaliser}{Definition~\ref{d:normaliser}}
\cs{HasTriangleCost}{Definition~\ref{d:tricost}}
\cs{existsmemfittingIdealblock}{Lemma~\ref{l:block}}
\cs{IsTraceIdeal}{Definition~\ref{d:trid}}
\cs{IsRibetElement}{Definition~\ref{d:rib}}
\cs{ribetfirstorder}{Theorem D}
\cs{existsrootofunityevallarge}{Lemma~\ref{l:rootone}}
\cs{existsrootsofunityellValle}{Lemma~\ref{l:rootmany}}
\cs{isTameTypeofroots}{Lemma~\ref{l:tametype}}
\cs{eqoneofmeminertiaAbove}{Theorem C}
\cs{inClassofisLatticeIn}{Lemma~\ref{l:inclass}}
\cs{relLengthobeyingCocyclesle}{Lemma~\ref{l:obeycup}}
\cs{eventuallyisBalanced}{Theorem C}
\cs{addconjGalmemevenClasses}{Lemma~\ref{l:even}}
\cs{hasTriangleCosttateTwist}{Theorem E}
\cs{evalConjPair}{Definition~\ref{d:evpair}}
\cs{memselmerZofmemselmerModule}{Lemma~\ref{l:selz}}
\cs{chartamegen}{Lemma~\ref{l:pointchar}}
\cs{pointCharidOrConjInv}{Lemma~\ref{l:pointchar}}
\cs{lengthleofdropties}{Lemma~\ref{l:drop}}
\cs{lengthleofunseenclass}{Lemma~\ref{l:unseen}}
\cs{tilting}{Lemma~\ref{l:tilting}}
\cs{existsisPencil}{Theorem B}
\cs{existsanticyclotomictwist}{Theorem A}

\begin{document}
\color{ink}

{\raggedright
{\sffamily\bfseries\fontsize{7.2}{9}\selectfont\color{leancol}\addfontfeatures{LetterSpace=8}READBACK NOTE 32\par}
\vspace{5pt}
{\fontsize{21.5}{25.5}\selectfont\bfseries Short Selmer groups after one anticyclotomic twist, balanced rows, and Ribet's lemma to first order\par}
\vspace{5pt}
{\fontsize{12.4}{16}\selectfont\itshape\color{inksoft} One unit of the Lean artifact \ArtifactName, read back into
mathematics\par}}
\vspace{8pt}
{\color{rulegrey}\hrule height 0.6pt}
\vspace{7pt}
\thispagestyle{empty}

\textbf{Summary.} A good twist must make the Selmer modules of the rows short, for $m$ powers at once, and must
not make the relevant $L$-values too divisible. This note proves that such a twist exists (Theorem~A): given
finitely many \emph{balanced} representations and finitely many families of values, one anticyclotomic character
of $\ell$-power order, ramified at boundedly many new primes, has all its powers \m{\psi,\dots,\psi^{m}} within the
bounds. The primes come from Theorem~B: boundedly many primes \emph{tied} to Frobenius elements make every
Selmer module short; the proof plays a mirror lemma against a tilting lemma for pencils. The passage from ties to
an honest twist is an argument on valuations of polynomials at roots of unity. Theorem~C shows that rows are
balanced. The second half is Ribet's lemma to first order (Theorem~D): from congruences between a block and a
string of characters it produces, by Cramer's rule used twice and Fitting ideals, a Hecke-type element whose top
coefficient has bounded valuation; Theorem~E supplies its hypothesis for the Tate twists from a prime triangle
and the odd part of a Tate--Shafarevich group.

\statement{leancol}{\fontsize{9.6}{13.2}\selectfont\textbf{What this note is.} The artifact \ArtifactName{} is a
body of Lean code: the compact edition of a formal proof, due to L.~Alpöge with the AI system Claude, of the
Riemann hypothesis and of its generalisation to Dirichlet $L$-functions. That headline is a claim. The authors
report that every file compiles and that Lean's kernel accepted the proofs; an independent replay is recorded in
the handout's record of checks when it reports. No mathematician has read the argument through, and nothing here
says that the Riemann hypothesis is proved. This note is the second of three on the unit \emph{The choice of the
twists}: it treats the declarations of eight files (roots of unity, the mirror and tilting lemmas, the twist
theorem, balanced rows, Ribet's lemma), and recalls from the first note, \emph{Columns and rows as Galois
representations}, the definitions it needs; together \UnitInNote{} of the unit's \UnitMembers{} declarations
appear here. The statements are complete. The proofs of Theorems B and~E and of Lemmas \ref{l:rootmany},
\ref{l:linepoint}, \ref{l:fitting}, \ref{l:selz}, \ref{l:pointchar}, \ref{l:drop}, \ref{l:tilting},
\ref{l:tametype}, \ref{l:inclass}, \ref{l:amitsur}, \ref{l:block} and~\ref{l:even} follow the formal proofs. The
proofs of Theorems A, C and~D and of Lemmas \ref{l:rootone}, \ref{l:mirror}, \ref{l:unseen}, \ref{l:obeycup}
and~\ref{l:cramer} are outlines of the formal proofs: they give the steps and where each hypothesis enters, leave
the computations to the formal text, and say so where they stand. Results of the background library and of other units are quoted as Inputs with their
owners. Every statement was compared with its Lean text, argument by argument. The Lean text read is \PinText;
Appendix~A is the concordance of \ConcDecls{} declarations.}

\sect{The statements}{s:intro}
The conventions and the recalled notions are in §\ref{s:back}, the definitions of this note in §\ref{s:defs}.
Throughout, $\ell$ is a prime number.

\statement{arith}{\phantomsection\label{t:A}\textbf{Theorem A} (one anticyclotomic twist makes all Selmer modules short).\enspace
\emph{Let $E$ be a CM field as in the conventions, $\mathcal{I}$ a finite set and \m{(n_i)_{i\in\mathcal{I}}} finite
sets, \m{S\subseteq S_0} finite sets of finite places of~$E$, \m{z\in\inab{S_0\setminus\pd{E}{\ell}}}, and
\m{M\in\NN}. Let $\kappa$ and $T$ be finite sets, $\cO$ a coefficient ring, and for \m{i\in\mathcal{I}} let
\m{\varrho_i\colon\GE{E}\to M_{n_i}(\cO)} be a homomorphism, $L_i$ a set of relations, \m{\gamma_i\in\GE{E}}, and
\m{\beta_i}, $B_i$ elements of \m{[0,\infty]}. Let \m{\mu=(\mu_{k,t})_{k\in\kappa,t\in T}} be homomorphisms
\m{\GE{E}\to\Qlb^{\times}} and \m{(\mathrm{val}_k)_{k\in\kappa}} functions as in the definition of a family. Let
\m{c\ge0} and \m{\eta>0} be real numbers and \m{m\in\NN}. Suppose that}
\begin{enumerate}[label=\textup{(\roman*)},leftmargin=2.3em,itemsep=2pt]
\item \emph{\m{\IsBalanced{\ell}{S}{\varrho_i}{L_i}{\gamma_i}{\beta_i}{M}} for every~$i$, and
\m{\IsFamily{S_0}{\mu_k}{\mathrm{val}_k}} for every~$k$;}
\item \emph{\m{B_i\ne\infty} and
\m{\nlen_{\cO}\bigl(\beta_i+\sum_{b}v_{\cO}(\det(\varrho_i(z^{[b]})-1))\bigr)\le B_i} for every~$i$, the sum over the two
signs;}
\item \emph{\m{\vl(\mu_{k,t}(z)-1)\le c} for all $k$ and~$t$.}
\end{enumerate}
\emph{Then there are a finite set $I$ of finite places of~$E$ and a multiplicative homomorphism
\m{\psi\colon\GE{E}\to\Qlb} such that for every \m{1\le d\le m}:}
\begin{enumerate}[label=\textup{(\arabic*)},leftmargin=2.3em,itemsep=2pt]
\item \emph{\m{\vl(\mathrm{val}_k(I,\psi^{d}))\le\vl(\mathrm{val}_k(\emptyset,1))+8c+\eta} for every~$k$;}
\item \emph{\m{\IsAuxTwist{S_0}{I}{\psi^{d}}{\eta}};}
\item \emph{for every $i$ and all large coefficient rings $\cO'$: \m{\cO\subseteq\cO'}, and there is a
multiplicative homomorphism \m{\psi'\colon\GE{E}\to\cO'} with \m{\psi'(g)=\psi^{d}(g)} for all~$g$ and}
\[
\nlen_{\cO'}\,\Sel{\psi'\cdot\varrho_i}{\unionWithConj{S}{I}}{L_i\cup\{\delta_{\gamma_i}\}}{\cO'/\ell^{M}}\ \le\ B_i+\eta .
\]
\end{enumerate}}
\leanline{\lean{existsanticyclotomictwist}}

\statement{arith}{\phantomsection\label{t:B}\textbf{Theorem B} (tied conditions at new primes shrink Selmer modules).\enspace
\emph{Let $E$, $\mathcal{I}$, \m{(n_i)_i}, \m{S\subseteq S_0}, \m{z\in\inab{S_0\setminus\pd{E}{\ell}}} and \m{M\in\NN} be
as in Theorem~A. Let $A$ be the ring of integers of a finite extension of~\m{\Ql}, and for \m{i\in\mathcal{I}} let
\m{\varrho_i\colon\GE{E}\to M_{n_i}(A)} be a homomorphism, $L_i$ a set of relations, \m{\gamma_i\in\GE{E}} and
\m{\beta_i\in[0,\infty]} with \m{\IsBalanced{\ell}{S}{\varrho_i}{L_i}{\gamma_i}{\beta_i}{M}}. Let $W$ be a neighbourhood
of~$z$ in~\m{\GE{E}} and \m{a\in A} with \m{a\mid\ell^{M}}. For \m{i\in\mathcal{I}} let $U_i$ be the property
\m{\reduceModPow{\ell}{M}{\varrho_i}(g)=1} of \m{g\in\GE{E}}. Then there is \m{J_0\in\NN} such that for every
\m{N\in\NN} there are a finite place~$\fq_0$, a finite set $J$ of finite places with \m{\#J\le J_0}, and a
function $x$ with \m{\IsPencil{\ell}{S_0}{U}{W}{N}{\fq_0}{J}{x}} and}
\[
\len_A\,\Sel{\varrho_i}{\unionWithConj{S}{J}}{L_i\cup\tieRelations{J}{x}{i}\cup\{\delta_{\gamma_i}\}}{A/a}\ \le\ \beta_i\qquad\text{\emph{for all }}i\in\mathcal{I}.
\]}
\leanline{\lean{existsisPencil}}

\statement{arith}{\phantomsection\label{t:C}\textbf{Theorem C} (rows are balanced).\enspace
\emph{Let $P$ be parameters, $D$ admissible data for~$P$, \m{\iota\colon\CC\to\Qlb} the fixed isomorphism, $n$ a
finite set, $t$ a natural number, $\HT$ a type and \m{R\colon\AbsGal{D}\to\GL_n(\Qlb)} a homomorphism with
\m{\IsRow{D}{t}{\HT}{R}}. Let $w$ be a finite place of~\m{\ED{D}} that is chosen and lies in
\m{\placesAboveG{D}\setminus\pd{\ED{D}}{\ell}}, and \m{M\in\NN}. The set
\m{\Sigma=\pd{\ED{D}}{\ell}\cup\placesAboveG{D}} is finite. For all large coefficient rings~$\cO$ there are a
homomorphism \m{u\colon\AbsGal{D}\to M_n(\cO)}, an element \m{\gamma\in\AbsGal{D}} and \m{\beta\in[0,\infty]} such
that $u$ is a lattice of~$R$, \m{\IsBalanced{\ell}{\Sigma}{u}{\crysRels{u}}{\gamma}{\beta}{M}}, and}
\[
\nlen_{\cO}\Bigl(\beta+\sum_{b}v_{\cO}\bigl(\det(u(\tau_w^{[b]})-1)\bigr)\Bigr)\ \le\ 10\,\fD{D}\cdot\#n\cdot\crudeUnit{\ell}{D}.
\]}
\leanline{\lean{eventuallyisBalanced}}

\statement{arith}{\phantomsection\label{t:D}\textbf{Theorem D} (Ribet's lemma to first order).\enspace
\emph{Let $A$ be a discrete valuation ring with valuation~$v_A$, $\mathbb{L}$ a commutative $A$-algebra, $Y$,
$\sigma$, $W$ and $r$ finite sets, $J$ a finite totally ordered set, $\Gamma$ a group, \m{u=(u_x)_{x\in Y}} a family of
homomorphisms \m{u_x\colon\Gamma\to M_\sigma(A)}, $\mathcal{S}$ an $A$-subalgebra of the algebra of functions
\m{Y\to\mathbb{L}[X_j:j\in J]} and $\mathcal{J}$ an $\mathcal{S}$-submodule of that algebra. Let \m{t\ge1},
\m{\epsilon\colon r\sqcup\{0,\dots,t-1\}\to\sigma} a bijection, \m{\mathcal{X}\subseteq Y}, \m{\pi\in\mathcal{X}}, and for
\m{x\in\mathcal{X}} let \m{\varrho_x\colon\Gamma\to M_r(A)} be a homomorphism, \m{\gamma_x\in\Gamma} and \m{y_x\in A}; let
\m{\xi_e\colon\Gamma\to A}, \m{e<t}, be multiplicative homomorphisms. Let \m{p,p_0,C\in A}, let
\m{\tau,\phi\colon J\times\{+,-\}\to\Gamma} be functions, \m{\varepsilon=(\varepsilon_j)_{j\in J}} polynomials in~\m{A[X]},
\m{a\in\Gamma}, \m{\mathrm{ty}\colon J\to\{0,\dots,t-1\}} a monotone map, $\theta$ a finitely supported function
\m{\Gamma\to\NN}, \m{\beta\colon Y\to\mathbb{L}} a function, \m{b_0} and \m{(b_e)_{e<t}} elements of \m{\NN\cup\{\infty\}},
\m{(\delta_v)_{v\in W}} elements of~$\Gamma$, \m{(o_v)_{v\in W}} polynomials in~\m{A[X]}, and $s$ a finite subset
of~$\Gamma$. Suppose:}
\begin{enumerate}[label=\textup{(\alph*)},leftmargin=2.3em,itemsep=2pt]
\item \emph{\m{\xi_0=1}; for \m{x\in\mathcal{X}} and \m{g\in\Gamma}, \m{u_x(g)} is the block sum of \m{\varrho_x(g)} and the
diagonal matrix of the \m{\xi_e(g)}, transported by~$\epsilon$; and the $A$-span of \m{\varrho_\pi(\Gamma)} is \m{M_r(A)};}
\item \emph{\m{\varrho_x(\gamma_x)=y_x\cdot1}, \m{y_x-1\in A^{\times}} and \m{\xi_e(\gamma_x)=1} for all
\m{x\in\mathcal{X}} and \m{e<t};}
\item \emph{\m{p\ne0}; \m{\IsTraceIdeal{\mathcal{J}}{u}{p}{\tau}{\phi}{\mathrm{pr}_1}{\varepsilon}}, with \m{\mathrm{pr}_1} the
first projection of \m{J\times\{+,-\}}; and $p$ divides \m{\varepsilon_{j}(\xi_e(\phi_k))-1} for all $e$ and all
\m{k=(j,b)};}
\item \emph{for every \m{e<t}, the product \m{\prod_{k<t,\,k\ne e}(\xi_e(a)-\xi_k(a))} divides~$C$;}
\item \emph{\m{p_0\ne0}; \m{\beta(x)} is the image in~$\mathbb{L}$ of
\m{\det\bigl(\sum_{g}\ribetNormaliser{\gamma_\pi}{a}{y_\pi}{p_0C}{\theta}{g}\,(u_x(g)-1)\bigr)} for every \m{x\in Y};
\m{b_0<v_A(p_0)}; and the $A$-module of the \m{c\in\obeyingCocycles{u}{\varrho_\pi}{p_0}} that vanish at all
\m{\tau_k} and at~\m{\gamma_\pi} has length at most~$b_0$;}
\item \emph{\m{\xi_e(\delta_v)=1} for all $v$, $e$, and \m{(o_v\cdot(X-1))(u_x(\delta_v))=0} for all \m{v\in W},
\m{x\in Y};}
\item \emph{for every \m{e<t} and every $A$-submodule $Z$ of the functions \m{\Gamma\to A/p}: if every \m{c\in Z}
lies in \m{\obeyingCocycles{u}{\xi_e}{p}}, vanishes at all~$\delta_v$, and satisfies ``for every \m{m\in A/p}, if
\m{c(g)=(\xi_e(g)-1)m} for all~$g$, then \m{(\xi_e(a)-1)\,c(g)=0} for all~$g$'', then
\m{\HasTriangleCost{Z}{\tau}{\phi}{\{j:\mathrm{ty}(j)=e\}}{C}{b_e}}, and every \m{c\in Z} that vanishes on~$s$ is zero.}
\end{enumerate}
\emph{Then there is \m{c\in A} such that \m{\IsRibetElement{\mathcal{J}}{\beta}{\pi}{c}} (with $c$ read
in~$\mathbb{L}$) and}
\[
v_A(c)\ \le\ b_0+\sum_{e<t}b_e+t\Bigl(\sum_{v\in W}v_A(o_v(1))+2\,\#s\cdot v_A(C)\Bigr)+(\#J+1)\,v_A(C).
\]}
\leanline{\lean{ribetfirstorder}}

\statement{arith}{\phantomsection\label{t:E}\textbf{Theorem E} (the cost of a Tate twist).\enspace
\emph{Let $E$ be a CM field as in the conventions, $J$ a finite totally ordered set, $\cO$ a coefficient ring,
\m{\nu,t,h,c} natural numbers, \m{w=(w_j)_{j\in J}} finite places of~$E$, \m{\tau=(\tau_j)_j} and
\m{\phi=(\phi_j)_j} families in~\m{\GE{E}}, \m{\gamma\in\GE{E}}, \m{e<t}, and
\m{\mathrm{ty}\colon J\to\{e':0<e'<t\}} a monotone map with
\m{\IsPrimeTriangle{\nu}{w}{\tau}{\phi}{\gamma}{t}{\mathrm{ty}}{h}{c}}. Let $Z$ be an $\cO$-submodule of the functions
\m{\GE{E}\to\cO/\ell^{\nu}} and \m{b\in\NN\cup\{\infty\}}. Suppose that}
\begin{enumerate}[label=\textup{(\roman*)},leftmargin=2.3em,itemsep=2pt]
\item \emph{\m{Z\subseteq\Sel{\cO(-e)}{\pd{E}{\ell}\cup w(J)\cup\cj{w(J)}}{\emptyset}{\cO/\ell^{\nu}}};}
\item \emph{for every \m{z\in Z} and \m{m\in\cO/\ell^{\nu}}: if \m{z(g)=(\cyc(g)^{-e}-1)\,m} for all~$g$, then
\m{(\cyc(\gamma)^{-e}-1)\,z(g)=0} for all~$g$;}
\item \emph{\m{\len_{\cO}\bigl(Z/Z\cap\Ob{\crysRels{\cO(-e)}}{\cO/\ell^{\nu}}\bigr)+\len_{\cO}\Sha^{\mathrm{odd}}(\ell,E,\cO,e,\nu)\le b}.}
\end{enumerate}
\emph{Then \m{\HasTriangleCost{Z}{\withConjInv{\tau}{\cdot}}{\withConjInv{\phi}{\cdot}}{\{j:\mathrm{ty}(j)=e\}}{\ell^{h}}{b}}.}}
\leanline{\lean{hasTriangleCosttateTwist}}

In Theorem~A the representations $\varrho_i$ are, in the third note, lattices of the rows of the auxiliary pairs,
balanced by Theorem~C with \m{L_i} the crystalline relations, and the families are those of Theorem~D of the
first note; the conclusion (3) is then the Selmer length bound of the first note, and (1) bounds the inverse
ratios of $L$-values. In Theorem~D the points of~$Y$ are Hecke eigensystems, $\mathcal{X}$ the lifts, the
variables $X_j$ the tie parameters at the auxiliary primes, and the element $f$ of
\m{\IsRibetElement{\mathcal{J}}{\beta}{\pi}{c}} a Hecke element that kills a determinant vanishing on the lifts
modulo the ties; its coefficient $c$ at \m{\prod_jX_j} measures the congruence.

\sect{Conventions and recalled notions}{s:back}
\begin{conventions}\phantomsection\label{c:conventions}
The conventions are those of the first note of the unit. The natural numbers are \m{0,1,2,\dots}. $\ell$ is a
prime number; \m{\Ql}, \m{\Zl}, \m{\Qlb}, \m{\Zlb}, \m{\Cl}, \m{\Gl=\Gal(\Qlb/\Ql)} are as usual, with
\m{|\ell|=1/\ell}; \m{\iota\colon\CC\to\Qlb} is a fixed isomorphism of fields where complex numbers occur. \m{\Qb}
is a fixed algebraic closure of~$\mathbb{Q}$ with a fixed complex conjugation~$\mathsf{c}$. A \emph{CM field as
in the conventions} is a CM number field $E$ with a fixed embedding in~\m{\Qb} that contains a primitive root of
unity of order~$5040$; \m{\GE{E}=\Gal(\Qb/E)}; in Lemmas \ref{l:fitting} and~\ref{l:tilting} only a number field with a
fixed embedding is needed, or no field at all, as said there. $P$ denotes parameters and $D$ admissible data
for~$P$. The letters $n$, $r$, $m$, $d$, $\sigma$, $\kappa$, $T$, $N$, $X$, $Y$, $W$, $J$, $\mathcal{I}$, when they
index matrices or families, denote finite sets (in §\ref{s:ribet} some of them are arbitrary sets, as said there); \m{\#n} is the number of elements, \m{M_n(A)} the matrices over a commutative
ring indexed by~$n$. A \emph{coefficient ring} is the ring of integers $\cO$ (also $\cO'$) of a finite Galois
extension of~\m{\Ql} inside~\m{\Qlb}; ``for all large~$\cO$'' means: for all coefficient rings containing a
suitable one; ``for all large~$N$'' is meant in the same way for natural numbers. $A$, $A_0$, $A'$, $\Lambda$, $\mathbb{L}$ denote commutative rings, $\Gamma$ a group,
$B$ a module. A \emph{discrete valuation ring} is a domain; \m{v_A} is its valuation with \m{v_A(\varpi)=1}, with
values in \m{\NN\cup\{\infty\}}.
\end{conventions}

\subhead{2.1. Background.} The background notions are those of §2 of the first note: finite places $w$, $v$,
$\fq$ of~$E$, their conjugates \m{\cj{w}}, norms \m{\Nm(w)} and the sets \m{\pd{E}{M}}; places $\fP$ of~\m{\Qb},
decomposition and inertia groups \m{G_{\fP}}, \m{I_{\fP}}, the inertia \m{\inab{S}} above a set of places, $q$-Frobenius
elements, the tame generators \m{\tau_{\fP}} and \m{\tau_w}; \m{g^{\mathsf c}=\mathsf{c}g\mathsf{c}}, \m{g^{[+]}=g},
\m{g^{[-]}=(g^{\mathsf c})^{-1}}; the cyclotomic character \m{\cyc} and \m{\cyc_M}; characters \emph{from
$\mathrm{U}(1)$ relative to~$S$ with modulus~$I$}; relations $\theta$, the indicator relation $\delta_g$, the
module \m{\Ob{L}{B}} of the functions that obey a set $L$ of relations, the relations \m{\Rel{Z}} of a module of
functions, the cocycle relations, the Selmer module \m{\Sel{\varrho}{S}{L}{B}} (locally constant cocycles
for~$\varrho$ with values in~$B^n$, unramified outside~$S$, obeying~$L$) and the local cocycles; \m{A/a}; the length
\m{\len_A}; the normalised length \m{\nlen_{\cO}=\vl(\varpi)\cdot\len_{\cO}}, applied to a module or to an element
of \m{\NN\cup\{\infty\}}; the truncated valuation \m{\vl} on~\m{\Qlb}; the Tate twist \m{\cO(-e)}; lattices of a
representation \m{R\colon\GE{E}\to\GL_n(\Qlb)}; Galois lattices over~$\cO$ for a group over~\m{\Gl}, their
crystalline cocycles \m{\Zf(L)}, jumps, twisted duals \m{L^{\vee}(1)}, the classes \m{\mathcal{Q}(F,K,U)}, tame types
and cup-orthogonals \m{L^{\perp\mathcal{Q}}}. For a representation $\varrho$ and a multiplicative homomorphism
$\psi'$ with values in the same ring, \m{\psi'\cdot\varrho} is \m{g\mapsto\psi'(g)\varrho(g)}; when $\varrho$ has
coefficients in a subring, it is read in the larger ring. For \m{e,\nu\in\NN}, \m{\Sha^{\mathrm{odd}}(\ell,E,\cO,e,\nu)}
is the module of the cocycles in \m{\Sel{\cO(-e)}{\pd{E}{\ell}}{\emptyset}{\cO/\ell^{\nu}}} that vanish on the inertia groups
above~$\ell$ when \m{e=0}, and, when \m{e\ne0}, vanish at the elements of the decomposition groups above~$\ell$ with
\m{\cyc=1} and satisfy \m{z(g^{\mathsf c})=-z(g)} (background library; the name ``odd'' is apt for \m{e\ne0} only).
For a submodule \m{\mathcal{N}} of a free module $\Lambda^{s}$ of finite rank, its \emph{Fitting ideal}
\m{\operatorname{Fitt}(\mathcal{N})} is the ideal generated by the determinants of the square matrices whose rows
lie in~\m{\mathcal{N}} (the zeroth Fitting ideal of the quotient).

\begin{definition}[notions recalled from the first note]\label{d:recall}
The following are defined in the note \emph{Columns and rows as Galois representations}; the numbers are those
of its definitions.
\begin{enumerate}[label=\textup{(\arabic*)},leftmargin=2.3em,itemsep=2pt]
\item (3.1) For a family \m{u=(u_x)_x} of homomorphisms \m{\Gamma\to M_\sigma(A)}, a homomorphism
\m{\varrho\colon\Gamma\to M_d(A)} and \m{a\in A}, \m{\obeyingCocycles{u}{\varrho}{a}} is the module of the cocycles
\m{c\colon\Gamma\to(A/a)^d} for~$\varrho$ with \m{\sum_g\theta(g)c(g)=0} for every finitely supported
\m{\theta\colon\Gamma\to A} such that \m{\sum_g\theta(g)(u_x(g)-1)=0} for all~$x$.
\item (3.4, 3.5, 3.37) For \m{g\in\Gl}: \m{\coeffEmb{\cO}{g}} is the map \m{\cO\to\Cl}, \m{x\mapsto g(x)}, and
\m{\complexEmbOf{e}{g}} is \m{\iota^{-1}g^{-1}e}, for a ring homomorphism $e$ into~\m{\Qlb}. For a homomorphism
\m{u\colon\GE{E}\to M_n(\cO)} and \m{v\in\pd{E}{\ell}}, \m{\ZfAt{v}{u}} is the module of the crystalline cocycles of the
restriction of~$u$ to the decomposition group of the chosen place over~$v$, and \m{\crysRels{u}} is the set of the
crystalline relations of~$u$: the relations of the modules \m{\ZfAt{v}{u}}, extended by zero. For a homomorphism
\m{\varrho\colon\Gamma\to M_n(A)} into matrices over a \m{\Zl}-algebra, \m{\matrixLattice{\varrho}} is $A^n$ with
$\Gamma$ acting through~$\varrho$.
\item (3.10, 3.13) For admissible data~$D$: \m{\AbsGal{D}=\Gal(\Qb/\ED{D})}; \m{\placesAboveG{D}} is the set of the finite
places of~\m{\ED{D}} above the two primes \m{\Dg{0}{D}}, \m{\Dg{1}{D}}; \m{\badPlaces{D}{\ell}} the set of those
containing \m{\ell\gD{D}\lvl{D}}; and \m{\crudeUnit{\ell}{D}=v_\ell(\ordch{D})+2^{-\ell}}.
\item (3.15, 3.16) \m{\IsColumn{D}{t}{\HT}{S}{R}}: $R$ is continuous with non-degenerate trace form,
\m{\tr R(g^{\mathsf c})=\cyc(g)^{-(t-1)}\tr R(g^{-1})}, and at the places outside
\m{S\setminus(\pd{\ED{D}}{\ell}\cup\placesAboveG{D})} it is of finite type with the jumps~$\HT$, and good outside
\m{\placesAboveG{D}}. \m{\IsRow{D}{t}{\HT}{R}}: a column of level~$\emptyset$ of row type, with the clauses (W2)--(W4)
on the inertia above \m{\placesAboveG{D}}, on the lattices (they span all matrices and contain a scalar \m{u(g)}
with \m{\cyc(g)=\Dg{i}{D}^{\#n}}), and on a scalar inertial element of large residual order.
\item (3.23--3.25) \m{\AuxPrimeData{E}{q}} is the set of the data \m{l=(\fP,\gamma,\phi,h)} at an auxiliary prime, with
\m{h\colon\GE{E}\to\mathbb{Z}/q}. \m{\IsAnticyclotomicAt{T}{\fq}{l}}: $h$ is from $\mathrm{U}(1)$ relative to~$T$ with
modulus~\m{\{\fq\}}, \m{\Nm(\fq)\equiv1\bmod q}, $\fP$ lies over~$\fq$, $\phi$ is an \m{\Nm(\fq)}-Frobenius at~$\fP$,
\m{\fq\notin T}, and \m{\gamma=\tau_{\fP}} with \m{h(\gamma)=1}. \m{\IsAuxiliary{S}{I}}: every \m{v\in I} has prime norm,
and \m{v,\cj{v}\notin S}, \m{\cj{v}\notin I}.
\item (3.26) \m{\IsFamily{S}{\nu}{\mathrm{val}}}: the $\nu_i$ are continuous; for auxiliary~$I$ the value
\m{\mathrm{val}(I,\chi)} is a sum \m{\sum_ac_a\chi(g_a)} with integral~$c_a$; adding a place $\fq$ to~$I$ raises
\m{\vl(\mathrm{val}(I,\chi))} by at most~$8a$ when, at a Frobenius element $\phi$ of~$\fq$, all
\m{\vl(1-\nu_i(\phi)\chi(\phi)^{\pm1})} are at most~$a$; and \m{\mathrm{val}(\emptyset,1)\ne0}.
\item (3.27, 3.28) \m{\evenClasses{\cO}{S}{\nu}{e}}: the cocycles in
\m{\Sel{\cO(-e)}{\pd{E}{\ell}\cup S\cup\cj{S}}{\emptyset}{\cO/\ell^\nu}} with \m{c(g^{\mathsf c})=c(g)} that vanish at the
elements of the decomposition groups above~$\ell$ with \m{\cyc=1}. \m{\IsNewPlaces{\nu}{w}{\tau}{\phi}}: no $w_j$ is
above~$\ell$, \m{\cyc_{\ell^\nu}(\phi_j)=1}, and $\tau_j$ is the tame generator at a place over~$w_j$.
\m{\withConjInv{\tau}{k}=\tau_j^{[b]}} for \m{k=(j,b)}.
\item (3.29) \m{\IsPrimeTriangle{\nu}{w}{\tau}{\phi}{a}{t}{\mathrm{ty}}{h}{c}}: the clauses (T1)--(T5) of the first
note; (T1) is \m{\IsNewPlaces{\nu}{w}{\tau}{\phi}}, (T4) the step ``\m{\ell^{h}z(\tau_j)=0} for an even class
of type \m{\mathrm{ty}(j)} that vanishes at the earlier $\tau_i$, and at the earlier $\phi_i$ of the same type'',
(T5) the clause ``\m{4\delta(a)^2\mid\ell^{h}} and \m{\delta(a)\cdot2z(g)=\delta(g)\cdot2z(a)} for even classes
unramified outside~$\ell$ that vanish at the $\phi_j$ of type~$e$'', with \m{\delta(g)=\cyc(g)^{-e}-1}.
\item (3.30, 3.31) \m{\IsAuxTwist{S}{Q}{\psi}{c}}: \m{\IsAuxiliary{S}{Q}}, $\psi$ is from $\mathrm{U}(1)$ relative to
\m{\pd{E}{\ell}} with modulus~$Q$, no place of~$Q$ is above~$\ell$, and \m{\sum_{\fq\in Q}\vl(\psi(\sigma_{\fq})-1)\le c}
for tame generators \m{\sigma_{\fq}} at places over the \m{\fq\in Q}. \m{\unionWithConj{S}{J}} is the union of $S$,
$J$ and the conjugates of~$J$.
\item (3.32) For homomorphisms \m{h_j\colon\Gamma\to\mathbb{Z}/q} and \m{y_j\in A} (\m{j\in\sigma}) with \m{y_j^{q}=1},
the \emph{point character} is \m{g\mapsto\prod_jy_j^{h_j(g)}}, written \m{y^{h}}; for a ring homomorphism
\m{e\colon A'\to A} and \m{\varrho\colon\Gamma\to M_n(A')},
\m{\twistAtPoint{h}{y}{e}{\varrho}{g}=\prod_jy_j^{h_j(g)}\cdot e(\varrho(g))}; and
\m{\reduceModPow{l}{M}{\varrho}} is $\varrho$ modulo~$l^M$.
\item (3.33) \m{\IsBalanced{\ell}{S}{\varrho}{L}{\gamma}{\beta}{M}}: the clauses (B1)--(B5) of the first note: $\varrho$
modulo~$\ell^M$ is locally constant and unramified outside~$S$; it spans all matrices and the residue field has
more than \m{4\#n+2} elements; \m{\gamma\in\inab{S\setminus\pd{E}{\ell}}} and \m{\gamma^{[-]}} act by one scalar $s$ with
\m{s\pm1} units; coboundaries obey~$L$; and for some $g$ with \m{\varrho(g)=\alpha}, \m{\varrho(g^{[-]})=\alpha'} scalars
and every \m{a\mid\ell^M}: the sum over \m{v\in S} of the least lengths of the local cocycles modulo~$a$, plus
\m{\#n\cdot v_A(\alpha-\cyc(g^{\mathsf c})\alpha')}, is at most \m{(\#S+f_E)\len_A((A/a)^n)+\beta}.
\item (3.34, 3.35) \m{\IsPencil{\ell}{S_0}{U}{W}{N}{\fq_0}{J}{x}}: $x$ assigns to each place $v$ a datum
\m{l_v=(\fP_v,\gamma_v,\phi_v,h_v)\in\AuxPrimeData{E}{\ell^N}} and slopes \m{(s_{v,i})_{i}}; (P1) \m{\IsAuxiliary{S_0}{J}}
and \m{\fq_0\in J}; (P2) \m{\IsAnticyclotomicAt{\pd{E}{\ell}}{v}{l_v}} for \m{v\in J}; (P3) for \m{v\ne\fq_0} in~$J$,
\m{h_{v'}(\phi_v)} is $1$ for \m{v'=\fq_0} and $0$ for the other \m{v'\in J}; (P4) \m{\phi_{\fq_0}\in W}; (P5) $U_i$ holds
at \m{\phi_v^{[\pm]}} when \m{s_{v,i}\ne0}. \m{\tieRelations{J}{x}{i}} is the set of the relations
\m{\delta_{\gamma_v^{[b]}}-s_{v,i}\delta_{\phi_v^{[b]}}}, \m{v\in J}.
\end{enumerate}
\end{definition}
\leanline{\lean{obeyingCocycles}\sep \lean{coeffEmb}\sep \lean{complexEmbOf}\sep \lean{ZfAt}\sep \lean{crysRels}\sep \lean{matrixLattice}\sep
\lean{AbsGal}\sep \lean{placesAboveG}\sep \lean{badPlaces}\sep \lean{crudeUnit}\sep \lean{IsColumn}\sep \lean{IsRow}\sep
\lean{AuxPrimeData}\sep \lean{IsAnticyclotomicAt}\sep \lean{IsAuxiliary}\sep \lean{IsFamily}\sep \lean{evenClasses}\sep
\lean{IsNewPlaces}\sep \lean{withConjInv}\sep \lean{IsPrimeTriangle}\sep \lean{IsAuxTwist}\sep \lean{unionWithConj}\sep
\lean{pointChar}\sep \lean{twistAtPoint}\sep \lean{reduceModPow}\sep \lean{IsBalanced}\sep \lean{IsPencil}\sep
\lean{tieRelations}\sep all defined in the first note of the unit}

\sect{The definitions of this note}{s:defs}

\begin{definition}[evaluation at a conjugate pair]\label{d:evpair}
Let $E$ be a CM field as in the conventions, $n$ a finite set, $A$ a commutative ring, $B$ a commutative
$A$-algebra with the discrete topology, and \m{g\in\GE{E}}. Then \m{\evalConjPair{g}} is the $A$-linear map that sends a function
\m{z\colon\GE{E}\to B^n} to the pair \m{(z(g^{[+]}),z(g^{[-]}))\in B^n\times B^n}.
\end{definition}
\leanline{\lean{evalConjPair}}

\begin{definition}[flattening of blocks]\label{d:flat}
Let $N$ and $\sigma$ be finite sets and $\mathbb{L}$ a commutative algebra over a commutative ring~$\Lambda$. Then
\m{\blockFlattenAlgEquiv{N}} is the isomorphism of $\Lambda$-algebras from the \m{N\times N} matrices with entries in
\m{M_\sigma(\mathbb{L})} to the matrices over~$\mathbb{L}$ indexed by \m{N\times\sigma}.
\end{definition}
\leanline{\lean{blockFlattenAlgEquiv}}

\begin{definition}[the normaliser]\label{d:normaliser}
Let $\Gamma$ be a group, $\Lambda$ a commutative ring, \m{\gamma,a\in\Gamma}, \m{y,p\in\Lambda} and $\theta$ a finitely
supported function \m{\Gamma\to\NN}. Then \m{\ribetNormaliser{\gamma}{a}{y}{p}{\theta}{\cdot}} is the finitely supported
function \m{\Gamma\to\Lambda} which is the sum of the indicator of~\m{\gamma a}, of $-y$ times the indicator of~$a$, and of
\m{p\,\theta}; its value at \m{g\in\Gamma} is \m{\ribetNormaliser{\gamma}{a}{y}{p}{\theta}{g}}. Paired with a function
\m{g\mapsto u(g)-1} it gives \m{(u(\gamma a)-1)-y\,(u(a)-1)+p\sum_g\theta(g)(u(g)-1)}.
\end{definition}
\leanline{\lean{ribetNormaliser}}

\begin{definition}[triangle cost]\label{d:tricost}
Let $A$ be a commutative ring, $\Gamma$ a set, $J$ a totally ordered set, $B$ an $A$-module, $Z$ a submodule of the
functions \m{\Gamma\to B}, \m{\tau,\phi\colon J\times\{+,-\}\to\Gamma} functions, $I$ a finite subset of~$J$,
\m{C\in A} and \m{b\in\NN\cup\{\infty\}}. Then \m{\HasTriangleCost{Z}{\tau}{\phi}{I}{C}{b}} holds if there is a subset
\m{\Gamma_0\subseteq\Gamma} such that
\begin{enumerate}[label=\textup{(C\arabic*)},leftmargin=3em,itemsep=1pt]
\item for every \m{j\in I} and every \m{c\in Z} that vanishes on~$\Gamma_0$: if for all \m{i<j} one has
\m{c(\tau_{(i,+)})=c(\tau_{(i,-)})}, and also \m{c(\phi_{(i,+)})=c(\phi_{(i,-)})} when \m{i\in I}, then
\[C\,c(\tau_{(j,+)})=C\,c(\tau_{(j,-)});\]
\item for every submodule $Z'$ of the elements of~$Z$ that vanish on~$\Gamma_0$, if every \m{c\in Z'} vanishes at
all~$\tau_k$ and has \m{c(\phi_{(j,+)})=c(\phi_{(j,-)})} for \m{j\in I}, then
\[\len_A\bigl(Z/\{c\in Z:c|_{\Gamma_0}=0\}\bigr)+\len_A\bigl(Z'/Z'[C]\bigr)\le b,\]
with \m{Z'[C]} the module of the elements of~$Z'$ killed by~$C$.
\end{enumerate}
\end{definition}
\leanline{\lean{HasTriangleCost}}

\begin{definition}[trace ideals]\label{d:trid}
Let $A$ be a commutative ring, $\mathbb{L}$ a commutative $A$-algebra, $Y$, $J$, $K$ sets ($J$ is totally ordered
in the formal text, which the definition does not use), $\sigma$ a finite set,
$\Gamma$ a group, $\mathcal{S}$ an $A$-subalgebra of the functions \m{Y\to\mathbb{L}[X_j:j\in J]} and $\mathcal{J}$ an
$\mathcal{S}$-submodule of these functions. Let $A_1$ be a commutative $A$-algebra over which $\mathbb{L}$ is an
algebra, \m{u=(u_x)_{x\in Y}} a family of functions \m{u_x\colon\Gamma\to M_\sigma(A_1)}, \m{p\in A},
\m{\tau,\phi\colon K\to\Gamma}, \m{k\colon K\to J} and \m{\varepsilon=(\varepsilon_j)_{j\in J}} polynomials in~\m{A[X]}.
Then \m{\IsTraceIdeal{\mathcal{J}}{u}{p}{\tau}{\phi}{k}{\varepsilon}} holds if
\begin{enumerate}[label=\textup{(\arabic*)},leftmargin=2.3em,itemsep=1pt]
\item every constant function $X_j$ lies in~$\mathcal{S}$, and for all \m{g\in\Gamma} and all~$i$ the function
\m{x\mapsto} (the $i$-th coefficient of the characteristic polynomial of \m{u_x(g)}) lies in~$\mathcal{S}$;
\item for every \m{g\in\Gamma}, the function \m{x\mapsto p\,\tr u_x(g)} lies in~$\mathcal{J}$, and for every \m{i\in K}
so does
\m{x\mapsto\tr\bigl((u_x(\tau_i)-1)u_x(g)\bigr)-X_{k(i)}\,\tr\bigl(((X-1)\varepsilon_{k(i)})(u_x(\phi_i))\,u_x(g)\bigr)}.
\end{enumerate}
\end{definition}
\leanline{\lean{IsTraceIdeal}}

\begin{definition}[Ribet elements]\label{d:rib}
Let $A$, $\mathbb{L}$, $Y$, $\mathcal{S}$, $\mathcal{J}$ be as in Definition~\ref{d:trid} with $J$ finite (and totally ordered), let
\m{\beta\colon Y\to\mathbb{L}} be a function, \m{\mathcal{X}\subseteq Y}, \m{\pi\in\mathcal{X}} and \m{c\in\mathbb{L}}.
Then \m{\IsRibetElement{\mathcal{J}}{\beta}{\pi}{c}} holds if $\beta$, as a function with constant polynomials as
values, lies in~$\mathcal{S}$, and there is \m{f\in\mathcal{S}} with \m{f\beta\in\mathcal{J}} such that for every
\m{x\in\mathcal{X}}: \m{\beta(x)=0} and the coefficient in \m{f(x)} of every monomial that properly divides
\m{\prod_{j\in J}X_j} is zero; and the coefficient of \m{\prod_{j\in J}X_j} in \m{f(\pi)} is~$c$.
\end{definition}
\leanline{\lean{IsRibetElement}}

\sect{Roots of unity at which polynomials keep low valuations}{s:roots}

\begin{lemma}[one variable]\label{l:rootone}
Let $\kappa$ be a finite set, \m{y=(y_k)_{k\in\kappa}} real numbers and $r$ a real number. Suppose \m{y_k>0} for
all~$k$, \m{0<r}, \m{r<1} and \m{|\ell|<r^{\ell-1}}. Then for all large $N$ and, given~$N$, all large~$\cO$: for every
family \m{(F_k)_{k\in\kappa}} of polynomials in~\m{\cO[X]} such that for each $k$ some coefficient $a_{k,i}$ of~$X^i$
in~$F_k$ has \m{y_k\le|a_{k,i}|\,r^{i}}, there is \m{\zeta\in\cO} with \m{\zeta^{\ell^{N}}=1} and
\m{r\,y_k\le|F_k(\zeta-1)|} for all~$k$.
\end{lemma}
\begin{proof}
For \m{\rho>0} let \m{\|F\|_\rho=\max_i|a_i|\rho^{i}} be the Gauss norm, a multiplicative non-archimedean norm, and
put \m{\rho_j=r^{1/\ell^{j}}}, an increasing sequence below~$1$. Choose $a$ with \m{r^{a}<|\ell|}, let $N$ be large
(how large is fixed at the end), and let $\cO$ contain primitive roots of unity $\eta_j$ of order \m{\ell^{j+1}}
for \m{j<N}. Suppose no $\zeta$ works. Fix \m{j<N} and \m{\eta=\eta_j}; for every \m{i<\ell^{j+1}} some \m{k=k(i)} has
\m{|F_{k}(\eta^{i}-1)|<r\,y_{k}}. Let \m{\Pi=\prod_kF_k} and \m{\Pi=F_kQ_k}. Differentiating the identity
\m{(1+X)^{\ell^{j+1}}-1=\prod_i(X-(\eta^{i}-1))} and inserting, at each~$i$, the remainder of \m{F_{k(i)}} on division by
\m{X-(\eta^{i}-1)}, one writes \m{\ell^{j+1}(1+X)^{\ell^{j+1}-1}\,\Pi^{a(j+1)}} as a sum of terms each divisible by
\m{(1+X)^{\ell^{j+1}}-1}, plus terms carrying a factor \m{F_{k(i)}(\eta^{i}-1)}. Since
\m{|\eta^{i}-1|\le\rho_t} for \m{t\ge j} (the absolute value of \m{\eta_j-1} is \m{|\ell|^{1/(\ell^{j}(\ell-1))}<\rho_j}
by the hypothesis on~$r$), the linear factors have Gauss norm~$\rho_t$ at the radius~$\rho_t$, and the ultrametric
inequality yields an inequality between the Gauss norms of~$\Pi$ at the radii \m{\rho_j} and \m{\rho_{j+1}}:
\[
\log\|\Pi\|_{\rho_{j+1}}-\log\|\Pi\|_{\rho_j}\ \ge\ \frac{-\log r}{a\,(j+1)}\cdot(\text{a positive constant}).
\]
Summing over \m{j<N}, the left side telescopes and is bounded (\m{\|\Pi\|_\rho\le1} and
\m{\|\Pi\|_{\rho_0}\ge\prod_ky_k}), while the right side grows like the harmonic series. This is impossible for
large~$N$. The estimate of the Gauss norms is a computation of twenty lines in the formal proof, of which this
is the outline.
\end{proof}
\leanline{\lean{existsrootofunityevallarge}}

\begin{lemma}[several variables and all powers up to $m$]\label{l:rootmany}
Let $\kappa$ be a finite set, \m{\delta>0} a real number, \m{m,L,a,q} natural numbers and
\m{B=(B_k)_{k\in\kappa}} non-negative real numbers. Then for all large $N$ and, given~$N$, all large~$\cO$ the
following holds. Let $\sigma$ and $T$ be finite sets with \m{\#\sigma\le q}, \m{(D_n)_{n\in T}} monic polynomials
in~\m{\cO[X]} with \m{\sum_n\deg D_n\le L}, \m{(F_{k,d})_{k\in\kappa,1\le d\le m}} polynomials over~$\cO$ in the variables
indexed by~$\sigma$, and \m{(u_{k,d})} points of~$\cO^{\sigma}$ such that $\ell$ divides \m{(u_{k,d,j}-1)^{a}} for every
coordinate~$j$ and \m{\vl(F_{k,d}(u_{k,d}))\le B_k}. Then there is \m{\zeta\in\cO^{\sigma}} with
\[
\vl\bigl(F_{k,d}(\zeta^{d})\bigr)\le\delta+B_k\ \ \text{for all }(k,d),\qquad
\zeta_j^{\ell^{N}}=1\ \ \text{and}\ \ \vl(D_n(\zeta_j))\le\delta\ \ \text{for all }j,\ n,
\]
where \m{\zeta^{d}} is the tuple of the $d$-th powers.
\end{lemma}
\begin{proof}
Choose $r$ with \m{\max(|\ell|,\ell^{-\delta})<r^{\,q+L+\ell+ma+1}<1}. The tuple is built one coordinate at a time:
suppose the coordinates in a set \m{s\subseteq\sigma} have been replaced by roots of unity so that
\m{|F_{k,d}|} at the mixed point is at least \m{\ell^{-B_k}r^{\#s}}. For a further coordinate~$j$, restrict each
\m{F_{k,d}} to the line where only the $j$-th coordinate varies, substituting \m{(1+X)^{d}} for its $d$-th power;
the hypothesis \m{\ell\mid(u_{k,d,j}-1)^{a}} shows that the value at the old point is, up to the factor
\m{r^{da}}, a lower bound for the Gauss norm of radius~$r$ of the restriction. Lemma~\ref{l:rootone}, applied to
these restrictions and to the Taylor shift of \m{\prod_nD_n} (a monic polynomial, of Gauss norm at least
\m{r^{L}}), gives a root of unity for the $j$-th coordinate that loses at most one more factor~$r$. After
\m{\#\sigma\le q} steps all coordinates are roots of unity, and the choice of~$r$ turns the accumulated factors into
the bound~$\delta$.
\end{proof}
\leanline{\lean{existsrootsofunityellValle}}

\begin{lemma}[points on a line through a root of unity]\label{l:linepoint}
Let $A$ be a commutative ring, \m{t,u\in A}, and \m{k,M,N} natural numbers. Suppose \m{u^{\ell^{k}}=1} and
\m{k+M\le N}. Then $\ell$ divides \m{\bigl((1+t(u-1))-1\bigr)^{\ell^{k}}}, and \m{\ell^{M}} divides
\m{(1+t(u-1))^{\ell^{N}}-1}.
\end{lemma}
\begin{proof}
Modulo~$\ell$ one has \m{(x+y)^{\ell^{k}}\equiv x^{\ell^{k}}+y^{\ell^{k}}}. So
\m{(u-1)^{\ell^{k}}\equiv u^{\ell^{k}}-1=0}, and then \m{(1+t(u-1))^{\ell^{k}}\equiv1}. If \m{x\equiv1\bmod\ell} then \m{x^{\ell^{M}}\equiv1\bmod\ell^{M}}, and
\m{\ell^{N}} is a multiple of \m{\ell^{k+M}}.
\end{proof}
\leanline{\lean{dvdoneaddmulpowsubone}}

\sect{The mirror lemma and the sharp polynomial}{s:mirror}

\begin{lemma}[mirror lemma]\label{l:mirror}
Let $A$ be a commutative ring, $\Gamma$ a group, $n$ a finite set, $\mathfrak{a}$ an ideal of~$A$ with
\m{A/\mathfrak{a}} finite, \m{\varrho\colon\Gamma\to M_n(A)} a homomorphism, \m{\chi\colon\Gamma\to A} a multiplicative
homomorphism, \m{\mathsf{c}'\colon\Gamma\to\Gamma} a homomorphism, $H$ and $U$ subgroups of~$\Gamma$ with $H$ normal,
\m{\gamma,g_0,z,g\in\Gamma}, \m{s,\alpha,\alpha'\in A}, \m{y\colon\Gamma\to(A/\mathfrak{a})^n} a function, $i$ an
index, $W$ a set, and \m{x=(x_w)_{w\in W}} functions that assign to each \m{g'\in\Gamma} an additive map
\m{x_w(g')\colon(A/\mathfrak{a})^n\to\mathbb{Q}/\mathbb{Z}}. Suppose:
\begin{enumerate}[label=\textup{(\arabic*)},leftmargin=2.3em,itemsep=1pt]
\item the $A$-span of \m{\varrho(\Gamma)} is \m{M_n(A)}, and \m{\mathsf{c}'(\mathsf{c}'(g'))=g'} for all~$g'$;
\item if \m{x_w(g')=x_v(g')} for all \m{g'\in U}, then \m{w=v};
\item \m{\varrho(g')m=m} for \m{g'\in U} and \m{m\in(A/\mathfrak{a})^n}; \m{\mathsf{c}'(H)\subseteq H\subseteq U}; and
every commutator \m{[g',\gamma]} lies in~$H$;
\item \m{s-1} and \m{s+1} are units; \m{\varrho(\mathsf{c}'(\gamma))^{-1}=s}, \m{\varrho(\gamma)=s},
\m{\chi(\mathsf{c}'(\gamma))=\chi(\gamma)=1} and \m{y(\gamma)=0};
\item \m{z\in U} and \m{\mathsf{c}'(z)\in U}; $y$ is a cocycle for~$\varrho$; and the coordinate \m{y(g_0)_i} is a
unit;
\item \m{x_w(a\cdot g'bg'^{-1})(m)=x_w(a)(m)+x_w(b)\bigl(\chi(g')\varrho(g')^{-1}m\bigr)} for all $w$, \m{g'\in\Gamma},
\m{a,b\in U} and~$m$;
\item \emph{(isotropy on the coset \m{zH})} \m{x_w(zh)(t\,y(zh))=x_w(\mathsf{c}'(zh))(t\,y(\mathsf{c}'(zh)))} for all $w$,
\m{h\in H} and \m{t\in A};
\item \m{\varrho(g)=\alpha} and \m{\varrho(\mathsf{c}'(g))^{-1}=\alpha'} are scalars.
\end{enumerate}
Then \m{\#W} is at most the number of the \m{m\in(A/\mathfrak{a})^n} with \m{(\alpha-\chi(\mathsf{c}'(g))\alpha')\,m=0}.
\end{lemma}
\begin{proof}
Let \m{\tau=[g_0,\gamma]\in H}. By the cocycle identity and (4), \m{y(\tau)=(1-s)\,y(g_0)} up to the action of
elements that act trivially, so \m{y(\tau)_i} is a unit. Since $\varrho$ spans all matrices, the vectors
\m{Y\,y(\tau)}, \m{Y\in M_n(A)}, fill \m{(A/\mathfrak{a})^n}. From (6), (7) applied to the elements \m{zh},
\m{zh\cdot k\tau k^{-1}} and \m{zh\gamma}-conjugates, and from the fact that \m{s^2-1} is a unit, one derives for all
\m{t\in A}, \m{k\in\Gamma}, \m{b\in H}:
\[
x_w(b)\bigl(t\,\varrho(k)y(\tau)\bigr)=x_w(\mathsf{c}'(\tau))\bigl(\chi(\mathsf{c}'(k))\,t\,\varrho(\mathsf{c}'(k))^{-1}y(\mathsf{c}'(b))\bigr).
\]
Taking \m{b=\mathsf{c}'(\tau)} and comparing $k$ with \m{kg} shows, by~(8), that \m{x_w(\mathsf{c}'(\tau))} vanishes on
\m{(\alpha-\chi(\mathsf{c}'(g))\alpha')\cdot(A/\mathfrak{a})^n}; so it factors through the cokernel of that scalar. The
map \m{w\mapsto x_w(\mathsf{c}'(\tau))} is injective: by (6) with \m{b=[g',\gamma]} the values of $x_w$ on~$U$ are
determined by \m{x_w(\mathsf{c}'(\tau))}, and then (2) applies. The dual of the cokernel has as many elements as the
kernel of the scalar. The formal proof is one identity obtained by a linear combination of six instances of
(6) and~(7); the display is what it yields.
\end{proof}
\leanline{\lean{cardletorsionofisotropic}}

\begin{lemma}[one polynomial bounds all twisted Selmer lengths, sharply at one point]\label{l:fitting}
Let $E$ be a number field with a fixed embedding in~\m{\Qb}, $n$ and $\sigma$ finite sets, \m{M,q} natural
numbers, $A_0$ a commutative ring, $A$ a discrete valuation ring of characteristic~$0$,
\m{\varrho\colon\GE{E}\to M_n(A_0)} a homomorphism, \m{e_0\colon A_0\to A} a ring homomorphism, $S$ a finite set of
finite places of~$E$, \m{h=(h_j)_{j\in\sigma}} homomorphisms \m{\GE{E}\to\mathbb{Z}/q}, $L$ a set of relations with
values in \m{M_n(A_0)}, and \m{y_0\in A^{\sigma}}. Suppose that every $h_j$ is locally constant and that
\m{\reduceModPow{\ell}{M}{\varrho}} is locally constant. Then there is a polynomial $F$ over~$A$ in the variables
indexed by~$\sigma$ such that
\[
v_A(F(y_0))\le\len_A\Sel{\twistAtPoint{h}{y_0}{e_0}{\varrho}{\cdot}}{S}{e_0(L)}{A/\ell^{M}},
\]
and such that for every discrete valuation ring $A'$, ring homomorphism \m{e\colon A\to A'} and \m{y\in A'^{\sigma}},
\[
\len_{A'}\Sel{\twistAtPoint{h}{y}{e\circ e_0}{\varrho}{\cdot}}{S}{(e\circ e_0)(L)}{A'/\ell^{M}}\le v_{A'}\bigl(e(F)(y)\bigr).
\]
\end{lemma}
\begin{proof}
Work over the polynomial ring \m{\Lambda=A[X_j:j\in\sigma]} with the universal twist
\m{\twistAtPoint{h}{X}{e_0}{\varrho}{\cdot}}. By Lemma~7.4 of the first note (with \m{m=\ell^{M}} and the
neighbourhood on which $\varrho$ modulo~$\ell^M$ and all $h_j$ are trivial) there are finitely many coordinates
and a module $\mathcal{R}$ of relations over~$\Lambda$ such that, for every specialisation \m{X\mapsto y} into a
ring~$A'$, the Selmer module with values in \m{A'/\ell^{M}} is the module of the vectors killed by the specialised
relations. For a discrete valuation ring, the length of such a module is at most the valuation of the
determinant of any square matrix with rows among the relations, and some such matrix attains the length (an
input of the unit on Fitting ideals). Choose a square matrix over~$\Lambda$ with rows in~$\mathcal{R}$ whose
specialisation at~$y_0$ attains the length there, and let $F$ be its determinant.
\end{proof}
\leanline{\lean{existssharpfittingpoly}}

\sect{Pencils of tied primes: proofs of Theorems B and A}{s:pencils}

\begin{lemma}[Selmer modules cut by relations lie in Selmer groups with local conditions]\label{l:selz}
Let $E$ be a CM field as in the conventions, $n$ a finite set, $A$ a commutative ring, $B$ a commutative
$A$-algebra with the discrete topology, \m{\varrho\colon\GE{E}\to M_n(A)} a homomorphism, $S$, $P'$, $Q'$ sets of
finite places, \m{K\subseteq L} sets of relations, and \m{x\in\Sel{\varrho}{S}{L}{B}}. Suppose
\m{S\subseteq P'\cup Q'}. Then $x$ is a continuous cocycle of the discrete module $B^n$ that is unramified
outside \m{P'\cup Q'} and whose restriction to the decomposition group at any chosen place over a place of~$P'$ is
a local cocycle of $\varrho$ and~$K$ there: an element of the Selmer group, in the sense of the unit on the
formula of Greenberg and Wiles, with these local conditions at~$P'$ and no condition at~$Q'$.
\end{lemma}
\begin{proof}
Unfold both definitions: the cocycle identity, continuity (local constancy), the vanishing on inertia outside
\m{S\subseteq P'\cup Q'}, and the local condition, witnessed by $x$ itself since \m{K\subseteq L}. A totally complex
field has no real places, so there is no condition at infinity.
\end{proof}
\leanline{\lean{memselmerZofmemselmerModule}}

\begin{lemma}[the point character of a pencil]\label{l:pointchar}
Let \m{\IsPencil{\ell}{S_0}{U}{W}{N}{\fq_0}{J}{x}} hold, with \m{x(v)=(l_v,s_v)} and
\m{l_v=(\fP_v,\gamma_v,\phi_v,h_v)}. \textup{(a)} For \m{v,v'\in J}: \m{h_{v'}(\gamma_v)} is $1$ if \m{v'=v} and $0$
otherwise. \textup{(b)} Let $\Lambda$ be a commutative ring and \m{y=(y_v)_{v\in J}} elements with
\m{y_v^{\ell^{N}}=1}. Then the point character \m{y^{h}} takes the value $y_v$ at \m{\gamma_v^{[b]}} for both signs~$b$,
and, for \m{v\ne\fq_0}, the value \m{y_{\fq_0}} at \m{\phi_v^{[b]}}.
\end{lemma}
\begin{proof}
(a) $h_{v'}$ is from $\mathrm{U}(1)$ with modulus~\m{\{v'\}}, hence trivial on the inertia above the other places
of~$J$, and \m{h_v(\gamma_v)=1} by (P2). (b) follows from (a), (P3) and the fact that a character from
$\mathrm{U}(1)$ takes the same value at $g$ and at~\m{g^{[-]}}.
\end{proof}
\leanline{\lean{chartamegen}\sep \lean{pointCharidOrConjInv}}

\begin{lemma}[dropping ties costs only determinants]\label{l:drop}
Let $E$ be a CM field as in the conventions, $A$ a commutative ring and $\mathcal{I}$ a set. Suppose that
\m{\IsPencil{\ell}{S_0}{U}{W}{N}{\fq_0}{J}{x}} holds; write \m{x(v)=(l_v,s_v)} for the value of~$x$ at a place~$v$,
where \m{l_v=(\fP_v,\gamma_v,\phi_v,h_v)}. Let $n$ be a finite set, \m{\varrho\colon\GE{E}\to M_n(A)} a
homomorphism, \m{i\in\mathcal{I}}, and \m{k,M\in\NN}. Suppose that
\m{\reduceModPow{\ell}{M}{\varrho}(g)=1} for all \m{g\in\inab{S_0^{\complement}}}, that \m{k+M+M\le N}, and that
\m{U_i(g)} implies \m{\reduceModPow{\ell}{M}{\varrho}(g)=1}. Let $A'$ be a discrete valuation ring, \m{e\colon A\to A'}
a ring homomorphism, \m{u\in A'} with \m{u^{\ell^{k}}=1}, $I$ a set of finite places with \m{I\subseteq J}, such that
\m{u=1} if \m{\fq_0\in I}, and let \m{(w_{v,b})_{v\in I,b}} be elements of~\m{\GE{E}} with
\m{\reduceModPow{\ell}{M}{\varrho}(\phi_v^{[b]})=\reduceModPow{\ell}{M}{\varrho}(w_{v,b})} for all \m{v\in I} and both
signs~$b$. Let $S$ be a set of finite places and $L$ a set of relations. Put \m{y_{\fq_0}=u} and
\m{y_v=1+e(s_{v,i})(u-1)} for \m{v\in J\setminus\{\fq_0\}}, and
\m{\varrho_y=\twistAtPoint{h}{y}{e}{\varrho}{\cdot}}. Then
\begin{align*}
\len_{A'}\Sel{\varrho_y}{S}{e(L)}{A'/\ell^{M}}\ \le{}&\ \len_{A'}\Sel{\varrho_y}{S}{e(L\cup\tieRelations{I}{x}{i})}{A'/\ell^{M}}\\
&+\sum_{v\in I,\,b}v_{A'}\bigl(\det(u\cdot e(\varrho(w_{v,b}))-1)\bigr).
\end{align*}
\end{lemma}
\begin{proof}
The module with the ties is cut out of the module without them by the kernels of the maps
\m{c\mapsto c(\gamma_v^{[b]})-s_{v,i}\,c(\phi_v^{[b]})}, \m{v\in I}; so it suffices to bound, for each \m{(v,b)}, the
length of the image of this map by \m{v_{A'}(\det(u\,e(\varrho(w_{v,b}))-1))}, which holds as soon as the image is
killed by the matrix \m{u\,e(\varrho(w_{v,b}))-1} (the length of the kernel of a matrix on \m{(A'/\ell^{M})^n} is at
most the valuation of its determinant). By Lemma~\ref{l:linepoint} the $y_v$ are \m{\ell^{N}}-th roots of unity
modulo~\m{\ell^{M}}, so \m{\varrho_y} modulo \m{\ell^{M}} is the twist of~$\varrho$ by the point character. Let $c$ be a
Selmer cocycle and write \m{g\cdot m=\varrho_y(g)m+c(g)} for the affine action on \m{(A'/\ell^{M})^n} that it
defines. At the place \m{\fP_v} (and at its conjugate for \m{b=-}) the action is tame, because \m{\ell^{k+2M}} divides
\m{\ell^{N}} and \m{\Nm(v)\equiv1} modulo~\m{\ell^{N}}: the Frobenius element and the tame generator commute in the
image (an input of the unit on tame cocycles). Evaluating the commutation at~$0$ gives
\m{(\varrho_y(\phi_v^{[b]})-1)\,c(\gamma_v^{[b]})=(\varrho_y(\gamma_v^{[b]})-1)\,c(\phi_v^{[b]})}. Now
\m{\varrho(\gamma_v^{[b]})\equiv1} (inertia outside~$S_0$), \m{\varrho_y(\gamma_v^{[b]})=y_v} and, for \m{v\ne\fq_0},
\m{\varrho_y(\phi_v^{[b]})=u\,\varrho(\phi_v^{[b]})} by Lemma~\ref{l:pointchar}; with \m{y_v-1=s_{v,i}(u-1)} this gives
\m{(u\varrho(\phi_v^{[b]})-1)\bigl(c(\gamma_v^{[b]})-s_{v,i}c(\phi_v^{[b]})\bigr)=0}, using (P5) when \m{s_{v,i}\ne0}
(then \m{\varrho(\phi_v^{[b]})\equiv1}) and the triviality when \m{s_{v,i}=0}. For \m{v=\fq_0} one has \m{u=1} and
the same identity.
\end{proof}
\leanline{\lean{lengthleofdropties}}

\begin{lemma}[mirrors: a Selmer module that no new prime shortens is short]\label{l:unseen}
Let $E$ be a CM field as in the conventions, $A$ the ring of integers of a finite extension of~\m{\Ql},
$\mathcal{I}$ a set, and suppose \m{\IsPencil{\ell}{S_0}{U}{W}{N}{\fq_0}{J}{x}}. Let $n$ be a finite set,
\m{\varrho\colon\GE{E}\to M_n(A)} a homomorphism, \m{i\in\mathcal{I}}, $S$ a finite set of finite places, $L$ a set of
relations, \m{\gamma\in\GE{E}}, \m{\beta\in[0,\infty]} and \m{M\in\NN}. Suppose
\m{\IsBalanced{\ell}{S}{\varrho}{L}{\gamma}{\beta}{M}}, that $U_i$ is the property
\m{\reduceModPow{\ell}{M}{\varrho}(g)=1}, that \m{S\subseteq S_0}, and let \m{a\in A} be non-zero with
\m{a\mid\ell^{M}}. For a finite set $T$ of finite places put
\[
\mathcal{X}(T)=\Sel{\varrho}{\unionWithConj{S}{T}}{L\cup\tieRelations{J}{x}{i}\cup\{\delta_\gamma\}}{A/a}.
\]
Suppose that some \m{y\in\mathcal{X}(J)} has a coordinate \m{y(g)_e} that is a unit, and that for every finite place
\m{\fq\notin J} and every \m{l\in\AuxPrimeData{E}{\ell^{N}}}, with elements \m{\gamma_l,\phi_l}, such that for all
\m{t\in A} the function $x$ changed at~$\fq$ to \m{(l,t\cdot\mathbf{1}_i)} satisfies
\m{\IsPencil{\ell}{S_0}{U}{W}{N}{\fq_0}{J\cup\{\fq\}}{\cdot}}, one has
\[
\mathcal{X}(J)\subseteq\bigl(\evalConjPair{\phi_l}\bigr)^{-1}\Bigl(\evalConjPair{\gamma_l}\bigl(\mathcal{X}(J\cup\{\fq\})\bigr)\Bigr).
\]
Then \m{\len_A\mathcal{X}(J)\le\beta}.
\end{lemma}
\begin{proof}
Let $X$ be the discrete module \m{(A/a)^n} with the action through~$\varrho$, and let $U'$ be the intersection
of the kernels of \m{\reduceModPow{\ell}{M}{\varrho}} and of \m{\cyc_{\ell^{M+1}}}, an open subgroup acting trivially
on~$X$. By
Lemma~\ref{l:selz}, \m{\mathcal{X}(J)} lies in the Selmer group of~$X$ with the local conditions of $\varrho$ and
\m{L\cup\tieRelations{J}{x}{i}} at the places of \m{\unionWithConj{S}{J}}, and vanishing at~$\gamma$. The formula of
Greenberg and Wiles in the form with relations at split primes (Input~1 below, applied with \m{K=E},
\m{m=\ell^{M+1}}, the places \m{\unionWithConj{S}{J}} and \m{N'=\mathbb{Q}/\mathbb{Z}}; its hypothesis on trivialised
cocycles holds because $\gamma$ is central modulo~$U'$ and \m{s-1} is a unit) bounds the order of this Selmer
group, by Input~1(a), by the product of the orders of the local conditions, divided by a power of~\m{\#X}, times
the number of the \emph{relation families} \m{(n_i)}. By Input~1(b) each relation family defines, through
\m{x_w(g')=n_{\varpi'(g')}} for \m{g'\in U'}, a function on~$U'$ with values in the dual of~$X$ that satisfies
hypothesis~(6) of the mirror lemma and the reciprocity relations at the split places outside~$T'$; these functions
$x_w$, $w$ running over the relation families, are the family to which the mirror lemma is applied, and they are
separated by their values on~$U'$ because \m{n_{\varpi'(\phi_i)}=n_i}. As $\gamma$ acts by a
scalar $s$ with \m{s-1} a unit and is central modulo~$U'$, evaluation at~$\gamma$ maps the Selmer group onto~$X$
(coboundaries), which accounts for the condition at~$\gamma$.

The number of relation families is bounded by the mirror lemma, Lemma~\ref{l:mirror}, applied with
\m{\chi=\cyc}, \m{\mathsf{c}'=(\cdot)^{\mathsf c}}, $H$ the intersection with its conjugate of the group of commutators
with the normal closure of \m{\{\gamma,\gamma^{[-]}\}} inside~$U'$, the cocycle~$y$ with its unit coordinate,
\m{z=\gamma_{\fq_0}} times a commutator, and the scalars \m{\alpha,\alpha'} of (B5). Its isotropy hypothesis is where
the hypothesis on new primes enters. Given \m{h\in H}, Theorem~A of the first note provides a prime \m{\fq} of
prime norm outside all places in play, with a datum $l$ whose Frobenius element is so close to \m{zh} that the
pairings in question take the same values, and which is anticyclotomic relative to \m{\pd{E}{\ell}\cup J\setminus\{\fq_0\}};
changing $x$ at~$\fq$ gives pencils for all slopes. By hypothesis the values of $y$ at
\m{\phi_l^{[\pm]}} are the values at \m{\gamma_l^{[\pm]}} of some \m{x'\in\mathcal{X}(J\cup\{\fq\})}. The values of $x'$ on the
inertia at~$\fq$ are multiples of its value at a tame generator, and the reciprocity relation of the relation
family at the pair of split places \m{\fq,\cj{\fq}} then reads exactly as the isotropy identity~(7).

It remains to convert orders into lengths over~$A$ and to sum the local terms. At a place of \m{J\cup\cj{J}} the
local cocycles that obey the ties are determined by their value at the Frobenius element (the tame generator
detects ramification on the appropriate side), so the local term is at most \m{\len_A((A/a)^n)}; at the places
of~$S$ one takes the places of~\m{\Qb} that minimise the local length; the torsion of the mirror lemma has
length at most \m{\#n\cdot v_A(\alpha-\cyc(g^{\mathsf c})\alpha')}. The inequality (B5) at~$a$ is precisely what makes
the total at most~$\beta$. The formal proof has a hundred lines; this is its outline, and the bookkeeping of the
local terms was not re-derived here.
\end{proof}
\leanline{\lean{lengthleofunseenclass}}

\inp{Input 1 (relations at split primes; background library, unit of the formula of Greenberg and Wiles)}{Let
$K$ be a number field with a fixed embedding in~\m{\Qb} (a CM field or not), \m{\GE{K}=\Gal(\Qb/K)}, \m{m\ge1} an
integer, $X$ a discrete module of~\m{\GE{K}} over a commutative ring, killed by~$m$, $S$ a finite set of finite
places of~$K$, with a chosen place of~\m{\Qb} over every finite place~$v$ and a submodule $L_v$ of the functions
from~\m{\GE{K}} to~$X$ (the \emph{local condition} at~$v$), $N'$ an abelian group, and $U'$ an open normal subgroup
of~\m{\GE{K}}. Suppose that the inertia groups above the places outside~$S$ act trivially on~$X$, that $L_v$
contains the coboundaries of the decomposition group at the chosen place for every \m{v\in S}, that every complex
conjugation in~\m{\GE{K}} acts trivially on~$X$, that $U'$ acts trivially on~$X$ and on the $m$-th roots of unity,
and that $U'$ contains no complex conjugation. For a set $\Pi'$ of finite places let \m{Y(\Pi')} be the
\emph{Selmer group relaxed at~$\Pi'$}: the continuous cocycles \m{\GE{K}\to X} that vanish on the inertia groups
above the places outside \m{S\cup\Pi'} and on the complex conjugations, and whose restriction to the decomposition
group at the chosen place lies in~$L_v$ for every \m{v\in S}. Then there are an integer \m{q\ge1}, a finite set
$\iota'$, distinct finite places \m{(v_i)_{i\in\iota'}} of~$K$ not dividing~$q$ with arithmetic Frobenius elements
\m{\phi_i\in U'} at places of~\m{\Qb} over them (\emph{split places} in the sense of the library, for the
modulus~$m$, the level~$q$ and~$U'$), elements \m{(\gamma_i)_{i\in\iota'}} of~\m{\GE{K}}, a locally constant map
\m{\varpi'\colon\GE{K}\to\iota'} and a finite set \m{T'\supseteq S} of finite places, such that every place
of~$S$ divides~$q$, the \m{\phi_i} detect roots (every element of \m{\Qb^{\times}} whose $m$-th power is a unit
outside~$q$ fixed by~$U'$, and which every \m{\phi_i} fixes, is fixed by~$U'$), and the following hold.
\begin{enumerate}[label=\textup{(\alph*)},leftmargin=2.3em,itemsep=2pt]
\item \emph{The count.} If $X$ is finite and has at most \m{(\#X)^{\#\iota'+1}} trivialised cocycles (pairs of a
function $f$ from the places of~\m{\Qb} above $q$ and~$\infty$ to~$X$ and a continuous function
\m{c\colon\GE{K}\to X} that vanishes on the inertia groups above the places not dividing~$q$ and different from
the~$v_i$, with \m{g\,f(x)=f(gx)+c(g)}), then
\[
\#Y(\emptyset)\cdot(\#X)^{\#S+\#\{\text{infinite places of }K\}}\ \le\ \#\mathcal{R}_{\mathbb{Q}/\mathbb{Z}}\cdot\#X\cdot\prod_{v\in S}\#\bigl(L_v\cap Z^1_v\bigr),
\]
where \m{Z^1_v} is the module of the cocycles of the decomposition group at the chosen place over~$v$ and, for an
abelian group~$G'$, \m{\mathcal{R}_{G'}} is the set of the \emph{relation families}: the families
\m{(n_i)_{i\in\iota'}} of homomorphisms \m{X\to G'} with \m{\sum_in_i(x(\gamma_i))=0} for all
\m{x\in Y(\{v_i:i\in\iota'\})}.
\item \emph{The relations.} Every \m{(n_i)\in\mathcal{R}_{N'}} satisfies \m{n_{\varpi'(\phi_i)}=n_i} for all~$i$;
\m{n_{\varpi'(a\,gbg^{-1})}(x)=n_{\varpi'(a)}(x)+n_{\varpi'(b)}\bigl(\cyc_m(g)\,g^{-1}x\bigr)} for all
\m{g\in\GE{K}}, \m{a,b\in U'} and \m{x\in X}; and the reciprocity relation
\m{\sum_jn_{\varpi'(\phi'_j)}(x(\delta_j))=0} for every finite family of distinct finite places \m{v'_j\notin T'}
with places of~\m{\Qb} over them, tame generators $\delta_j$ for~$m$ and Frobenius elements \m{\phi'_j\in U'}
there, and every \m{x\in Y(\{v'_j\})}.
\end{enumerate}
The set $T'$ is what the further places must avoid. For a totally complex~$K$ the conditions on complex
conjugations are empty; the third note applies the statement over a totally real field, where they are not.
\emph{Reference:} to be confirmed; the catalogue of this write-up classes the statement as a substantial result
of the library, part of the lower bound in the formula of Greenberg and Wiles, without a printed source for
this form. The hypothesis of~(a) is supplied by a second result of the same unit (few trivialised cocycles when
an element central modulo~$U'$ acts on~$X$ without fixed vectors), for which the second reading located no
published derivation either.}

\begin{lemma}[tilting]\label{l:tilting}
Let $A$ be a discrete valuation ring, $m$ a finite set, $B$ an $A$-module, $\mathfrak{a}$ an ideal of~$A$ such that
\m{A/\mathfrak{a}} is artinian, \m{\Phi_0,\Phi_1\colon B\to(A/\mathfrak{a})^m} $A$-linear maps and $Y$ a submodule of~$B$.
Then there is a list $T$ of at most \m{\#m} elements of the residue field of~$A$ such that for every \m{s\in A}
whose residue is not in~$T$,
\[
\len_A\bigl(Y\cap\ker(s\Phi_0-\Phi_1)\bigr)\ \le\ \len_A\bigl(Y\cap\ker\Phi_0\cap\Phi_1^{-1}(\Phi_0(Y))\bigr).
\]
\end{lemma}
\begin{proof}
Let \m{Z=\sum_{s}Y\cap\ker(s\Phi_0-\Phi_1)}, the sum over all \m{s\in A}. On each summand \m{\Phi_1=s\Phi_0}; hence
\m{\Phi_1(Z)\subseteq\Phi_0(Z)}. Let \m{\mathcal{M}\subseteq A^m} be the preimage of \m{\Phi_0(Z)} under the reduction \m{A^m\to(A/\mathfrak{a})^m},
\m{v\mapsto\bar v}; it is a free $A$-module of rank at most~\m{\#m}. As $\mathcal{M}$ is projective and $\Phi_0$
maps $Z$ onto \m{\Phi_0(Z)}, there is an $A$-linear map \m{\sigma'\colon\mathcal{M}\to Z} with
\m{\Phi_0(\sigma'(v))=\bar v} for all \m{v\in\mathcal{M}}; and, as \m{\Phi_1(Z)\subseteq\Phi_0(Z)} and $\mathcal{M}$
maps onto \m{\Phi_0(Z)}, there is an endomorphism $\psi$ of~$\mathcal{M}$ with
\m{\overline{\psi(v)}=\Phi_1(\sigma'(v))} for all~$v$. Let $T$ be the list of the roots, in the residue field, of the
reduction of the characteristic polynomial of~$\psi$; it has at most \m{\#m} elements. Let the residue of~$s$
not be in~$T$. Then \m{\det(s-\psi)} is a unit, so \m{s-\psi} is an automorphism of~$\mathcal{M}$, and the identity
\m{(s\Phi_0-\Phi_1)(\sigma'(v))=\overline{(s-\psi)(v)}} shows that
\m{(s\Phi_0-\Phi_1)(Z)\supseteq\Phi_0(Z)}. (The map $\Phi_1$ need not vanish on \m{Z\cap\ker\Phi_0}, so the two
kernels on~$Z$ may differ; only their lengths are compared.) Both images lie in \m{(A/\mathfrak{a})^m}, a module of
finite length, and in \m{\NN\cup\{\infty\}}
\begin{align*}
\len_A\bigl(Z\cap\ker(s\Phi_0-\Phi_1)\bigr)+\len_A\bigl((s\Phi_0-\Phi_1)(Z)\bigr)&=\len_AZ\\
&=\len_A\bigl(Z\cap\ker\Phi_0\bigr)+\len_A\Phi_0(Z).
\end{align*}
The second terms are finite, and that on the left is at least that on the right. Hence
\[
\len_A\bigl(Z\cap\ker(s\Phi_0-\Phi_1)\bigr)\ \le\ \len_A\bigl(Z\cap\ker\Phi_0\bigr).
\]
Finally \m{Y\cap\ker(s\Phi_0-\Phi_1)} is contained in~$Z$, so the left side here is the left side of the lemma;
and \m{Z\cap\ker\Phi_0} is contained in \m{Y\cap\ker\Phi_0\cap\Phi_1^{-1}(\Phi_0(Y))}, because
\m{\Phi_1(Z)\subseteq\Phi_0(Z)\subseteq\Phi_0(Y)}.
\end{proof}
\leanline{\lean{tilting}}

\begin{proof}[Proof of Theorem~B]
By induction on the factorisation of~$a$ into irreducible elements. If $a$ is a unit, the Selmer modules are
zero; take \m{J_0=1}, and for given~$N$ one helper prime $\fq_0$ from Theorem~A of the first note (applied with
\m{R=S_0} and the neighbourhood $W$ of~$z$), with slopes~$0$: this is a pencil. Now suppose the statement holds
for~$a$ with a bound~$J_1$, let $\varpi$ be a uniformiser with \m{\varpi a\mid\ell^{M}}, and fix~$N$. All Selmer
modules below are finite (Lemma~7.4 of the first note). Among the pencils \m{(J,x)} with the given helper prime
whose lengths modulo~$a$ are at most~$\beta_i$ and for which \m{\#J} plus the sum over~$i$ of the lengths
modulo~\m{\varpi a} is at most a bound \m{J_0}, choose one for which the sum of the lengths modulo~\m{\varpi a} is
least. The bound \m{J_0} depends on $J_1$, the ranks \m{\#n_i}, \m{\len_A(A/\varpi a)} and the (finite) lengths of
the Selmer modules of the \m{\varrho_i} modulo~\m{\varpi a} that are unramified outside~$S$, but not on~$N$: it is
chosen so that the pencil of the induction hypothesis qualifies. For that pencil, the cocycles
modulo~\m{\varpi a} that vanish at the tame generators \m{\gamma_v^{[\pm]}}, \m{v\in J}, are unramified at~$J$ and
its conjugates (Lemma~7.3 of the first note), so the length modulo~\m{\varpi a} is at most
\m{2\,\#n_i\,\#J\cdot\len_A(A/\varpi a)}, the length of the target of the evaluation at these generators, plus the
length of that Selmer module; and \m{\#J\le J_1}. We claim that the chosen pencil has all lengths
modulo~\m{\varpi a} at most~$\beta_i$. Fix~$i$. If no element of the module modulo~\m{\varpi a} has a unit
coordinate, the module is the image of the module modulo~$a$ under multiplication by~$\varpi$, of length at
most~$\beta_i$. Otherwise, if the inclusion of Lemma~\ref{l:unseen} holds at every admissible new prime, that
lemma gives the bound. In the remaining case there are a new prime~$\fq$ and a datum~$l$ at which the inclusion
fails. Apply Lemma~\ref{l:tilting} to \m{\Phi_0=\evalConjPair{\gamma_l}}, \m{\Phi_1=\evalConjPair{\phi_l}} on the
module for \m{J\cup\{\fq\}} with the old ties, modulo~$a$ and modulo~\m{\varpi a}: the residue field has more than
\m{4\#n_i+2} elements (B2), so there is a unit slope $t$ whose inverse avoids both lists of bad residues. For the
pencil \m{(J\cup\{\fq\},x')}, with $x'$ equal to $x$ changed at~$\fq$ to \m{(l,t^{-1}\mathbf{1}_i)}, the tied module is
contained in \m{\ker(\Phi_0-t^{-1}\Phi_1)}, so its length is at most that of
\m{\ker\Phi_0\cap\Phi_1^{-1}(\Phi_0(\cdot))}, which is contained in the module for~$J$ (vanishing at the tame generators
of~$\fq$ makes a cocycle unramified there) and, modulo~\m{\varpi a}, properly so because the inclusion failed. The
modules of the other rows only shrink, the lengths modulo~$a$ stay at most~$\beta_i$, and the count of primes
grows by one while the sum of lengths drops by at least one. This contradicts the minimality.
\end{proof}

\begin{proof}[Proof of Theorem~A]
This is an outline of the formal proof: its steps and the hypotheses used at each are those of the formal text;
the estimates inside the steps are left to it.

\emph{Step 1: a pencil.} Let $W$ be a neighbourhood of~$z$ on which every \m{\mu_{k,t}} stays within~$c$ of its value
(so that \m{\vl(\mu_{k,t}(\phi)-1)\le c}) and on which the \m{\varrho_i} modulo~\m{\ell^{M}} agree with their values at $z$
and at~\m{z^{[-]}}. Let $J_0$ be the bound of Theorem~B for \m{a=\ell^{M}}, choose \m{\delta>0} with
\m{\delta\,(8J_0+2)<\eta}, and choose $k$ and \m{N\ge k+2M} large as required by Lemma~\ref{l:rootmany} in the two
uses below. Theorem~B gives a pencil \m{(\fq_0,J,x)} for~$N$ with \m{\#J\le J_0} and tied lengths at most~$\beta_i$.

\emph{Step 2: polynomials.} For a tuple \m{y=(y_v)_{v\in J}} of \m{\ell^{N}}-th roots of unity the point character
\m{y^{h}} is from $\mathrm{U}(1)$ relative to \m{\pd{E}{\ell}} with modulus the set of the $v$ with \m{y_v\ne1}; every
subset of~$J$ is auxiliary for~$S_0$. By (F2), over a large coefficient ring \m{\cO'\supseteq\cO} each value
\m{\mathrm{val}_k(I,y^{h})}, \m{I\subseteq J}, is a polynomial \m{V_{k,I}(y)} with integral coefficients. By
Lemma~\ref{l:fitting} there are polynomials \m{F_i} whose valuations bound the Selmer lengths of the twists
\m{y^{h}\varrho_i} with the relations \m{L_i\cup\{\delta_{\gamma_i}\}}, and polynomials \m{G_i} that do the same with
the ties at \m{J\setminus\{\fq_0\}} added and are sharp at the point \m{y=1}.

\emph{Step 3: along a line.} Consider the line \m{y_{\fq_0}=X}, \m{y_v=1+s_v(X-1)}, with $s_v$ the slope of~$v$ for the
row~$i$ at hand; the pencil carries slopes for the rows only, and for a family one takes \m{s_v=0}, that is
\m{y_v=1} for \m{v\ne\fq_0}. At \m{X=1} the restricted polynomials have valuations at most $B_i$ (for $G_i$: sharpness,
the bound~$\beta_i$ of the pencil, Lemma~\ref{l:drop} for \m{I=\{\fq_0\}} with \m{w=z^{[\pm]}}, and hypothesis~(ii)) and
at most \m{\vl(\mathrm{val}_k(\emptyset,1))+8c} (for \m{V_{k,\{\fq_0\}}}: (F3) at the helper prime, whose Frobenius element
lies in~$W$). Lemma~\ref{l:rootmany} in one variable gives \m{u\in\cO'} with \m{u^{\ell^{k}}=1} at which these
valuations grow by at most~$\delta$ and at which the characteristic polynomials of the
\m{\varrho_i(\phi_v^{[b]})^{-1}} and the polynomials \m{X-\mu_{k,t}(\phi_v)^{\pm1}} have valuation at most~$\delta$.

\emph{Step 4: the point on the line.} Let $y$ be the point of the line at \m{X=u}; by Lemma~\ref{l:linepoint} its
coordinates are \m{\ell^{N}}-th roots of unity modulo the relevant powers, and $\ell$ divides
\m{(y_v-1)^{\ell^{k}}}. By Lemma~\ref{l:drop} with \m{I=J\setminus\{\fq_0\}} and \m{w=\phi_v^{[b]}}, the Selmer length of
\m{y^{h}\varrho_i} is at most the tied length plus at most \m{2J_0} determinant terms, each of valuation at
most~$\delta$ by Step~3; so \m{\vl(F_i(y))\le B_i+\delta+2J_0\delta} in normalised terms. By (F3), adding the places
of~\m{J\setminus\{\fq_0\}} one at a time raises \m{\vl(V_{k,\cdot}(y))} by at most~$8\delta$ each: at the Frobenius element
of~$v$ the point character takes the value \m{y_{\fq_0}=u} (Lemma~\ref{l:pointchar}), and \m{|1-\mu u^{\pm1}|} was
controlled in Step~3.

\emph{Step 5: roots of unity, for all powers.} Apply Lemma~\ref{l:rootmany} with the variables \m{v\in J}, the
powers \m{1\le d\le m}, the polynomials \m{F_i} and \m{V_{k,J}} (the same for every~$d$), the point~$y$ of Step~4
for every \m{(k,d)}, and the monic polynomials \m{D_d=X^{d}-1}, \m{1\le d\le m}, of total degree \m{m(m+1)/2}.
It gives \m{\zeta\in\cO'^{J}} with \m{\zeta_v^{\ell^{N}}=1} for all~$v$,
\[
\vl(\zeta_v^{d}-1)\le\delta\qquad\text{for all }v\in J\text{ and all }1\le d\le m,
\]
and the valuations of the \m{F_i} and \m{V_{k,J}} at \m{\zeta^{d}} within~$\delta$ of those at~$y$, for every~$d$.
The polynomial \m{X-1} alone would not do: for a root of unity \m{\zeta_v} of $\ell$-power order and $\ell$
dividing~$d$, \m{\vl(\zeta_v^{d}-1)} exceeds \m{\vl(\zeta_v-1)}, and \m{\zeta_v^{d}=1} is not excluded. Put \m{I=J} and
\m{\psi=\zeta^{h}}; then \m{\psi^{d}=(\zeta^{d})^{h}}, a point character at \m{\ell^{N}}-th roots of unity.
Conclusion (1) holds because the total loss is at most \m{\delta(8J_0+2)<\eta}; (2) because $J$ is auxiliary,
\m{\psi^{d}} is from $\mathrm{U}(1)$ with modulus~$J$ (Step~2), and
\m{\sum_v\vl(\psi^{d}(\gamma_v)-1)=\sum_v\vl(\zeta_v^{d}-1)\le\#J\cdot\delta\le\eta}
(Lemma~\ref{l:pointchar}); (3) by the second property of~\m{F_i} in Lemma~\ref{l:fitting}, for every coefficient
ring containing~$\cO'$. The same $I$ and $\psi$ serve all $d$ and all~$i$. The bound in (2) for every
power~$d$, not only for \m{d=1}, is what makes the twist $m$-fold in the third note.
\end{proof}

\sect{Balanced rows: proof of Theorem C}{s:balanced}

\begin{lemma}[classes of polynomials vanishing on roots of unity are tame]\label{l:tametype}
Let $\cO$ be a coefficient ring, $N$ and $r$ natural numbers, $K$ a property of polynomials over~$\cO$, and
$\mathfrak{M}$ a finite multiset of elements of~\m{\Zlb}. Suppose that $\mathfrak{M}$ has at most $N$ elements, that
every \m{y\in\mathfrak{M}} satisfies \m{y^{\ell^{a}-1}=1} for some \m{1\le a\le r}, and that $K$ holds for every
\m{Q\in\cO[X]} with \m{Q(y)=0} for all \m{y\in\mathfrak{M}}. Then $K$ is of tame type \m{(\ell,r,Nr)}.
\end{lemma}
\begin{proof}
A root of unity of order dividing \m{\ell^{a}-1} lies in the unramified extension of degree~$a$ of~\m{\Ql}; its
minimal polynomial over~\m{\Zl} is monic of degree at most \m{a\le r}. So the product of the minimal polynomials of
the elements of~$\mathfrak{M}$ is monic over~\m{\Zl}, of degree at most~\m{Nr}, and vanishes on~$\mathfrak{M}$; and
\m{X^{M'}-1} vanishes on~$\mathfrak{M}$ for the least common multiple $M'$ of the \m{\ell^{a}-1}, \m{a\le r}.
\end{proof}
\leanline{\lean{isTameTypeofroots}}

\begin{lemma}[the local lattices of a column]\label{l:inclass}
Let $D$ be admissible data, \m{\iota\colon\CC\to\Qlb} the fixed isomorphism, $\cO$ a coefficient ring, $n$ a finite
set, $\HT$ a type, \m{u\colon\AbsGal{D}\to M_n(\cO)} a homomorphism, \m{R\colon\AbsGal{D}\to\GL_n(\Qlb)} a
homomorphism, $t$ a natural number, $S$ a set of finite places, \m{w\in\pd{\ED{D}}{\ell}}, and $\mathcal{R}$ a
property of homomorphisms \m{\AbsGal{D}\to\GL_n(\Qlb)}. Suppose \m{\IsColumn{D}{t}{\HT}{S}{R}} and that $R$ has the
property~$\mathcal{R}$. Let \m{\fP_w} be the chosen place over~$w$, with its chart \m{e_w}, let
\m{\Lambda_w=\matrixLattice{u|_{G_{\fP_w}}}}, and \m{\HT_w(g)=\HT(\complexEmbOf{e_w|_{\ED{D}}}{g})} for \m{g\in\Gl}. If
$u$ is a lattice of~$R$, then for every \m{g\in\Gl} the lattice \m{\Lambda_w} has the jumps \m{\HT_w(g)} at
\m{\coeffEmb{\cO}{g}}, and \m{\Lambda_w} belongs to the class \m{\mathcal{Q}(F,K,U)}, where: $F$ is the trivial filter
(no condition) if \m{w\in\placesAboveG{D}} and \m{F=\top} otherwise; \m{K(s,Q)} holds if \m{w\in\placesAboveG{D}} and for
every $R'$ with the property~$\mathcal{R}$ there is a polynomial that kills $s$ on~$R'$ at~\m{\fP_w} and at all of
whose roots $Q$ vanishes; and \m{U(g')} is the set of the elements of \m{\HT_w(g)} for a $g$ with
\m{\coeffEmb{\cO}{g}=g'}.
\end{lemma}
\begin{proof}
Let \m{u(g)=Y^{-1}R(g)Y}. The map \m{x\mapsto Yx} from $\cO^n$ to~\m{\Qlb^n} intertwines \m{\Lambda_w} with the
restriction of~$R$ to \m{G_{\fP_w}}. By the definition of a column, at the place \m{\fP_w} above~$\ell$ the
representation $R$ is potentially crystalline with the jumps~$\HT$, and crystalline unless
\m{w\in\placesAboveG{D}}. These properties and the killing polynomials pass to a lattice along a jointly injective
family of equivariant maps, and the jumps are compatible with base change from~$\cO$ to~\m{\Qlb} (inputs of the
unit on crystalline representations).
\end{proof}
\leanline{\lean{inClassofisLatticeIn}}

\begin{lemma}[obeying cocycles and the cup-orthogonal]\label{l:obeycup}
Let $X$, $n$, $d$ be finite sets, $A$ the ring of integers of a finite extension of~\m{\Ql}, $\Gamma$ a
topological group over~\m{\Gl} with homomorphism~$j$, $\Gamma'$ a group, \m{\eta\colon\Gamma\to\Gamma'} a homomorphism,
\m{u=(u_x)_{x\in X}} homomorphisms \m{u_x\colon\Gamma'\to M_n(A)}, \m{\varrho\colon\Gamma'\to M_d(A)} a homomorphism,
\m{(F,K,U)} as in the definition of a class, and \m{m\in\NN}. Suppose that $j$ is an open embedding and that
\m{\matrixLattice{u_x\circ\eta}} belongs to \m{\mathcal{Q}(F,K,U)} for every~$x$. Let
\m{L=\matrixLattice{\varrho\circ\eta}}. If $L$ belongs to \m{\mathcal{Q}(F,K,U)}, then
\begin{align*}
&\len_A\Bigl(\obeyingCocycles{u}{\varrho}{\ell^{m}}\big/\{c:c\circ\eta\in\Ob{\Rel{\Zf(L)}}{A/\ell^{m}}\}\Bigr)\\
&\qquad\le\ \len_A\Bigl(\Zf(L^{\vee}(1))\big/\Zf(L^{\vee}(1))\cap(L^{\vee}(1))^{\perp\mathcal{Q}(F,K,U)}\Bigr),
\end{align*}
the first quotient being by the intersection with the set shown.
\end{lemma}
\begin{proof}
This is an outline of the formal proof.
There is a perfect pairing between \m{L/\ell^{m}} and the reduction of the twisted dual, so that the cocycles
with values in \m{(A/\ell^{m})^d} which obey the relations of \m{\Zf(L)} are those orthogonal to \m{\Zf(L^{\vee}(1))}
under the cup product; and for potentially crystalline lattices an index of orthogonal complements on one side
divides the corresponding index on the other (inputs of the unit on local duality; this is where $j$ must be an
open embedding). It therefore suffices to show that the cup product of an obeying cocycle~$c$, restricted
along~$\eta$, with a class of \m{\Zf(L^{\vee}(1))} that lies in the cup-orthogonal of the class is a coboundary. By
Lemma~4.1 of the first note, $c$ is the value of an equivariant map on a subrepresentation $Y_0$ of
\m{M_n(A)^{X}\times(A/\ell^{m})^d} that projects injectively to the first factor; \m{Y_0}, restricted along~$\eta$,
is a Galois lattice which embeds jointly into the lattices \m{\matrixLattice{u_x\circ\eta}}, hence belongs to the
class, and it carries a crystalline cocycle whose image is~$c$ (the coboundary of a point of the graph, divided
by a power of~$\ell$). Lifting the map \m{Y_0\to L/\ell^{m}} to~$L$ gives a pairing of $Y_0$ with \m{L^{\vee}(1)}
modulo~\m{\ell^{m}} in the sense of the definition of the cup-orthogonal, and the definition gives the
coboundary.
\end{proof}
\leanline{\lean{relLengthobeyingCocyclesle}}

\begin{proof}[Proof of Theorem~C]
This is an outline of the formal proof; the local length counts in (B5) rest on the inputs named.
The finiteness of~$\Sigma$ is clear. Let $\cO$ be large enough for (W3) and such that $R$ has a lattice
\m{u\colon\AbsGal{D}\to M_n(\cO)} (a continuous representation has one over every large coefficient ring; an input
of the background library). Let $F'$ be the algebra map of Lemma~5.6 of the first note. By (W3), \m{u(\AbsGal{D})}
spans \m{M_n(\cO)} and \m{u(g_1)=\alpha} is a scalar for some $g_1$ with \m{\cyc(g_1)=\Dg{i}{D}^{\#n}}. By (W4) there are
an inertial \m{g_0} above \m{\placesAboveG{D}\setminus\pd{\ED{D}}{\ell}} and $y$ with \m{R(g_0)=y}, and then
\m{u(g_0)=y}, \m{y\in\cO}, and \m{y^{j}-1} is a unit of~$\cO$ for \m{1\le j\le4\#n+2}.

\emph{A scalar and its conjugate.} If \m{u(g)=y'} is a scalar, then \m{y'} is a unit and
\m{u(g^{[-]})=\cyc(g)^{t-1}\,y'} (with \m{t-1} truncated): by the trace identity of a column, for every~$k$ the
traces of \m{u(k)\,u(g^{[-]})} and of \m{\cyc(g)^{t-1}y'\,u(k)} agree, and \m{u(\AbsGal{D})} spans all matrices, on
which the trace form is non-degenerate.

\emph{The clauses.} Take \m{\gamma=g_0}. (B1): $u$ is continuous (Lemma~5.7(b) of the first note), and it is trivial
on the inertia outside~$\Sigma$ because a column of level~$\emptyset$ is good outside \m{\placesAboveG{D}}. (B2): the
span was shown; the residue of~$y$ has order greater than \m{4\#n+2}. (B3): \m{u(g_0)=y=u(g_0^{[-]})}, because
\m{\cyc(g_0)=1} on inertia prime to~$\ell$, and \m{y\pm1} are units since \m{y^2-1} is. (B4): coboundaries are
crystalline cocycles, so they obey \m{\crysRels{u}}. (B5) holds with \m{g=g_1}, \m{\alpha'=\cyc(g_1)^{t-1}\alpha} and
\[
\beta=\sum_{v\in\placesAboveG{D}\setminus\pd{\ED{D}}{\ell}}v_{\cO}\bigl(\det(u(\gamma_v)-1)\bigr)+\#n\cdot v_{\cO}\bigl(\alpha-\cyc(g_1^{\mathsf c})\,\alpha'\bigr),
\]
where \m{\gamma_v} is the element that (W2) provides at the chosen place of~\m{\Qb} over~$v$. Indeed, let
\m{a\mid\ell^{M}}. At a place $v$ above~$\ell$ the local cocycles modulo~$a$ that obey \m{\crysRels{u}} obey the
relations of the crystalline cocycles of the local lattice, and their length is at most
\m{\rk_{\cO}\ZfAt{v}{u}\cdot\len(\cO/a)} (an input of the unit on Fitting ideals); by Lemma~\ref{l:inclass} and
Lemma~11.5 of the first note the rank is \m{\#n} plus the number of negative jumps at the embeddings read at~$v$.
Summing over~$v$, Lemma~4.4 of the first note and (R4) bound the total number of negative jumps by \m{\#n} times
the number \m{\fD{D}} of infinite places. At a place $v$ of \m{\placesAboveG{D}} prime to~$\ell$, (W2) and an input
of the unit on tame cocycles bound the local length by \m{\len((\cO/a)^n)} plus the length of the kernel of
\m{u(\gamma_v)-1}, which is at most the valuation of its determinant. Adding up gives (B5).

\emph{The size of~$\beta$.} Put \m{t'=\max(t,1)}. As \m{\cyc(g_1^{\mathsf c})=\cyc(g_1)} and
\m{\alpha'=\cyc(g_1)^{t-1}\alpha} with \m{t-1} truncated, \m{\alpha-\cyc(g_1^{\mathsf c})\alpha'=\alpha\,(1-\cyc(g_1)^{t'})}
exactly. Here $\alpha$ is a unit, and \m{\cyc(g_1)^{t'}-1=\Dg{i}{D}^{\#n\,t'}-1} has valuation at most
\m{\crudeUnit{\ell}{D}} by Lemma~5.7(a) of the first note, applied with the exponent \m{\#n\,t'}, which lies
between $1$ and \m{\Dr{D}\tD{D}} because \m{1\le\#n\le\Dr{D}} and \m{t'\le\tD{D}} for a row. By (W2) each \m{\vl(\det(u(\gamma_v)-1))} is at most
\m{\#n\,\crudeUnit{\ell}{D}}, and there are at most \m{4\fD{D}} places~$v$ (Lemma~5.5(b) of the first note). At the
chosen place~$w$, (W2) gives \m{\gamma_w=\tau_w} and the same bound for \m{\det(u(\tau_w^{[b]})-1)}, for both signs.
In normalised length the left side of the theorem is therefore at most
\m{(4\fD{D}+1+2)\,\#n\,\crudeUnit{\ell}{D}\le10\,\fD{D}\,\#n\,\crudeUnit{\ell}{D}}.
\end{proof}

\leanline{\lean{eqoneofmeminertiaAbove}\sep the triviality of a column on the inertia outside \m{\pd{\ED{D}}{\ell}\cup\placesAboveG{D}} and its level, used for (B1)}

\sect{Ribet's lemma to first order: proof of Theorem D}{s:ribet}
In this section $\Lambda_0$ is a commutative ring, $\mathbb{L}$ a commutative ring, and index sets may be finite
sets of any kind.

\begin{lemma}[Amitsur's formula for block matrices]\label{l:amitsur}
Let $\sigma$ and $N$ be finite sets, $\mathbb{L}$ a commutative algebra over a commutative ring~$\Lambda_0$,
$\mathfrak{A}$ a $\Lambda_0$-algebra (not necessarily commutative) and \m{\varrho\colon\mathfrak{A}\to M_\sigma(\mathbb{L})} a
homomorphism of $\Lambda_0$-algebras. Suppose that every coefficient of the characteristic polynomial of every
\m{\varrho(w)}, \m{w\in\mathfrak{A}}, lies in the image of~$\Lambda_0$. Let \m{P'\in M_N(\mathfrak{A})} and \m{k\in\NN}.
Then the $k$-th coefficient of the characteristic polynomial of the matrix
\m{\blockFlattenAlgEquiv{N}\bigl((\varrho(P'_{ij}))_{i,j}\bigr)} lies in the image of~$\Lambda_0$.
\end{lemma}
\begin{proof}
By Input~2 it suffices to check the property on elements that span \m{M_N(\mathfrak{A})}: the matrices with one
non-zero entry. For such a matrix the block image has, up to a power of~$X$, the characteristic polynomial of
the image of one entry of~$\mathfrak{A}$ (the trace of a product is insensitive to cyclic rotation), or the
polynomial \m{X^{\#N\#\sigma}}.
\end{proof}
\leanline{\lean{amitsurcharpolymemmatrix}}

\inp{Input 2 (Amitsur's formula, integrality consequence; background library, unit of the linear algebra of pseudocharacters)}{If
\m{\varrho\colon\mathfrak{A}\to M_\sigma(\mathbb{L})} is a homomorphism of $\Lambda_0$-algebras and $\mathfrak{A}$ is spanned
over~$\Lambda_0$ by elements whose images have characteristic polynomials with coefficients in the image
of~$\Lambda_0$, then every \m{\varrho(w)} has this property; and then every $w$ has an ``adjugate'': an element $v$
and \m{F'} in the image of~$\Lambda_0$ with \m{\det\varrho(w)=F'} and \m{\varrho(v)\varrho(w)=\varrho(w)\varrho(v)=F'}.
\emph{Reference} (second reading): G.~Chenevier, \emph{The $p$-adic analytic space of pseudocharacters of a
profinite group and pseudorepresentations over arbitrary rings}, §1, equation (1.5). \emph{Relation:} a weaker
form: the descent of the coefficients on a spanning family follows from the integral formula; no formula for
the coefficients is stated.}

\begin{lemma}[Ribet's construction: Cramer's rule twice]\label{l:cramer}
Let $A$ be a commutative ring, $\Gamma$ a group, $Y$ a set, $\sigma$ a non-empty finite set, $N$ a finite set,
$\Lambda_0$ a commutative ring, $\mathbb{L}$ a commutative ring that is an algebra over $A$ and over~$\Lambda_0$, and
$\mathcal{S}$ a $\Lambda_0$-subalgebra of the functions \m{Y\to\mathbb{L}}. Let \m{u=(u_x)_{x\in Y}} be homomorphisms of
$A$-algebras \m{u_x\colon A[\Gamma]\to M_\sigma(A)}, \m{\kappa\colon A\to\mathcal{S}} a ring homomorphism,
$\mathcal{J}$ an $\mathcal{S}$-submodule of the functions \m{Y\to\mathbb{L}}, \m{w\in A[\Gamma]^{N}},
\m{V\colon A[\Gamma]^{N}\to A[\Gamma]} a homomorphism of left \m{A[\Gamma]}-modules, and \m{\beta\colon Y\to\mathbb{L}} a
function. Suppose that $\kappa$ followed by the inclusion of~$\mathcal{S}$ is the structure map of the $A$-algebra
of functions; that for all \m{g\in\Gamma} and~$k$ the function \m{x\mapsto} ($k$-th coefficient of the characteristic
polynomial of \m{u_x(g)}, read in~$\mathbb{L}$) lies in~$\mathcal{S}$; and that \m{\beta(x)} is the image of
\m{\det u_x(V(w))} for all~$x$. Then \m{\beta\in\mathcal{S}}, and for every \m{h\in A[\Gamma]}, every
\m{M'\in M_N(A[\Gamma])} and every left ideal $\mathfrak{a}$ of \m{A[\Gamma]} such that each \m{V(M'_p)} (the image of the
$p$-th row) lies in the sum of \m{\mathfrak{a}\cdot V(w)} and of the left ideal generated by the $r$ with
\m{(x\mapsto\tr(u_x(r)u_x(g)))\in\mathcal{J}} for all~$g$, there is \m{f\in\mathcal{S}} with \m{f\beta\in\mathcal{J}} and
\[
f(x)=\det\bigl(\blockFlattenAlgEquiv{N}(u_x(M'))\bigr)\cdot\tr u_x(h)\qquad\text{for every }x\text{ with }u_x(m\,h)=0\text{ for all }m\in\mathfrak{a}.
\]
\end{lemma}
\begin{proof}
This is an outline of the formal proof.
Let \m{\rho'\colon\mathcal{S}[\Gamma]\to M_\sigma(Y\to\mathbb{L})} be the $\mathcal{S}$-algebra homomorphism with
\m{\rho'(g)(x)=u_x(g)}. By hypothesis and Input~2, all characteristic polynomials of~$\rho'$ have coefficients
in~$\mathcal{S}$. The adjugate of \m{V(w)} gives \m{\beta=\det\rho'(V(w))\in\mathcal{S}} and an element $\eta'$ with
\m{\rho'(\eta')\rho'(V(w))=\beta}. Write \m{V(M'_p)=r_p+m_pV(w)} with $r_p$ in the trace ideal and \m{m_p\in\mathfrak{a}}.
By Lemma~\ref{l:amitsur} the block matrix of $M'$ has an adjugate \m{W'} with entries in \m{\mathcal{S}[\Gamma]} and
determinant \m{F''\in\mathcal{S}}. Put \m{a'=F''-\sum_{p,q}w_q\,W'_{qp}\,m_p} and \m{f=\tr\rho'(a'h)}, an element
of~$\mathcal{S}$ (a coefficient of a characteristic polynomial). Where \m{u_x(mh)=0} for \m{m\in\mathfrak{a}}, the sum
does not contribute and \m{f(x)=F''(x)\tr u_x(h)}. Finally \m{a'V(w)=\sum_{p,q}w_qW'_{qp}\,r_p} by Cramer's rule for
\m{W'} and the linearity of~$V$, so \m{f\beta=\tr\rho'(\eta'\,h\,a'\,V(w))}-type sums of traces of multiples of the
$r_p$, which lie in~$\mathcal{J}$ by the definition of the trace ideal.
\end{proof}
\leanline{\lean{existsribetelementcramer}}

In the next lemma, for an $A$-linear map \m{\Phi\colon A^n\to\Lambda^{s}}, a function \m{z\colon\Gamma\to A^n}, a
submodule $\mathcal{N}$ of~$\Lambda^{s}$ and \m{p\in A}, the \emph{module of forms} \m{\mathfrak{F}(\Phi,z,\mathcal{N},p)}
consists of the functions \m{g\mapsto\Psi(z(g))} on~$\Gamma$, for the $A$-linear maps \m{\Psi\colon A^n\to A/p} that
vanish on \m{\Phi^{-1}(\mathcal{N})} (background library).

\begin{lemma}[a Fitting ideal contains a multiple of the product of the tie parameters]\label{l:block}
Let $A$ be a discrete valuation ring, $\Gamma$ a set, $J$ a finite totally ordered set, $n$ a finite set, $s$ a
finite subset of~$\Gamma$, \m{\Lambda=A[X_j:j\in J]}, \m{\Phi\colon A^n\to\Lambda^{s}} an $A$-linear map,
\m{z\colon\Gamma\to A^n} a function, $\mathcal{N}$ a $\Lambda$-submodule of~$\Lambda^{s}$, \m{p,C\in A},
\m{\tau,\phi\colon J\times\{+,-\}\to\Gamma} functions, $I$ a finite subset of~$J$ and \m{b\in\NN\cup\{\infty\}}.
Suppose that \m{p\ne0}; that \m{\HasTriangleCost{\mathfrak{F}(\Phi,z,\mathcal{N},p)}{\tau}{\phi}{I}{C}{b}}; that
\m{\Phi(pv)\in\mathcal{N}} for all~$v$; that \m{\Phi(z(g))-\mathbf{1}_g\in\mathcal{N}} for all \m{g\in s}; that
\m{\Phi(z(\tau_k))-X_j\,\Phi(z(\phi_k))\in\mathcal{N}} for all \m{k=(j,b')}; and that for all \m{j\in I} and \m{i<j},
either \m{i\in I} or \m{X_i\Lambda^{s}\subseteq\mathcal{N}}. Then there is \m{G\in\Lambda} with
\[
v_A(G(0))\le b+(\#I+\#s)\,v_A(C)\qquad\text{and}\qquad G\cdot\prod_{j\in I}X_j\in\operatorname{Fitt}(\mathcal{N}).
\]
\end{lemma}
\begin{proof}
Let $\Gamma_0$ be the set of the triangle cost. Three facts on Fitting ideals are used (background library):
\emph{(cost)} for a module $\mathcal{V}$ of functions \m{\Gamma\to A/p} there are \m{\mathcal{K}\supseteq\mathcal{N}}
containing \m{\Phi(v)} whenever all forms in~$\mathcal{V}$ kill~$v$, and \m{e\in A} with
\m{v_A(e)\le\len(\mathfrak{F}/\mathfrak{F}\cap\mathcal{V})} and \m{e\operatorname{Fitt}(\mathcal{K})\subseteq\operatorname{Fitt}(\mathcal{N})};
\emph{(triangle)} if \m{r_jv_j} lies in a module plus the span of the earlier $v_i$, for all \m{j\in I}, then
\m{\prod_{j\in I}r_j} multiplies the Fitting ideal of the module with the $v_j$ adjoined into that of the module;
\emph{(constants)} \m{\operatorname{Fitt}(\mathcal{N}+\mathfrak{b}\Lambda^{s})\subseteq\operatorname{Fitt}(\mathcal{N})+\mathfrak{b}}.
First apply (cost) to the functions vanishing on~$\Gamma_0$: this gives \m{\mathcal{K}} and $e$. Then (triangle)
with \m{r_j=CX_j} and \m{v_j=\Phi(z(\phi_{(j,+)}))-\Phi(z(\phi_{(j,-)}))}: by the hypothesis on the ties,
\m{X_jv_j} is congruent modulo~$\mathcal{N}$ to the difference of the \m{\Phi(z(\tau_{(j,\pm)}))}, and (C1), read
through the duality between forms and membership (a vector lies in a module containing \m{p\,\Phi(A^n)} if every
form vanishing on the module kills it), says that $C$ times this difference lies in the module generated by the
earlier ones; the hypothesis on \m{i\notin I} disposes of the earlier indices outside~$I$. Then (cost) again, for
the $C$-torsion, with a constant $f'$; by (C2) \m{v_A(e)+v_A(f')\le b}. The matrix \m{C\cdot\mathbf{1}} of size
\m{\#s} has its rows in the last module, so \m{C^{\#s}} lies in its Fitting ideal, and (constants) for the ideal of
polynomials without constant term gives an element $d_4$ of the Fitting ideal of the middle module with constant
term \m{f'C^{\#s}}. Put \m{G=e\,C^{\#I}d_4}.
\end{proof}
\leanline{\lean{existsmemfittingIdealblock}}

\begin{proof}[Proof of Theorem~D]
This is an outline. The formal proof has 150 lines; here is its architecture, with the places where each hypothesis enters. Let
\m{\Lambda=A[X_j:j\in J]} and work in the group algebra \m{\Lambda[\Gamma]}; let \m{\chi_x}, \m{\psi_x} and \m{\eta_e} be
the algebra homomorphisms from \m{\Lambda[\Gamma]} to matrices over~$\Lambda$ and to~$\Lambda$ induced by \m{u_x},
\m{\varrho_x} and \m{\xi_e}. On~$\mathcal{X}$, \m{\chi_x} is the block sum of \m{\psi_x} and the \m{\eta_e}, by~(a).

\emph{Two elements of the group algebra.} By (b) the element \m{k_0=\prod_{x\in\mathcal{X}}(y_x-1)^{-1}(y_x-\gamma_x)}
has \m{\eta_e(k_0)=1} for all~$e$ and \m{\psi_x(k_0)=0} for all \m{x\in\mathcal{X}}. By (d) there are $q_e$ with
\m{\eta_{k}(q_e)=C} for \m{k=e} and $0$ otherwise (the interpolation polynomial at the values \m{\xi_k(a)}, evaluated
at~$a$).

\emph{Coordinates and relations.} Choose a finite set \m{s_0\subseteq\Gamma} such that the \m{u(g)-1}, \m{g\in s_0},
span all \m{u(g)-1} over~$A$; index coordinates by \m{W\sqcup s_0\sqcup s} (the elements \m{\delta_v}, $s_0$ and~$s$),
and write each \m{u(g)-1} as a combination \m{z(g)} of the coordinates. Let $V$ be the \m{\Lambda[\Gamma]}-linear map
sending a coordinate vector to the corresponding combination of the \m{g-1}, and \m{\omega} the vector of the
normaliser of~(e). Let $\mathcal{N}$ be the module of the coordinate vectors whose image under $V$ lies in the
sum of the trace ideal of~$\mathcal{J}$ and of \m{\ker(\eta_0)\cdot V(\omega)}. It contains: the cocycle relations
\m{z(gg')-z(g)-g\,z(g')}; the linear relations among the \m{u(g)-1}; the ties
\m{z(\tau_k)-X_j\varepsilon_j(\phi_k)z(\phi_k)} and all multiples of~$p$ (by (c)); and \m{(g-1)\omega}.

\emph{The block.} Evaluating the variables at~$0$ and applying \m{\psi_\pi}, a vector killed by~$\mathcal{N}$ becomes
an element of \m{\obeyingCocycles{u}{\varrho_\pi}{p_0}} that vanishes at the \m{\tau_k} and, by the normaliser
relation and the unit \m{y_\pi-1}, at~\m{\gamma_\pi}. By (e) these have length at most \m{b_0<v_A(p_0)}, and since
\m{\varrho_\pi} spans all matrices there is a square matrix $B_0$ with rows in~$\mathcal{N}$ and
\m{v_A(\det\psi_\pi(B_0)|_{X=0})\le b_0} (an input on Fitting ideals).

\emph{The characters.} Fix~$e$. The forms on the image of~$\mathcal{N}$ under~\m{\eta_e} give cocycles in
\m{\obeyingCocycles{u}{\xi_e}{p}} that vanish at the \m{\delta_v}; and the third condition of (g) holds for them:
if such a cocycle were of the form \m{(\xi_e(g)-1)m} without being killed by \m{\xi_e(a)-1}, the relation
\m{(g-1)\omega\in\mathcal{N}} together with (d) and \m{b_0<v_A(p_0)} would make the unit \m{y_\pi-1} a sum of two
non-units. So (g) applies: it gives the triangle cost and, by the detection on~$s$, elements of~$\mathcal{N}$
that eliminate the coordinates of~$s_0$. Let \m{\mathfrak{q}} be the ideal of~$\Lambda$ spanned by the monomials
that do not divide \m{\prod_jX_j}. Lemma~\ref{l:block} is applied with \m{I=\{j:\mathrm{ty}(j)=e\}} and with the
image of~$\mathcal{N}$ enlarged by \m{(\mathfrak{q}:\prod_{\mathrm{ty}(i)<e}X_i)\,\Lambda^{s}}; the variables of the
earlier types are not inverted: for \m{\mathrm{ty}(i)<e} the multiples of~$X_i$ lie in this colon module, which is
the alternative ``\m{X_i\Lambda^{s}\subseteq\mathcal{N}}'' of the lemma, and for \m{i<j} with \m{\mathrm{ty}(j)=e} the
monotonicity of~\m{\mathrm{ty}} gives \m{\mathrm{ty}(i)=e} or \m{\mathrm{ty}(i)<e}. The lemma gives
\m{\Gamma_e\in\Lambda} with
\[
v_A(\Gamma_e(0))\ \le\ b_e+\bigl(\#\{j:\mathrm{ty}(j)=e\}+\#s\bigr)\,v_A(C)
\]
and \m{\Gamma_e\prod_{\mathrm{ty}(j)=e}X_j} in the Fitting ideal of the enlarged module; it is, up to the colon
ideal, a combination of determinants of matrices $B_e$ with rows in the image of~$\mathcal{N}$.

\emph{Assembly.} For each choice of such matrices form the matrix $M'$ with the rows \m{o_v(\delta_v)} at the
coordinates of~$W$ (in~$\mathcal{N}$ by (f)), the eliminating rows, and \m{\sum_eq_e\cdot}(the rows of $B_e$). On the
characters its determinant is \m{\prod_vo_v(1)\cdot C^{\#s}} times \m{\det B_e}. Apply Lemma~\ref{l:cramer} to
\m{h=k_0\cdot(\ldots)\cdot q_0} and \m{B_0+k_0(M'-B_0)}, with the ideal \m{\mathfrak{a}=\ker\eta_0} (for
\m{x\in\mathcal{X}} one has \m{u_x(mh)=0} for \m{m\in\mathfrak{a}}, because \m{\xi_0=1} is a constituent of~$u_x$ and
$k_0$ kills the block): it gives $f$ with \m{f\beta\in\mathcal{J}}, and summing over the choices, an $f$ whose
value at \m{x\in\mathcal{X}} is, modulo monomials not dividing \m{\prod_jX_j},
\[
\det\psi_x(B_0)\cdot C\cdot\Bigl(\prod_vo_v(1)\,C^{\#s}\Bigr)^{t}\cdot\prod_e\Gamma_e\cdot\prod_jX_j .
\]
So no monomial properly dividing \m{\prod_jX_j} occurs, and the coefficient $c$ at~$\pi$ has
\[
v_A(c)\ \le\ b_0+\sum_e v_A(\Gamma_e(0))+t\,\Bigl(\sum_vv_A(o_v(1))+\#s\,v_A(C)\Bigr)+v_A(C).
\]
The numbers \m{\#\{j:\mathrm{ty}(j)=e\}} add up to~\m{\#J}, so the bounds for the \m{\Gamma_e(0)} add up to
\m{\sum_eb_e+(\#J+t\,\#s)\,v_A(C)}, and the total is exactly the bound of the theorem. That \m{\beta\in\mathcal{S}} is part of Lemma~\ref{l:cramer}, and \m{\beta(x)=0} on~$\mathcal{X}$ because the
matrix of the normaliser has a zero column at the constituent~\m{\xi_0=1}.
\end{proof}

\sect{The cost of a Tate twist: proof of Theorem E}{s:tate}

\begin{lemma}[symmetrisation gives even classes]\label{l:even}
Let $E$ be a CM field as in the conventions, $\cO$ a coefficient ring, \m{e,\nu} natural numbers, $S$ a set of
finite places of~$E$ and \m{c\colon\GE{E}\to\cO/\ell^{\nu}} a function. Suppose that \m{e\ne0}, that
\m{c\in\Sel{\cO(-e)}{\pd{E}{\ell}\cup S\cup\cj{S}}{\emptyset}{\cO/\ell^{\nu}}}, and that $c$ vanishes at every element of a
decomposition group above~$\ell$ with \m{\cyc=1}. Then \m{g\mapsto c(g)+c(g^{\mathsf c})} lies in
\m{\evenClasses{\cO}{S}{\nu}{e}}.
\end{lemma}
\begin{proof}
Conjugation by~$\mathsf{c}$ preserves the cyclotomic character, the set \m{\pd{E}{\ell}\cup S\cup\cj{S}} and the set of
the elements of decomposition groups above~$\ell$ with \m{\cyc=1}. So \m{c\circ(\cdot)^{\mathsf c}} is a cocycle of the
same kind, and the sum is visibly invariant.
\end{proof}
\leanline{\lean{addconjGalmemevenClasses}}

\begin{proof}[Proof of Theorem~E]
Take for $\Gamma_0$ the set of the elements of inertia groups above~$\ell$ if \m{e=0}, and of the elements of
decomposition groups above~$\ell$ with \m{\cyc=1} otherwise. Write \m{I=\{j:\mathrm{ty}(j)=e\}}; it is empty for
\m{e=0}.

\emph{(C1).} Let \m{j\in I} and \m{z\in Z} vanish on~$\Gamma_0$, with \m{z(\tau_i)=z(\tau_i^{[-]})} for \m{i<j}, and
\m{z(\phi_i)=z(\phi_i^{[-]})} for \m{i<j} in~$I$. By Lemma~\ref{l:even}, \m{z^{+}=z+z\circ(\cdot)^{\mathsf c}} is an even
class of type \m{e=\mathrm{ty}(j)}. For a cocycle of \m{\cO(-e)} and an element $g$ with \m{\cyc(g)\equiv1},
\m{z(g^{-1})=-z(g)}; with Lemma~7.2 of the first note this turns the hypotheses into \m{z^{+}(\tau_i)=0} and
\m{z^{+}(\phi_i)=0}. Clause (T4) gives \m{\ell^{h}z^{+}(\tau_j)=0}, that is \m{\ell^{h}z(\tau_j)=\ell^{h}z(\tau_j^{[-]})}.

\emph{(C2), first term.} A function of $Z$ that obeys \m{\crysRels{\cO(-e)}} vanishes on~$\Gamma_0$. Indeed, an
element of~$\Gamma_0$ is conjugate to an element~$s$ of the decomposition group of a chosen place above~$\ell$
with \m{\cyc(s)=1}, inertial if \m{e=0}; for \m{e\ne0} such an $s$ need not be inertial, and on the decomposition
group the module \m{\cO(-e)} is the character \m{\cyc^{-e}}, not the trivial one. A crystalline cocycle of
\m{\cO(-e)} vanishes at such an~$s$: for \m{e>0} it is a coboundary after inverting~$\ell$ (the finite part of
\m{H^1} of a negative Tate twist vanishes), hence it vanishes where \m{\cyc^{-e}=1}; for \m{e=0} it is unramified.
This is an input of the unit on the Bloch--Kato condition; the second reading relates it to Example~3.9 of Bloch
and Kato (a weaker form). The cocycle identity transports the vanishing along the conjugation.
Hence \m{\len(Z/\{z:z|_{\Gamma_0}=0\})\le\len(Z/Z\cap\Ob{\crysRels{\cO(-e)}}{\cO/\ell^{\nu}})}.

\emph{(C2), second term.} Let \m{Z'} be as in (C2). By Lemma~7.3 of the first note its elements, which vanish at
all \m{\tau_k}, are unramified outside~$\ell$. If \m{e=0} they lie in \m{\Sha^{\mathrm{odd}}(\ell,E,\cO,0,\nu)}, and
\m{\len(Z'/Z'[\ell^{h}])\le\len Z'}. If \m{e\ne0}, the map \m{z\mapsto z\circ(\cdot)^{\mathsf c}-z} sends $Z'$ into
\m{\Sha^{\mathrm{odd}}(\ell,E,\cO,e,\nu)}, and its kernel is killed by~\m{\ell^{h}}. Indeed, let $z$ be in the
kernel. Then \m{z^{+}=2z} is an even class unramified outside~$\ell$, and it vanishes at the $\phi_j$ of type~$e$:
the hypothesis on the $\phi$-pairs gives \m{z(\phi_j)=z(\phi_j^{[-]})=-z(\phi_j)}, that is \m{2z(\phi_j)=0} (and
not \m{z(\phi_j)=0}, which is why (T5) is applied to~$z^{+}$ and not to~$z$). By (T5) for the even class~$z^{+}$,
\m{\delta(\gamma)\cdot4z(g)=\delta(g)\cdot4z(\gamma)} for all~$g$: the element \m{4\delta(\gamma)\,z} of~$Z$ is of the
form \m{g\mapsto(\cyc(g)^{-e}-1)m}. Hypothesis (ii) then gives \m{4\delta(\gamma)^2z=0}, and \m{4\delta(\gamma)^2}
divides~\m{\ell^{h}} by (T5). So \m{\len(Z'/Z'[\ell^{h}])\le\len\Sha^{\mathrm{odd}}(\ell,E,\cO,e,\nu)}, and (iii) gives the bound~$b$.
\end{proof}

\sect{Where the results are used}{s:use}
The results are used in the third note of the unit, \emph{A fit family of twists exists}. Theorem~A, with
Theorem~C for its hypothesis, gives there the one twist for each pair of a prime and an isomorphism that makes
all Selmer length bounds hold at once. Theorem~D is the engine of the bound for the denominator of the
projection to theta lifts: it is applied to the Galois representations of a base level, with the lifts as the
set~$\mathcal{X}$; its hypothesis (g) is supplied by Theorem~E for the characters \m{\cyc^{-e}} of the string, and
its hypothesis (e) by Theorem~A of this note through the Selmer length bound. Lemmas \ref{l:tametype},
\ref{l:inclass} and~\ref{l:obeycup} enter the estimate of the crystalline losses above~$\ell$. Theorem~B and the
lemmas of §§\ref{s:roots}--\ref{s:pencils} are used only inside this note.

\sect{Notes on the formal text}{s:notes}
\begin{enumerate}[label=\textup{\arabic*.},leftmargin=1.6em,itemsep=5pt]
\item \emph{Letters.} Besides the changes of the first note: the formal text writes {\leanfont O} for the rings
called $A$, $A_0$, $A'$ here, {\leanfont L} for the algebra~$\mathbb{L}$ and for sets of relations, {\leanfont S}
for the subalgebra~$\mathcal{S}$, {\leanfont 𝒥} for~$\mathcal{J}$, {\leanfont 𝓧} for~$\mathcal{X}$, {\leanfont ι} for the
index set~$\mathcal{I}$, {\leanfont G} for the group~$\Gamma$, {\leanfont R} for the rings $A$ of
Lemma~\ref{l:obeycup} and $\Lambda_0$ of Lemma~\ref{l:amitsur}, {\leanfont Λ} for the algebra~$\mathfrak{A}$,
{\leanfont c} for the involution~$\mathsf{c}'$ of Lemma~\ref{l:mirror}, {\leanfont I} for its ideal~$\mathfrak{a}$,
{\leanfont n} for the function~$\theta$ of Theorem~D, {\leanfont L} for the set~$\Gamma_0$ of
Definition~\ref{d:tricost}, and {\leanfont P₀} for the helper prime~$\fq_0$.
\item \emph{Arguments.} The arguments were read from the text of the files. For seven declarations (the four
lemmas on pencils, Theorem~D, Theorem~C and Lemma~\ref{l:inclass}) the dossier gives a list reconstructed by hand,
because the files declare the same variables twice or the statement contains a local definition; none was
confirmed by Lean's own answer. The program reads the argument of the function of Definition~\ref{d:normaliser}
(the group element) as an argument, and so do we; in Theorem~C it reads the quantifier ``for all large~$\cO$'' as
two arguments. The instance ``discrete valuation ring'' is not used by Definition~\ref{d:tricost}, and
``group'' not by Definition~\ref{d:trid}; they are not listed there.
\item \emph{Statements with local definitions.} Lemmas \ref{l:pointchar}(b), \ref{l:drop}, \ref{l:unseen},
\ref{l:inclass}, \ref{l:obeycup} and Theorem~B have a local definition inside their formal statement; they were
read to their true end in the source.
\item \emph{Theorem~C.} The finiteness of the set~$\Sigma$ is an argument of the formal statement with a default
proof; the set enters the conclusion as a finite set made from that proof.
\item \emph{Theorem~A.} The hypotheses \m{S\subseteq S_0} and \m{z\in\inab{S_0\setminus\pd{E}{\ell}}} stand in a
declaration of variables before the statement and are included in it. The same $I$ and $\psi$ serve every $d$
and every~$i$; only the coefficient ring $\cO'$ and \m{\psi'} depend on them. The bounds $B_i$ are extended
non-negative reals, $c$ and $\eta$ non-negative reals.
\item \emph{Theorem~D.} The index set of the characters is \m{\{0,\dots,t-1\}} with \m{t\ne0}; the map
\m{\mathrm{ty}} is a monotone map whose target the formal text leaves to be inferred. Hypothesis (g) quantifies
over all submodules~$Z$; its conclusion is used for the modules of forms of the proof. The formal bound is
\m{b_0+(\sum_eb_e+t(\sum_vv_A(o_v(1))+2\#s\,v_A(C))+(\#J+1)v_A(C))}; the outline above reproduces it: each
character contributes \m{(\#\{j:\mathrm{ty}(j)=e\}+\#s)\,v_A(C)} through Lemma~\ref{l:block}, the assembly
\m{t\,\#s\,v_A(C)} and one more \m{v_A(C)}.
\item \emph{Leaves of the structured proofs.} Two leaf steps of the dossier quote nine lines (one long
simultaneous choice in Lemma~\ref{l:rootmany}, one chain of rewritings in Lemma~\ref{l:unseen}); the bound of the
workspace is eight.
\item \emph{Outlines.} The proofs of Lemmas \ref{l:rootone}, \ref{l:mirror}, \ref{l:unseen}, \ref{l:obeycup},
\ref{l:cramer} and of Theorems A, C and~D are outlines of long formal computations, as said at the head or at
the end of each. For them I compared the hypotheses used at each step with the formal text but did not re-derive
the identities. Two points were re-derived in the revision of this note: in Theorem~A the roots of unity are
chosen for the polynomials \m{X^{d}-1}, \m{1\le d\le m}, as in the formal proof; and in Theorem~D the multiples of
\m{v_A(C)} of the outline add up to those of the statement (note~6).
\item \emph{Input~1.} The formal statement of the relations at split primes is long (it has a local definition
and was read to its end by the second reading); it is stated above for a number field, with all its hypotheses,
because the third note uses it over a totally real field. Its hypotheses on the local conditions, on the complex
conjugations and on the chosen places stand in a declaration of variables before the formal statement and are
included in it. The write-up has not yet read its proof, and no printed source for this form is known to the
catalogue; the same holds for the count of trivialised cocycles that supplies the hypothesis of its part~(a).
\item \emph{The tilting lemma.} The formal proof compares lengths: with the notation of the proof above, it shows
\m{(s\Phi_0-\Phi_1)(Z)\supseteq\Phi_0(Z)} and concludes by additivity of the length. The map $\Phi_1$ does not factor
through~$\Phi_0$ on~$Z$ in general (for \m{A=\Zl}, \m{\mathfrak{a}=(\ell)}, \m{\#m=1}, \m{B=Y=(\Zl/\ell)^2},
\m{\Phi_0(x,y)=x}, \m{\Phi_1(x,y)=y} one has \m{Z=B}, and the kernel of \m{s\Phi_0-\Phi_1} is not \m{\ker\Phi_0}).
\end{enumerate}

\subhead{How this note was written.} From the dossier of the unit (its second chunk): the statements in words and
the structured proofs, which a program compared line by line with the source (\UnitSteps{} leaf steps in the
\UnitProofs{} proofs of the whole unit). I read the eight Lean files of these declarations completely, and for
the long proofs the reader's step-by-step rendering beside them. The statements of this note were compared with
the Lean statements argument by argument, and the order of the quantifiers is that of the formal text (in
Theorem~A: first $I$ and~$\psi$, then $d$, then~$i$, then the coefficient ring; in Theorem~B: first $J_0$,
then~$N$). The proofs of Lemmas \ref{l:rootmany}, \ref{l:linepoint}, \ref{l:fitting}, \ref{l:selz}, \ref{l:pointchar},
\ref{l:drop}, \ref{l:tilting}, \ref{l:tametype}, \ref{l:inclass}, \ref{l:amitsur}, \ref{l:block}, \ref{l:even} and of
Theorems B and~E follow the formal ones; the others are outlines (note~8). I did not open the proofs of the inputs. The reference
of Input~2 and its place are those of the second reading of this write-up; I did not check them in the
literature.

\head{Appendix A. Concordance with the Lean artifact}
{\fontsize{8.6}{11.6}\selectfont\renewcommand{\leanfont}{\ttfamily\fontsize{7.7}{10}\selectfont\color{leancol}}
\setlength{\LTleft}{0pt}\setlength{\LTright}{\fill}\setlength{\tabcolsep}{0pt}
\begin{longtable}{@{}p{30mm}p{56mm}p{43mm}r@{}}
\textbf{In this note} & \textbf{Lean declaration} & \textbf{File} & \textbf{Lines}\\\midrule\endhead
\ConcordanceRows
\end{longtable}}

\subhead{Two statements as they are written in the artifact.}
\quoted{existsanticyclotomictwist}
\quoted{hasTriangleCosttateTwist}
In these texts {\leanfont IsBalanced}, {\leanfont IsFamily}, {\leanfont IsAuxTwist}, {\leanfont IsPrimeTriangle} and
{\leanfont HasTriangleCost} are the predicates written above with the macros of the same names; {\leanfont
idOrConjInv b z} is \m{z^{[b]}}, {\leanfont unionWithConj S I} is \m{\unionWithConj{S}{I}}, {\leanfont twistRep}
forms \m{\psi'\cdot\varrho_i}, {\leanfont mapRel} carries relations to the larger ring, {\leanfont .single (γ i) 1}
is \m{\delta_{\gamma_i}}, {\leanfont tateTwist ℓ E O e} is \m{\cO(-e)}, {\leanfont shaOdd} is \m{\Sha^{\mathrm{odd}}}, and
{\leanfont relLength S T} is the length of \m{S/S\cap T}.

\head{Notation}
\begingroup\fontsize{8.6}{10.9}\selectfont\raggedright
\setlength{\columnsep}{9mm}\setlength{\parskip}{0pt}
\begin{multicols}{2}
\nt{$\ell$, $\Ql$, $\Zl$, $\Qlb$, $\Zlb$, $\Cl$, $\Gl$, $\iota$}{a prime and the $\ell$-adic objects; the fixed isomorphism \m{\CC\to\Qlb}}
\nt{$\Qb$, $\mathsf{c}$, $g^{\mathsf c}$, $g^{[b]}$}{an algebraic closure of~$\mathbb{Q}$; complex conjugation; \m{\mathsf{c}g\mathsf{c}}; $g$ or \m{(g^{\mathsf c})^{-1}}}
\nt{$E$, $\GE{E}$, $f_E$}{a CM field as in the conventions; \m{\Gal(\Qb/E)}; its number of infinite places}
\nt{$P$, $D$, $\ED{D}$, $\fD{D}$, $\Dg{i}{D}$, $\gD{D}$, $\lvl{D}$, $\Dr{D}$, $\tD{D}$, $\ordch{D}$}{parameters, data and its numbers, as in the first note}
\nt{$w$, $v$, $\fq$, $\fq_0$, $\cj{w}$, $\Nm(w)$, $\fP$, $\fP_w$, $G_{\fP}$, $I_{\fP}$, $\tau_{\fP}$, $\tau_w$}{places, norms, decomposition and inertia groups, tame generators}
\nt{$\pd{E}{M}$, $\inab{S}$, $\cyc$, $\cyc_M$}{places containing~$M$; inertia above~$S$; cyclotomic characters}
\nt{$\cO$, $\cO'$, $\varpi$, $\cO(-e)$, $\vl$, $v_\ell$, $\nlen$, $\len$}{coefficient rings, a uniformiser, Tate twists, valuations and lengths}
\nt{$A$, $A_0$, $A'$, $A_1$, $\Lambda$, $\Lambda_0$, $\mathbb{L}$, $\mathfrak{A}$, $B$, $v_A$}{rings, algebras, a module, the valuation of a discrete valuation ring}
\nt{$\Gamma$, $\Gamma'$, $\Gamma_0$, $j$, $\eta$}{groups; a subset; the homomorphism to~\m{\Gl}; a homomorphism (also an idele character and a real number)}
\nt{$M_n(A)$, $\GL_n$, $\#n$, $\tr$, $\det$, $\rk$}{matrices, cardinality, trace, determinant, rank}
\nt{$\theta$, $\delta_g$, $\Ob{L}{B}$, $\Rel{Z}$, $\Sel{\varrho}{S}{L}{B}$}{relations, the indicator relation, obeying functions, relations of a module, Selmer modules}
\nt{$\Zf(L)$, $L^{\vee}(1)$, $L^{\perp\mathcal{Q}}$, $\mathcal{Q}(F,K,U)$, $\Sha^{\mathrm{odd}}$}{crystalline cocycles, twisted dual, cup-orthogonal, classes, the module of §2.1}
\nt{$\operatorname{Fitt}(\mathcal{N})$, $\mathfrak{F}(\Phi,z,\mathcal{N},p)$}{Fitting ideal; module of forms}
\nt{\m{\obeyingCocycles{u}{\varrho}{a}}, \m{\crysRels{u}}, \m{\matrixLattice{\varrho}}, \m{\coeffEmb{\cO}{g}}, \m{\complexEmbOf{e}{g}}}{Definition~\ref{d:recall}(1),(2)}
\nt{\m{\AbsGal{D}}, \m{\placesAboveG{D}}, \m{\badPlaces{D}{\ell}}, \m{\crudeUnit{\ell}{D}}}{Definition~\ref{d:recall}(3)}
\nt{\m{\IsColumn{D}{t}{\HT}{S}{R}}, \m{\IsRow{D}{t}{\HT}{R}}}{columns and rows (Definition~\ref{d:recall}(4))}
\nt{\m{\AuxPrimeData{E}{q}}, \m{\IsAnticyclotomicAt{T}{\fq}{l}}, \m{\IsAuxiliary{S}{I}}, \m{\IsFamily{S}{\nu}{\mathrm{val}}}}{Definition~\ref{d:recall}(5),(6)}
\nt{\m{\evenClasses{\cO}{S}{\nu}{e}}, \m{\IsNewPlaces{\nu}{w}{\tau}{\phi}}, \m{\withConjInv{\tau}{k}}}{Definition~\ref{d:recall}(7)}
\nt{\m{\IsPrimeTriangle{\nu}{w}{\tau}{\phi}{a}{t}{\mathrm{ty}}{h}{c}}}{prime triangles (Definition~\ref{d:recall}(8))}
\nt{\m{\IsAuxTwist{S}{Q}{\psi}{c}}, \m{\unionWithConj{S}{J}}, \m{\twistAtPoint{h}{y}{e}{\varrho}{g}}, \m{\reduceModPow{l}{M}{\varrho}}, $y^{h}$}{Definition~\ref{d:recall}(9),(10)}
\nt{\m{\IsBalanced{\ell}{S}{\varrho}{L}{\gamma}{\beta}{M}}}{balanced representations (Definition~\ref{d:recall}(11))}
\nt{\m{\IsPencil{\ell}{S_0}{U}{W}{N}{\fq_0}{J}{x}}, \m{\tieRelations{J}{x}{i}}}{pencils and ties (Definition~\ref{d:recall}(12))}
\nt{\m{\evalConjPair{g}}, \m{\blockFlattenAlgEquiv{N}}, \m{\ribetNormaliser{\gamma}{a}{y}{p}{\theta}{g}}}{Definitions \ref{d:evpair}--\ref{d:normaliser}}
\nt{\m{\HasTriangleCost{Z}{\tau}{\phi}{I}{C}{b}}, \m{\IsTraceIdeal{\mathcal{J}}{u}{p}{\tau}{\phi}{k}{\varepsilon}}, \m{\IsRibetElement{\mathcal{J}}{\beta}{\pi}{c}}}{Definitions \ref{d:tricost}--\ref{d:rib}}
\nt{$R$, $R'$, $u$, $\varrho$, $\varrho_i$, $\varrho_x$, $\varrho_y$, $\psi$, $\psi'$, $\chi$, $\xi_e$, $\mu_{k,t}$, $\nu$}{representations, lattices and characters ($\nu$ also a natural number)}
\nt{$\HT$, $\HT_w$, $\Lambda_w$, $e_w$, $\mathcal{R}$, $\mathfrak{M}$}{types; the local lattice and chart at~$w$; a property; a multiset}
\nt{$l=(\fP,\gamma,\phi,h)$, $l_v$, $h_v$, $\gamma_v$, $\phi_v$, $\gamma_l$, $\phi_l$, $s_{v,i}$, $\mathbf{1}_i$, $\mathbf{1}_g$}{data at auxiliary primes; slopes; indicator vectors}
\nt{$I$, $J$, $J_0$, $J_1$, $S$, $S_0$, $T$, $W$, $Q$, $P'$, $Q'$, $\Sigma$, $\mathcal{X}(T)$}{sets of places and index sets; bounds for the number of primes; the Selmer modules of Lemma~\ref{l:unseen}; $W$ also a neighbourhood, $T$ a list of residues, $Q$ a polynomial}
\nt{$\mathcal{I}$, $\kappa$, $Y$, $X$, $\mathcal{X}$, $\pi$, $\sigma$, $n$, $n_i$, $r$, $d$, $m$, $N$, $M$, $k$, $q$, $t$, $e$, $i$, $b$}{index sets and natural numbers; $X$ also a variable, $\pi$ a point, $\kappa$ also a ring homomorphism, $e$ also a ring homomorphism, $b$ also a bound or a sign}
\nt{$\beta$, $\beta_i$, $B_i$, $b_0$, $b_e$, $c$, $\eta$, $\delta$, $C$, $p$, $p_0$, $a$, $\alpha$, $\alpha'$, $s$}{bounds, real numbers and ring elements ($a$, $s$ also group elements or a finite set, $c$ also a cocycle or a ring element, $\delta$ also the function of (T5))}
\nt{$\gamma$, $\gamma_i$, $\gamma_x$, $z$, $g$, $g'$, $g_0$, $g_1$, $h$, $\tau$, $\phi$, $\zeta$, $u$, $y$, $y_v$, $y_x$, $y_0$}{elements of groups, roots of unity and points ($u$ also a lattice, $y$ also a cocycle, $h$ also an element of a group algebra)}
\nt{$F$, $F_k$, $F_{k,d}$, $F_i$, $G$, $G_i$, $D_n$, $V_{k,I}$, $\mathrm{val}_k$, $a_{k,i}$, $\rho_j$, $\|F\|_\rho$, $\Pi$, $Q_k$, $\eta_j$}{polynomials and Gauss norms of §\ref{s:roots} and of the proof of Theorem~A; $F$ also a filter, $G$ also the element of Lemma~\ref{l:block}}
\nt{$\mathsf{c}'$, $H$, $U$, $U'$, $U_i$, $\mathfrak{a}$, $x_w$, $\tau=[g_0,\gamma]$, $Y$}{objects of the mirror lemma and its use ($U$ also a function to finite sets, $Y$ also a matrix or a submodule)}
\nt{$\Phi_0$, $\Phi_1$, $\Phi$, $\Psi$, $\psi$, $\sigma'$, $\mathcal{M}$, $\bar v$, $Z$, $Z'$, $Z'[C]$}{linear maps and modules of §§\ref{s:pencils}, \ref{s:ribet}; the free module of the tilting lemma; reduction modulo~$\mathfrak{a}$; the $C$-torsion}
\nt{$K$, $\GE{K}$, $X$, $L_v$, $N'$, $U'$, $Y(\Pi')$, $\Pi'$, $\iota'$, $v_i$, $\phi_i$, $\gamma_i$, $\varpi'$, $T'$, $Z^1_v$, $\mathcal{R}_{G'}$, $G'$, $n_i$, $v'_j$, $\phi'_j$, $\delta_j$}{objects of Input~1: the field, the module, local conditions, the relaxed Selmer group, split places, relation families}
\nt{$\mathfrak{q}$, $t'$, $D_d$}{the ideal of the monomials not dividing \m{\prod_jX_j}; \m{\max(t,1)}; the polynomial \m{X^{d}-1}}
\nt{$\mathcal{S}$, $\mathcal{J}$, $\varepsilon_j$, $\epsilon$, $\mathrm{pr}_1$, $\mathrm{ty}$, $\delta_v$, $o_v$, $\omega$, $k_0$, $q_e$, $s_0$, $B_0$, $B_e$, $\Gamma_e$, $M'$, $P'$, $W'$, $F'$, $F''$, $a'$, $\eta'$, $\rho'$, $r_p$, $m_p$, $w$, $V$, $f$}{objects of §\ref{s:ribet}: subalgebra and submodule, polynomials, the bijection, maps, group-algebra elements, matrices}
\nt{$\mathcal{N}$, $\mathcal{K}$, $\mathcal{V}$, $\mathcal{R}$, $\mathfrak{b}$, $e$, $f'$, $d_4$, $r_j$, $v_j$, $L$, $K$, $\Lambda^{s}$}{modules of relations, ideals and constants; $L$ also a set of relations or a lattice, $K$ also a property of polynomials}
\nt{$z^{+}$, $\delta(g)$, $\chi_x$, $\psi_x$, $\eta_e$}{the symmetrisation; \m{\cyc(g)^{-e}-1}; the algebra homomorphisms of the proof of Theorem~D}
\nt{$\top$, $\mathbb{Q}/\mathbb{Z}$, $\mathrm{U}(1)$}{the trivial filter; the dualising group; the unitary group of rank one}
\end{multicols}
\endgroup

\head{References}
{\fontsize{9.4}{12.6}\selectfont
The reference is the one named by the second reading of this write-up, with the place given there; I did not
open it.\par
\smallskip
G.~Chenevier, \emph{The $p$-adic analytic space of pseudocharacters of a profinite group and
pseudorepresentations over arbitrary rings}.\par}

\end{document}
